\documentclass{book}%
\usepackage{graphicx} % pour l'image de la coupe d'Onésime
\usepackage[french]{babel}
\usepackage{fontspec}%
\usepackage{xunicode}%
%\usepackage{yfonts}%
%\usepackage{hyperref}% pour le courriel
\usepackage{amsmath}%
%\usepackage{cancel}% pour \inRas
\usepackage{quoting}%
\usepackage{ulem}% pour souligner en pointillé
\usepackage{comment}%
\usepackage{xcolor}%
\usepackage{graphicx}
%scale = {hscale, vscale}
%\usepackage[scale={0.75,0.8}]{geometry}%
%\setlength{\hoffset}{-18pt}        
%\setlength{\oddsidemargin}{20pt} % Marge gauche sur pages impaires ATTENTION écrit 2 fois


\newfontfamily{\policeFr}{GFS Bodoni}%
%\newfontfamily{\policeLat}{Book Antiqua}%
\newfontfamily{\policeLat}{GFS Neohellenic}%
%\newfontfamily{\policeAn}{Comic Sans MS}%
\newfontfamily{\policeAn}{Courier 10 Pitch}%
%\newfontfamily{\policeAll}{Fraktur}%
%\newfontfamily{\policeAll}{fraktur-bt.ttf}%
\newfontfamily{\policeAll}{Bertholdr Mainzer Fraktur}%
%\newfontfamily{\policeIt}{BodoniFLF}%
\newfontfamily{\policeIt}{GFS Artemisia}%
%\newfontfamily{\policeEsp}{Great Moravia}%
%\newfontfamily\policeGr{GFSBodoni.otf}[BoldFont = GFSBodoniBold.otf, ItalicFont = GFSBodoniIt.otf, BoldItalicFont = GFSBodoniBoldIt.otf]%
%\setmainfont{GFSBodoni.otf}[BoldFont = GFSBodoniBold.otf, ItalicFont = GFSBodoniIt.otf, BoldItalicFont = GFSBodoniBoldIt.otf]
\newfontfamily\policeGr{GFS Bodoni}%
\setmainfont{GFS Bodoni}
%\setdefaultlanguage[variant=ancient]{greek}%
\usepackage{polyglossia}%
\setdefaultlanguage{french}%
%\setotherlanguage[variant=ancient]{greek}%
%\setmainlanguage[variant=ancient]{greek} 
%\setotherlanguage[variant=ancient]{greek} 
%\setmainfont{Libertinus Serif} 
%\newfontfamily{\greekfont}{Liberation Serif}[
%  Script=Greek,
%  NFSSFamily=liberationserif,
%]

\usepackage[maxlevel=3]{csquotes}%
%\MakeOuterQuote{"}% pour les guillemets français
%\usepackage{verse}%


\newcommand{\italique}[1]{\textit{#1}}%
%\newcommand{\citComm}[1]{\begin{quoting}[vskip=0pt] #1 \end{quoting}}% vskip = \topsep + \partopsep par défaut
\newcommand{\espace}[1]{\hspace{#1}}%
\newcommand{\cesure}[0]{ [...] }%
\newcommand{\nieme}[1]{#1\italique{e}}%
\newcommand{\lettreBreve}[1]{$\breve{\mbox{#1}}$}
\newcommand{\synizese}[1]{$\widearc{\mbox{#1}}$}
\newcommand{\inRas}[2]{$\overset{\text{#2}}{\text{\sout{#1}}}$}
\newcommand{\auDessus}[2]{$\overset{\text{#2}}{\text{#1}}$}
%\newcommand{\etc}[0]{κτλ}%
%\newcommand{\cit}[0]{\espace{1cm}} % indent. pour une citation
\newcommand{\expo}[1]{\textsuperscript{#1}}%
\newcommand{\alinea}[0]{ \ \ }%
\newcommand{\Syn}[0]{ Συναγωγή }%
\newcommand{\SynB}[0]{ Συναγωγή\expo{B} }% la deuxième version pour la lettre A
\newcommand{\ClemAlex}{ Clemens Alexandrinus }
\newcommand{\Eustathe}{ Eustathius }
\newcommand{\AelDion}{ Aelius Dionysius }
\newcommand{\EtymMagnum}{ Etymologicum magnum }
\newcommand{\EtymGen}{ Etymologicum genuinum }
\newcommand{\EtymSym}{ Etymologicum Symeonis }
\newcommand{\Apostol}{ Michael Apostolius }
\newcommand{\Phryn}{ Phrynichus \italique{Praep. soph.} }
\newcommand{\Stobee}{ Joannes Stobaeus }
\newcommand{\EpimAlphab}{ Epimerismi alphabetici in Homerum (ed. Cramer) }
\newcommand{\GregCor}{ Gregorius Corinthius }
\newcommand{\Antiatt}{ Antiatticista (ed. Bekker) }
\newcommand{\Harp}{ Harpocration }% Autre édition de référence : Dindorf
\newcommand{\Rpc}[0]{R\expo{pc}}%
\newcommand{\Rac}[0]{R\expo{ac}}%
\newcommand{\Rdeux}[0]{R\expo{2}}%
\newcommand{\Mon}[0]{ M }%
\newcommand{\PsiUn}[0]{PSI-1194}
\newcommand{\PsiDeux}[0]{POxy-3839}
\newcommand{\PsiTrois}[0]{POxy-3840}
\newcommand{\PsiQuatre}[0]{POxy-4935}

\newcommand{\francais}[1]{{\policeFr #1}}%
\newcommand{\latin}[1]{{\policeLat #1}}%
\newcommand{\allemand}[1]{{\policeAll #1}}%
%\newcommand{\allemand}[1]{{\textfrak{#1}}}%
%\newcommand{\grec}[1]{{\policeGr #1}}%
\newcommand{\grec}[1]{#1}%
\newcommand{\anglais}[1]{{\policeAn \scriptsize{#1}}}%
\newcommand{\italien}[1]{{\policeIt #1}}%
\newcommand{\espagnol}[1]{{\policeEsp #1}}%

\newcommand{\Nauck}[0]{N\expo{2}}%
\newcommand*{\auteur}[1]{\textsc{#1}*}%
\newcommand{\cf}[0]{\italique{cf.}\xspace}%
\newcommand{\chap}[1]{\^{#1}}%

\newlength{\larg}%
\newcommand{\decalage}[1]{%
\settowidth{\larg}{#1}%
\espace{\larg}%
}
\newcommand{\decalageN}[1]{%
\settowidth{\larg}{#1}%
\espace{-\larg}%
}

\usepackage{xcolor}
\newcommand{\fantome}{{\color{white}v}}

\newcommand{\lecon}[1]{\textbf{#1}}
%\newcommand{\douteux}[1]{\dotuline{#1}} % pour les lettres douteuses : un point en-dessous
\newcommand{\douteux}[1]{$\underset{\cdot}{\text{#1}}$}
\newcommand{\difference}[1]{\dotuline{#1}}

\newcommand{\marqueur}[1]{\textsuperscript{#1}}%

\newcommand{\eff}[1]{%
{\color{orange}#1}%
}%

\newcommand{\LSJ}[0]{LSJ}%


\newcommand{\perso}[1]{\decalageN{\textbf{#1} --- }\textbf{#1}}% --- }
\newcommand{\Note}[1]{\section*{#1}%
\addcontentsline{toc}{section}{#1}}
\newcommand{\note}[1]{\footnote{#1}}
\newcommand{\biblio}[2]{\hspace*{-1.5cm}\textbf{[#1]}\hspace*{1.5cm}#2\\

}

\pagenumbering{Roman}


\usepackage{minipage-marginpar}
\usepackage{marginnote}
\usepackage[marginparwidth=3cm]{geometry}% pour les marges
%\usepackage{geometry}

\newcommand{\personnageGrec}[1]{%
	\reversemarginpar%
	\marginnote{%
		\textbf{\large{#1}}%\espace{1.6cm}
}}
\newcommand{\personnageFr}[1]{%
	\reversemarginpar%
	\marginnote{%
		\textbf{\large{#1}}%\espace{1.6cm}
}}

\newcommand{\citComm}[1]{\fantome\\
\espace{1cm}\begin{minipagewithmarginpars}{12.4cm}%
#1
\end{minipagewithmarginpars}%
}
\newcommand{\trimIamb}[1]{%
\espace{1cm}\begin{minipagewithmarginpars}{8.4cm}%
#1
\end{minipagewithmarginpars}%
}
\newcommand{\trimIambFr}[1]{%
	\espace{1cm}\begin{minipagewithmarginpars}{12cm}%
		#1
	\end{minipagewithmarginpars}%
}

\begin{document} 
\renewcommand{\contentsname}{Annexes}
\title{Annexes}
\author{}\date{}
\maketitle
\setcounter{page}{0}


\tableofcontents
\setcounter{page}{0}
\clearpage

%\clearpage
\setcounter{page}{1}

\Note{Notice sur Alphonse Willems par Maurice Leroy}

{\centering La quasi-intégralité de la traduction proposée des \italique{Thesmophories} ainsi que les notes qui composent ces annexes sont tirées de l'œuvre d'Alphonse Willems. J'y joins en guise de préambule quelques éléments biographiques, bien faible hommage que je puis rendre à cet immense savant\note{Les travaux sur Aristophane de Willems ont été réunis par ses élèves après sa mort dans son \italique{Aristophane, traduction avec notes et commentaires critiques} (1919). Les articles de Willems furent du vivant de leur auteur victimes de la conspiration du silence ; la verve facétieuse et irrévérencieuse de Willems ne ménageait certainement pas assez la susceptibilité des sommités aristophaniennes de l’époque. Ces blackboulés du talent préférèrent ignorer superbement tous les articles d’Alphonse Willems pour préserver leur petit milieu où règne l’esprit courtisan, microcosme d’entre-soi à tendance fortement onaniste… ou comment Aristophane fut sacrifié par ceux-là mêmes qui prétendaient œuvrer à sa diffusion et sa compréhension. En second lieu, Willems fut littéralement pillé sans vergogne par le Van Daele des Belles Lettres ; le traducteur-plagiaire de cette édition francophone de référence avait effectivement tout intérêt à ce que son lecteur ne sache rien de l’œuvre d’Alphonse Willems. On doit au grand Victor Coulon et au géant Taillardat d'avoir réhabilité Willems. Encore récemment, Anne Lebeau faisait republier ses traductions de \italique{Lysistrata} et de la \italique{Paix}, aux notes malheureusement tronquées.}.}\\
{\centering \italique{Notice sur Alphonse Willems,\\
membre de l'Académie,\\
Né à Saint-Josse-ten-Noode le 21 février 1839,\\
décédé à Schaerbeek le 27 novembre 1912.\\
Par Maurice LEROY.}\\
\italique{Extrait de l'}Annuaire de l'Académie royale de Belgique \italique{(1968).}}\\


\begin{figure}[!h]
	%\hspace*{-1cm}
	\centering
	\includegraphics[scale=0.68]{Willems.png}
	\caption{Alphonse Willems}
	\label{Willems}
\end{figure}

Alphonse-Charles-Joseph Willems était né à Saint-Josse-ten-Noode le 21 février 1839 d'une famille aisée — ce qui lui permit, tout au long de sa carrière, de s'adonner sans arrière-pensées aux études qui avaient sa prédilection — originaire de Bar-le-Nassau ; son père était médecin et tâta aussi de la politique locale : il fut bourgmestre de sa commune de 1842 à 1847. 
Toutefois, c'est au Collège des Joséphites de Melle, qui jouissait alors d'une réputation enviable, qu'il fit ses années d'humanités, études solides comme l'atteste la vaste culture dont, très jeune, il fit preuve mais dont l'empreinte dogmatique le marqua peu puisque dans la suite, à l'instar de plusieurs libéraux notables de son époque, il se convertit au protestantisme\note{On sent peut-être l'« antipapiste » dans une chronique de 1872 [13] où Willems parle du « pouvoir néfaste » qui, à Rome, a laissé détruire tant d'admirables monuments de l'Antiquité, — Les chiffres entre crochets renvoient à la Bibliographie imprimée à la suite de la présente notice.}.

En 1855, à l'âge de seize ans, il entre à l'Université Libre de Bruxelles et, trois ans plus tard, est proclamé, à la deuxième session de 1858, Docteur en Philosophie et Lettres ; l'année suivante, il suivra les cours de la Candidature en Droit mais n'ira pas plus loin dans ses études juridiques. 
Il convient de rappeler qu'à cette époque l'enseignement supérieur en Belgique était, en ce qui concerne la philologie, d'un niveau peu élevé : rien n'avait préparé le jeune docteur -- grade qui s'obtenait après deux ans de candidature et un an seulement d'un doctorat très peu spécialisé et qui ne comportait pas la présentation d'un mémoire — à une exégèse des œuvres littéraires, moins encore aux problèmes de la critique textuelle, matières dont il devait néanmoins devenir plus tard un maître incontesté. 

Le désir d'élargir ses connaissances et le goût inné pour le beau, joints à un travail personnel poursuivi avec enthousiasme et persévérance, comblèrent les lacunes.
La soif d'écrire le prit dès le début de ses études universitaires et ses publications attirèrent rapidement l'attention.
À la vérité, plusieurs domaines de recherche sollicitèrent son intérêt : néerlandiste à ses débuts, Alphonse Willems se sentit très vite captivé par les choses de l'art dans la connaissance desquelles — nous songeons surtout à la peinture et aux livres — il se montra plus qu'un amateur ; toutefois, à partir de la quarantaine, alors qu'il était en pleine possession de ses moyens, c'est aux études helléniques qu'il consacra le meilleur de ses forces et qu'il apporta la contribution la plus décisive.\\


\italique{Le néerlandiste.} Une tradition familiale\note{Nous devons ce détail, comme d'autres concernant la jeunesse de Willems, à notre Confrère F. Vanlangenhove : voir plus loin en tête de la Bibliographie.} rapporte que Willems, alors qu'il était élève à Melle, fut un jour injustement puni pour avoir conversé en flamand avec un de ses camarades : cet incident renforça peut-être en lui le désir d'étudier et d'approfondir la littérature néerlandaise.
Quoi qu'il en soit, dans le premier travail qu'il livra à l'impression — un article paru en deux fois dans la \italique{Revue Trimestrielle} de 1856 [1] et consacré à l'écrivain hollandais P. Hooft — ses premiers mots sont pour évoquer « les puissants motifs qui devraient nous engager, nous Belges, à étudier la langue flamande » et on lit un peu plus loin et dans la conclusion des déclarations que ne renieraient pas les chefs les plus intransigeants du mouvement nationaliste flamand («a-t-il [= le peuple flamand] besoin d'emprunter sa langue et sa littérature à un peuple étranger[= la France], à un peuple qui a toujours exercé une influence si pernicieuse sur son pays ? »), voire les partisans de la « grande Néerlande » (« nous considérons les Belges et les Hollandais, sinon comme un seul et même peuple, du moins comme deux peuples frères que des circonstances malheureuses ont pu seules désunir »); on voudra bien se souvenir que l'auteur de ces lignes avait alors dix-sept ans et parlait avec la fougue de sa jeunesse.

L'année suivante, c'est en néerlandais qu'il publie un pamphlet [4] où il s'élève contre ces écrivains de France ou de Belgique qui, dans la conviction que tout ce qui n'est pas français est « barbare », méconnaissent ou dédaignent les productions littéraires flamandes\note{Alph. Willems n'a pas mentionné ces deux travaux [1, 4] dans la bibliographie, par ailleurs complète, qu'il fournit au volume commémoratif du cinquantenaire de l'Université de Bruxelles : \italique{1834- 1884. L'Université de Bruxelles. Notice historique} par L. Vanderkindere (Bruxelles,1884), p. 214. Ils sont également absents de la bibliographie qui figure à la p. 79 de la brochure publiée par l'Université,à l'occasion de sa participation à l'Exposition organisée à Bruxelles en 1897, sous le titre : \italique{Université Libre de Bruxelles, 63\expo{ème} Année Académique. Notice sur l'Exposition Universitaire en 1897}, Bruxelles, 1897.} ; mais en même temps paraît une étude [3] sur la version néerlandaise du Roman de Renard et les controverses qui s'étaient élevées à son propos : s'il y est question à nouveau de la méconnaissance des lettres flamandes en Belgique, on sent cette fois, dans la discussion des sources, une maturité rare chez un jeune homme de cet âge.
C'est le même mérite que l'on reconnaît à la réédition qu'il donne en 1858, en la munissant d'une introduction et de variantes, du célèbre pamphlet contre l'église catholique romaine (1569) de Marnix de Sainte-Aldegonde [5] : contrairement à ses prédécesseurs, Willems reproduit scrupuleusement le texte de 1597, le dernier paru du vivant de l'auteur. 

Cet estimable travail philologique couronna aussi sa carrière de néerlandiste car, si nous faisons abstraction d'un bref article de 1871 [10] où il publia, retrouvés sur une feuille de parchemin qui avait servi de feuille de garde, 150 vers du \italique{Cassamus} (version flamande du roman français \italique{Le vœu du paon}), ses travaux ultérieurs furent consacrés à d'autres domaines, sans en excepter les lettres et la philologie françaises dont il entretint à plusieurs reprises les lecteurs de l'\italique{Indépendance Belge} [9,11,12,15,16] ; n'était-il pas aussi devenu entre-temps un des collaborateurs de la \italique{Nouvelle biographie générale} [7] publiée à Paris chez Firmin Didot ? 

Toutefois lorsque, le 30 octobre 1880, l'Université de Bruxelles proposa à Willems d'assurer, avec le titre de Professeur extraordinaire, l'enseignement, créé à son intention, de l'histoire de la littérature flamande\note{Il n'est pas sans intérêt de mentionner que c'est au jeune étudiant Willems que l'Université s'était adressée en 1856 pour qu'il fournît la version néerlandaise du discours d'allure polémique (on était alors au plus fort des discussions entre catholiques et libéraux à propos de l'encyclique de 1832 qui avait condamné les libertés des cultes, de la presse et d'association inscrites dans la Constitution Belge) qu'avait prononcé P.-Th. Verhaegen le 6 octobre 1856 à l'occasion de la séance solennelle d'ouverture des cours de l'année académique 1856-1857 [2].}, il accepta mais, au bout de trois ans, définitivement conquis par l'hellénisme, il renonça à ce cours sans pour autant qu'il reniât ses premières attaches : comme membre de l'Académie, il acceptera plus tard de faire rapport sur un mémoire de philologie néerlandaise [56] et, à deux reprises, de représenter ce corps savant à des Congrès de langue et littérature néerlandaises\note{Bruxelles, août 1906; Leyde, août 1908.}.
Ainsi, lorsque, au moment de la disparition du maître, Auguste Vermeylen soulignait, à juste titre, les services qu'il avait rendus à la langue et à la culture flamandes, c'est à l'activité de ses années de jeunesse qu'il se référait bien plus qu'à celle de son âge mûr. \\


\italique{Le collectionneur-bibliophile.} Une traduction du néerlandais, parue en 1859, montre l'intérêt que, tout jeune encore, Willems manifestait pour les choses de l'art\note{L'auteur de l'étude, P. Scheltema, présente dans sa préface la « traduction, dont a bien voulu se charger un jeune savant, très-ami des arts, M. A. Willems ».} : il s'agissait d'un essai sur la vie et l'œuvre de Rembrandt [6], sujet auquel il allait revenir dix ans plus tard dans un feuilleton de l'\italique{Indépendance Belge} [8]. 
En 1873, c'est d'une étude sur Frans Hals, placée en tête d'un très bel album d'eaux-fortes, qu'il fournissait la version française [14].
En fait, c'est là une préoccupation qui ne sera jamais absente de son activité et qui fait partie intégrante du personnage.
Son mariage même — qui eut lieu en 1862 à Corfou — n'avait pu que donner plus d'élan à son goût pour les Beaux-Arts : sa femme, Marie Tardieu, était en effet issue d'une lignée de graveurs français dont une branche était venue s'établir en Belgique peu avant 1830\note{Le frère de Madame Willems, Charles Tardieu (1838-1909), journaliste, critique d'art et wagnérien lucide (il devint membre de la classe des Beaux-Arts de l'Académie en 1893) fut longtemps rédacteur en chef et co-directeur du journal \italique{L'Indépendance Belge} aux suppléments littéraires duquel Willems donna plusieurs chroniques [8, 9, 11, 12, 13, 15, 16].}. 

En dehors de ces quelques écrits, sa passion pour l'art s'est concrétisée par la formation d'une galerie de tableaux où, à côté de deux Rubens et d'un J.-B. van Loo (école française du \nieme{XVIII} siècle), figuraient d'importantes œuvres modernes : Hipp. Boulenger, H. de Braekeleer, Ém. Claus, C. Corot, G. Decamps, Eug. Fromentin, V. Gilsoul, H. Leys, J.-B. Madou, G. Moreau, Th. Rousseau, Alfr. Stevens, Alfr. Verwée, Ant. Vollon, etc.\note{Le catalogue de la vente publique (qui eut lieu à Bruxelles le 12 mai 1921) de cette collection comptait 50 tableaux et 1 marbre (P. Devigne). — Plusieurs de ces toiles figurent aujourd'hui parmi les joyaux du Musée des Beaux-Arts de Bruxelles : le « Saint Michel terrassant les démons de Rubens, un portrait de van Loo, un autre d'Ém. Claus, un paysage de Corot, un « Paris » de Th. Rousseau ainsi que le tableau d'Henri Leys intitulé « La furie française à Anvers en 1583 ».}.

Comme bibliophile, Willems s'était acquis une réputation vraiment internationale et ce, non seulement par la collection qu'il s'était formée\note{La bibliothèque de Willems fut dispersée aux enchères en 1914, au cours de deux ventes qui eurent lieu, l’une à l'Hôtel Drouot à Paris les 4, 5, 6 et 7 mai, l'autre à Gand, Hôtel Cour St-Georges, les 25, 26, 27, 28 et 29 mai ; le catalogue de la première comportait 611 numéros, celui de la seconde 1851.} mais encore par les travaux de base qu'il publia dans ce domaine ; il avait en effet mené à bien l'étude minutieuse, approfondie, exhaustive de l'œuvre des Elzevier, cette illustre dynastie d'imprimeurs.
Un essai sur la première édition des Maximes de La Rochefoucauld [17] préluda à la parution de deux gros volumes (867 pages) sur Les Elzevier [18] qui sont incontestablement un chef-d'œuvre d'érudition mais d'une érudition qui ne se laisse jamais noyer par les détails : généalogie et biographie des quelque quatorze membres de la famille qui travaillèrent dans les officines de Leyde, de La Haye, d'Amsterdam et d'Utrecht (ainsi que dans les succursales de l'étranger), recensement et description minutieuse des quelque 2115 volumes que produisirent leurs ateliers.
On trouve dans ces colonnes\note{L'intérêt de Willems pour les Elzevier ne le quitta pas ; quelques années plus tard, il collaborait à l'établissement par G. Berghman — qui avait notamment poussé l'étude des fleurons et ornements typographiques — d'un supplément à son ouvrage [26] et accordait, sous forme de lettre, une préface au Catalogue d'une mise en vente exceptionnelle, par le libraire Damascène Morgand, de plus de 1300 ouvrages sortis pour la plupart des presses de la célèbre famille d'imprimeurs [22] ; en 1908 enfin, il fournissait, d'après une pièce des archives de Leyde, des précisions sur la date de naissance de Louis Elzevier [54].} non seulement les données techniques relatives à chaque volume mais encore le rappel des circonstances qui ont présidé à leur sortie de presse ; de même est levé le mystère qui planait sur pas mal de pseudonymes ou d'anonymes.
Et il n'est pas moins intéressant de noter, comme l'a fait Henri Grégoire, que « presque autant que des Elzevier, il est question, dans ces pages, des savants et des lettrés qui furent en relations avec les princes de l'art typographique, Daniel Heinsius et Saumaise -- le livre retentit de leurs querelles — Nicolas Heinsius surtout, sont, chemin faisant, dessinés d'un trait pittoresque.
Tel portrait révèle la sympathie de Willems pour ces grands humanistes qu'il fréquentait dès lors ».\\

\italique{L'helléniste}. Peu après lui avoir confié l'enseignement de l'histoire de la littérature flamande, l'Université de Bruxelles demandait à Willems, le 13 janvier 1881, de se charger de celui des textes grecs en Candidature.
Nommé Professeur ordinaire le 17 juin 1882 et déchargé du cours de littérature flamande le 4 août 1883, Alphonse Willems s'était vu confier encore, pour les années de Doctorat, les auteurs grecs (26 juillet 1884) et l'histoire de la littérature grecque (10 novembre 1888).
Il allait désormais, avec un éclat qu'on n'avait point connu jusqu'à lui, assurer l'enseignement du grec jusqu'au moment où, l'âge de la retraite ayant sonné (1909), il transmettrait le flambeau à Henri Grégoire\note{Complétons ici les renseignements concernant sa carrière à l'Université de Bruxelles. Il fut élu Délégué de sa Faculté au Conseil d'Administration pour les années académiques 1883-1884 et 1884-1885 ; ses Collègues lui confèrent la Présidence de la Faculté de Philosophie et Lettres à plusieurs reprises : 1886- 1887, 1887-1888, 1891-1892 et 1901-1902. En outre, de 1898 à 1900, il fut Délégué du Conseil d'Administration pour la direction de la Bibliothèque et Président de la Commission de la Bibliothèque. Enfin, c'est par décision du 24 juillet 1909 que lui fut accordé le titre de Professeur honoraire. 
	
Ajoutons que, titulaire de la Médaille commémorative du \nieme{75} anniversaire depuis le 25 janvier 1906 et de la Médaille de 1\expo{re} classe depuis le 18 mai 1906, il s'était successivement vu octroyer, dans l'Ordre de Léopold, les grades suivants : Chevalierle 5 octobre 1892, Officier le 7 septembre 1903 et Grand Officier le 1\expo{er} janvier 1912.}. 

Son enseignement de quelque trente années se caractérisa par la volonté constante de comprendre et de situer exactement les textes en prenant bien garde de ne pas les noyer sous des flots d'érudition.
Il expliqua des tragédies d'Euripide, des odes de Pindare, des poésies de Théocrite ou de l'Anthologie et, surtout, les comédies d'Aristophane.
Son premier cours avait porté sur l'\italique{Hippolyte} d'Euripide : il en sortit un volume qui, accepté par l'Académie (les rapporteurs avaient été Aug. Wagener, P. Willems et L. Roersch) en 1882, parut dans les Mémoires sous le titre \italique{Notes et corrections sur l'Hippolyte d'Euripide} [19] et le classa d'emblée parmi les maîtres de l'hellénisme.
Ce titre modeste (\italique{Notes sur...}) préfigurait celui d'une quinzaine de communications qu'il devait faire à l'Académie de 1894 à 1911 mais qui, celles-ci, porteront toutes sur Aristophane.

Rappelons en effet que notre Compagnie l'avait élu membre correspondant de la Classe des Lettres et des Sciences Morales et Politiques le 13 mai 1896 et membre titulaire le 9 mai 1900 ; assidu aux séances, il accepta aussi de faire à diverses reprises, avec la conscience qui le caractérisait, rapport sur des mémoires soumis à la Classe [23, 27, 28, 30, 31, 35, 36, 41, 44, 51, 56, 60]\note{Notons aussi qu'à plusieurs reprises Willems prit la parole aux réunions de la Société pour le progrès des études philologiques et historiques comme en témoigne l'\italique{Annuaire-Bulletin} de cette Société : 1899, pp. 96 et 103 ; 1900, pp. 9-13 ; 1902. pp. 13-15.}. 

Cette vocation d'helléniste pourrait paraître bien tardive (Willems avait quarante-quatre ans lorsqu'il publia son premier travail en la matière) ; il est vraisemblable au contraire qu'elle l'avait hanté dès ses débuts et avait mûri au cours des années.
Un événement fut très probablement déterminant : le voyage en Orient et en Grèce qu'il entreprit en 1862 à une époque où une expédition de ce genre relevait encore quelque peu de l'aventure ; ce séjour en Grèce\note{C'est par erreur, semble-t-il, que Grégoire a dit de Willems qu'il avait été consul de Belgique dans l'Athènes de l'Interrègne » car le Ministère des Affaires Étrangères ignore tout de cette fonction qu'il aurait remplie ; d'ailleurs Lonchay parle simplement de « son long séjour à Athènes » sans faire allusion à un quelconque poste consulaire.} se plaça pendant les années tumultueuses (mais en fut-il jamais d'autres dans ce pays ?) 1862 et 1863 où fleurirent les mutineries et qui virent l'écroulement de la monarchie othonienne et l'avènement de la dynastie danoise.
Mais pendant ces quelques mois qui durent être particulièrement lumineux pour le jeune homme — c'est, nous l'avons dit, à Corfou, où sa fiancée et sa future belle-mère étaient venues le rejoindre, qu'eut lieu son mariage — que de notations prises sur le vif, que de contacts avec un peuple athénien si peu différent de celui du cinquième siècle avant notre ère, que d'observations accumulées dont il tirera un si judicieux profit pour l'exégèse du grand comique. 
Et on peut imaginer la joie que Willems vieillissant, devenu le plus expert des aristophaniens, dut éprouver lorsque, au printemps de 1905, il reprit contact avec la Grèce à l'occasion du Congrès international d'archéologie d'Athènes. 

Revenons un instant à ces années qui suivirent le retour en Belgique d'un Alphonse Willems tout baigné d'hellénisme : c'est alors qu'il entreprit de patientes recherches sur l'activité des Elsevier et la familiarité qu'il en acquit avec les grands humanistes contribua sans doute aussi à le rapprocher du grec. 
Quoi qu'il en soit, s'il accepta sans trop d'hésitation la chaire de grec que lui offrait l'Université — mais qu'il n'avait point sollicitée — c'était comme une mission qu'il avait conscience de remplir : jusqu'à sa retraite, il s'en acquitta avec un scrupule exemplaire, préparant et polissant minutieusement ses leçons ; Grégoire rapporte ce trait de modestie : « Willems avouant à ses élèves qu'il avait oublié le manuscrit de sa leçon, et renonçant à la faire ce jour-là, ne voulant dicter qu'une traduction bien étudiée. Un autre eût peut-être profité de cette occasion pour jouer, sans grand péril, la difficulté ». 

C'est donc à Aristophane que Willems consacra le meilleur de sa production scientifique et cette prédilection remonte peut-être à ses années d'études: n'est-ce pas déjà une citation du grand comique que nous trouvons en deuxième page de l'étude sur la version flamande du Roman de Renard qu'il avait publiée à l'âge de dix-huit ans ?\note{[3], p. 286 : « nous sommes tentés de murmurer avec le comique grec, οὐ γὰρ πείσεις, οὐδ᾽ ἢν πείσῃς : tu ne me persuaderas pas, quand tu le prouverais » (Il s'agit du vers 600 du \italique{Ploutos}).}
Après ses \italique{Notes} sur Euripide mises au point en 1882, le maître semble s'être accordé un temps de réflexion (il y a un «trou » de onze années dans sa bibliographie), plus exactement le loisir d'étudier à fond le comique athénien devenu son auteur de prédilection et d'acquérir avec son œuvre la familiarité indispensable.
En tout cas, de 1894 à sa mort (1912), toutes ses publications vont porter sur Aristophane et elles constituent, tant pour l'établissement du texte que pour l'explication de maint passage jugé jusqu'alors difficile voire désespéré un commentaire tellement lumineux et décisif que les rectifications et interprétations qui y sont proposées ont été quasi unanimement adoptées par les éditeurs postérieurs. 

À cette méditation continue sur Aristophane -- commentaire et traduction — qui s'est exprimée dans quelque vingt-sept mémoires\note{Dix-sept d'entre eux (dont douze portent comme titre \italique{Notes sur...}) ont paru dans les \italique{Bulletins de l'Académie Royale de Belgique} [20, 21, 25, 29, 32, 33, 34, 37, 38, 42, 45, 48, 50, 52, 53, 57, 59), dix (dont deux \italique{Notes sur...}) dans divers revues et recueils [24, 39, 40, 43, 46, 47, 49, 55, 58, 61].} et dont une partie était restée inédite\note{Willems avait entrepris en 1909 de publier séparément la traduction des pièces d'Aristophane ; trois seulement avaient paru de son vivant dans la \italique{Revue de l'Université de Bruxelles} [55, 58, 61] mais, par bonheur, le manuscrit des huit autres pièces était également prêt pour l'impression.} — mais l'ensemble, recueilli avec un soin pieux par deux des disciples du maître : Ch. Lacomblé et Léon Leclère, a paru dans les trois beaux volumes [62, 63, 64] de l'édition posthume de 1919\note{C'est à elle qu'il convient de se référer pour citer l'œuvre de Willems, car les éditeurs ont pu utiliser des additions et corrections que l'auteur avait notées sur ses exemplaires personnels ; c'est pourquoi, dans notre bibliographie, nous avons chaque fois, pour les mémoires portant sur Aristophane, indiqué la référence à l'édition de 1919.} -- il nous semble distinguer trois mérites essentiels. 

Et tout d'abord le respect du texte reçu. 
Partant de l'idée que les œuvres de l'Antiquité avaient subi au cours de la tradition manuscrite un grand nombre de corruptions, les philologues du \nieme{XIX} siècle s'étaient trop souvent cru autorisés à y apporter des corrections arbitraires ; avec beaucoup de bon sens, Willems s'est élevé contre cette façon de voir avant même que les trouvailles papyrologiques ne vinssent confirmer l'exactitude de ses vues.
Servi par une connaissance érudite des sources mais aussi par un sens profond de la langue, il n'a pas de peine à démontrer l'inanité de tant de conjectures, fussent-elles proposées par des hellénistes en renom ; sans fard, disons même sans aménité — on sent chez lui l'indignation de voir son poète favori massacré par des infidèles ! — il s'attaque aux coryphées de la philologie trop sûrs d'eux-mêmes et trop prompts à estimer « corrompu » tout passage dont le sens dépasse leur entendement ; c'est miracle de voir avec quelle ingénuité Willems abat des constructions pseudo-savantes, avec quelle verve il repousse de prétendues émendations, avec quelle vigueur il rejette l'argument d'autorité.
Si son nom apparaît si rarement dans l'apparat critique des éditions modernes, c'est que sa tâche a le plus souvent consisté à restaurer, là où elle avait été abîmée par des exégètes trop ingénieux, la leçon des manuscrits ; il conviendrait cependant de ne pas oublier ce bon combat qu'il a mené pour l'authenticité de la tradition. 

Le deuxième mérite de Willems est qu'il donne de nombreux passages déclarés obscurs une explication aussi simple que convaincante, et ce grâce à sa connaissance précise des mœurs athéniennes, une aperception juste de la psychologie populaire, un recours aux données de l'archéologie et aux \italique{realia}. Ici, il a été grandement servi par son expérience de la Grèce moderne ; laissons encore la parole à Henri Grégoire :  «Les éditeurs d'Aristophane étaient, et sont encore, ignorants de bien des détails de mœurs et surtout de cuisine, qu'il suffisait d'étudier dans la vie orientale. Rien de plus savoureux que certaines dissertations de Willems sur des points d'antiquité vivante, si l'on peut dire. Discutant avec des érudits de cabinet, il avait la supériorité du voyageur qui a respiré les relents des \italique{bakals}, goûté aux chefs-d'œuvre des rôtisseries levantines, fréquenté le marché aux poissons, où les contemporains d'Aristophane allaient chercher les mêmes délicatesses qui plaisent aujourd'hui à leurs descendants.
Il est délicieux de l'entendre parler en connaisseur d'\italique{aphyes} ou de \italique{pinnotères}, vanter le `caillé d'amouille', la cervelle rôtie dans une feuille de figuier, et remontrer à M, van Leeuwen que la φυστὴ μᾶζα, « loin de ressembler à une omelette soufflée, n'était qu'une simple galette comme celle qui, à Athènes, servait de pain de munition ». 

Enfin il y a la traduction et reconnaissons tout de suite que cette partie de son œuvre n'est pas à l'abri des critiques ; dans le but louable de rendre avec plus de précision l'original, Willems s'est livré à une recherche lancinante du mot juste, à la chasse au terme qui, sans équivoque, rendrait en français le mot grec correspondant. 
Il l'a fait en recourant aux couches les plus diverses du vocabulaire : archaïsmes, régionalismes, termes techniques, argots ; si, prises séparément, beaucoup de ces trouvailles sont incontestablement heureuses et donnent le ton juste, l'impression d'ensemble ne laisse pas d'être quelque peu disparate, si bien que ce manque d'unité trahit dans une certaine mesure Aristophane dont la langue, au dire même de Willems, était précisément d'un attique pur et châtié.
Mais, ces réserves mises à part, la traduction de Willems n'en reste pas moins une des meilleures que nous possédions dans une langue moderne et, en tout état de cause, elle constitue un commentaire suivi fort précieux des Comédies : car on peut être assuré que rien n'y avait été laissé au hasard et que chaque ligne du texte avait été soigneusement pesée avant d'être couchée sur papier.\\

La forte personnalité de Willems transparaît dans son œuvre ; il maniait sans pitié l'ironie et sa verve polémique était intarissable : combien de philologues aux noms pourtant illustres ont expié dans ses écrits l'audace de leurs conjectures ou leur méconnaissance des mœurs antiques ! Exigeant envers lui-même, il ne pouvait que l'être aussi en famille comme au milieu de ses élèves\note{Albert Stassart (qui écrivit sous le nom de Georges Rency) avait été l'élève de Willems en 1897-1899. Dans ses \italique{Souvenirs de ma vie littéraire} parus en 1940, il rappelle les leçons de ce maître éminent « qui nous expliquait si bien, avec autant d'amour que de sagace érudition, Théocrite, Euripide et le divin Aristophane ». Or, en 1910, alors que le maître venait de prendre sa retraite, le même Rency avait eu l'inélégance de publier une nouvelle, \italique{Rose d'Octobre} — qui est par ailleurs du plus mauvais goût mélodramatique — où il mettait en scène un professeur de grec à l'Université, pris tout entier par une traduction d'Aristophane, fervent de \italique{realia} et bibliophile par surcroît ; ce savant, baptisé Jérôme Gaspardot, avait sacrifié, par son égoïsme aveugle et sa passion pour les vieux livres, sa femme et sa fille, ne s'étant en somme aperçu de leur existence que lorsque l'une et l'autre furent mortes de consomption : on voit le genre. Il est difficile de ne pas reconnaître en Gaspardot notre Alphonse Willems -- qui, en son temps, avait dû étriller Stassart-Rency pour la piètre valeur de sa thèse de doctorat sur le roman grec ! -- et le nom Gaspardot pourrait aussi être un trait décoché par Rency à son condisciple Camille Gaspar, autre élève, mais brillant celui-là, de Willems ; c'est ce que vient de montrer G. Cambier dans un article de la \italique{Revue Nationale}, t. XXXIX, 1967, pp. 275-277 : \italique{Alphonse Willems et Georges Rency}.} car il s'était fait une haute idée de sa mission d'enseignement et il savait charmer ses élèves en les instruisant, ayant à cœur de leur former le jugement et le goût. 

Voici comment s'exprimait l'un d'entre eux, Ch. Lacomblé : « Pour qui apportait dans ses études autre chose que le souci de l'examen et de la place à conquérir, le maître devenait le guide le plus bienveillant, le plus précieux ; aimable, accueillant, prodiguant les trésors de sa haute culture, s'intéressant aux progrès, s'associant aux travaux de ceux en qui il espérait éveiller quelque vocation».
Et encore: « les qualités maîtresses de son esprit et de son enseignement dominèrent son travail : le sérieux, le jugement, le respect du texte, l'étude profonde et consciencieuse, la défiance des autorités, le courage d'aborder les difficultés de front et de source, l'horreur de l'à peu près, du recours présomptueux à la conjecture, du pédantisme serein et satisfait, qui change ce qu'il ne comprend pas et écrit `ce n'est pas grec', parce qu'il ne sait pas dire `\italique{non liquet}'. 

Atteint par la limite d'âge en 1909, Alphonse Willems abandonna son enseignement mais non ses recherches ; c'est alors en effet qu'il mit au point sa traduction d’Aristophane et en commença — sans pouvoir la mener au terme — la publication. 
Les derniers mois de sa vie furent assombris par la maladie de sa femme plongée dans un état d'inconscience qui le désespérait ; elle lui survécut de peu. Quant à lui, c'est brusquement qu'il fut emporté, en pleine lucidité, en plein travail, le 27 novembre 1912 dans sa demeure d'allure modeste mais qu'il avait patiemment et amoureusement peuplée d'un trésor de peintures et de livres précieux.\\

\Note{Biblographie d'Alphonse Willems (établie par Maurice Leroy)}

On trouvera des éléments de la biographie d'Alphonse Willems dans les articles suivants : Henri GRÉGOIRE, \italique{Alphonse Willems} dans \italique{Bulletin de la Société pour le progrès des études philologiques et historiques}, Année 1913, pp. 121-132. (texte qui a été repris dans \italique{Université de Bruxelles. Annuaire pour l'année académique 1913-1914}, pp.65-73). 

\italique{Mort de M. Alphonse Willems} dans \italique{Revue de l'Université de Bruxelles}, t. XVIII (1912-1913), pp. 276-281 (y figurent notamment les discours prononcés aux funérailles par P. Errera, H. Lonchay et Aug. Vermeylen). 

Ch. Lacomblé, \italique{Notes d'un vieil élève} dans l'édition posthume de la traduction d'Aristophane, t. 1, pp. XXIX-XL [62]. 

Maurice Leroy, \italique{Alphonse Willems} dans \italique{Biographie Nationale}, t. XXXI, 2 (1962), col, 737- 739.\\

Notre Confrère Fernand Vanlangenhove -- Madame Vanlangenhove est la petite-fille d'Alphonse Willems —  nous a procuré le portrait du maître et a bien voulu nous éclairer sur plus d'un point de sa biographie : nous l'en remercions vivement ici. 

Pour reconstituer la bibliographie de Willems, nous nous sommes aidés des deux volumes commémoratifs des cinquantième et soixante-quinzième anniversaires de l'Université de Bruxelles, celui de Vanderkindere cité ci-dessus et celui du C\expo{te} Goblet d'Alviella, 1884-1909. 
L'Université de Bruxelles pendant son troisième quart de siècle (Bruxelles,1909), pp. 206-207, de la \italique{Bibliographie Nationale} (cf. ci-après [6], de la notice qui figure dans la brochure sur l'Exposition de 1897 et, pour tout ce qui regarde Aristophane, des trois volumes de l'édition posthume de 1919 [62,63,64].

Dans la liste qui suit, dressée par ordre chronologique, l'abréviation \italique{Bull.} renvoie, pour les années 1894 à 1898, aux \italique{Bulletins de l'Académie Royale des Sciences, des Lettres et des Beaux-Arts de Belgique}, pour les années 1899 à 1912 au \italique{Bulletin de la Classe des Lettres et des Sciences Morales et Politiques et de la Classe des Beaux-Arts de l'Académie Royale de Belgique}.\\


\biblio{1}{De la littérature néerlandaise au \nieme{XVII} siècle. Pierre Hooft, \italique{Revue Trimestrielle}, t. X, 1856, pp. 190-226 et t. XI, 1856, pp. 121-147.}%
\biblio{2}{Vrye hoogeschoolte Brussel. Redevoering in openbare zitting namens de Hoogeschool uitgesprokendoor den heer Verhaegen, bestierder-opzigter, by de heropeningder leergangen, den 6° october 1956. Vertaald door Alp. Willems. 15 pp., Bruxelles, Al. Cadot, 1856.}%
\biblio{3}{Étude sur le poëme Van den Vos Reinaerde, d'après le travail récent de M. Jonckbloet. \italique{Messager des Sciences historiques ou Archives des Arts et de la Bibliographie de Belgique}, année 1857, pp. 285-321.}%
\biblio{4}{Wat men in Frankrijk zoo al over onze letterkunde denkt. 23pp., Bruxelles, F. van Meenen, 1857.}%
\biblio{5}{De Werken van Ph. van Marnix van Sint Aldegonde. De Bijenkorf der H. Roomsche Kerke, met Inleiding en Varianten. 2 vol., XIV-324 et 264 pp., Bruxelles, Fr. van Meenen, 1858. — La préface (pp. V-XIV) est signée A. W. ; le texte reproduit l'édition de 1597.}%
\biblio{6}{Rembrand, Discours sur sa vie et son génie, avec un grand nombre de documents historiques, par le D\expo{r} P. Scheltema. Traduit par A. Willems, revu et annoté par W. Burger, 87 pp., Bruxelles, A. Labroue, 1859. — Ce texte avait d'abord paru en trois livraisons dans la \italique{Revue Universelle des Arts}, t. VIII, 1858, pp. 273-209, 369- 301 et 485-516. D'autre part, la \italique{Bibliographie Nationale. Dictionnaire des écrivains belges et catalogue de leurs publications. 1830-1880} signale (t. IV, Bruxelles, 1910, p. 350) que W. Burger est le pseudonyme de Thoré et ajoute : « Une nouvelle édition, corrigée et augmentée, a été publiée et annotée par W. Burger, Paris, \italique{V}e\italique{ Renouard}, 1866. In-89, XI-156 p. Le nom du traducteur a disparu du titre ».}%
\biblio{7}{Collaboration à la \italique{Nouvelle biographie générale publiée sous la direction du Dr Hœfer}. Paris, Firmin Didot, 46 vol., 1857 à 1866 ; dans la liste des collaborateurs imprimée à la fin du t. 46, on lit : « Willems (Alphonse), à Bruxelles » mais, parmi les quelque 45.600 colonnes de texte qu'offre cette œuvre considérable, il serait peu aisé de repérer les collaborations de Willems, d'autant plus que pas mal d'articles ne sont pas signés.}%
\biblio{8}{Rembrandt Harmensz. Van Rijn, sa vie et ses œuvres, par C. Vosmaer. L'\italique{Indépendance Belge}, 26 juin 1869.}%
\biblio{9}{Le Dictionnaire de Littré. L'\italique{Indépendance Belge}, 6 janvier 1870.}%
\biblio{10}{Het Brusselsch fragment van den Cassamus. \italique{De taal- en letterbode}, t, II, 1871, pp. 158-166.}%
\biblio{11}{Les races romanes, par H. Kern. L'\italique{Indépendance Belge}, 10 et 11 novembre 1872.}%
\biblio{12}{G. de Ville-Hardouin, La conquête de Constantinople. L'\italique{Indépendance Belge},24 décembre 1872.}%
\biblio{13}{Rome, par Francis Wey. L'\italique{Indépendance Belge}, 28 décembre 1872.}%
\biblio{14}{Eaux-Fortes d'après Frans Hals par William Unger. Avec une étude sur le maître et ses œuvres par C. Vosmaer. Traduction française d'Alph. Willems. Leyde, A. W. Sijthoff, 1873, 36 pp. in-f°, 20 pl.}%
\biblio{15}{Joinville, Histoire de Saint-Louis. L'\italique{Indépendance Belge}, 24 décembre 1873.}%
\biblio{16}{L'homme au masque de fer. L'\italique{Indépendance Belge}, 17 février 1874.}%
\biblio{17}{La première édition des Maximes de La Rochefoucauld imprimée par les Elzevier en 1664. Notice bibliographique, 16 pp., Bruxelles, G. A. Van Trigt, 1879. — Ce texte a été réimprimé dans l'édition de La Rochefoucauld des \italique{Grands Écrivains de la France, Appendice au tome premier} (Paris, Hachette, 1883). — Ce t. I avait paru en 1868.}%
\biblio{18}{Les Elzevier. Histoire et Annales typographiques. 2 vol, CCIX-607 pp., 3 pl et 1 tableau généalogique. Bruxelles, G.-A. Van Trigt — Paris, Ad. Labitte, 1880 (cf. [26]).}%
\biblio{19}{Notes et corrections sur l'Hippolyte d'Euripide. \italique{Mémoires in-8°de l'Académie Royale de Belgique}, tome XXXVI, 1883, 74 pp. Bruxelles, G.-A. Van Trigt.}%
\biblio{20}{Notes sur les \italique{Guêpes} d'Aristophane. À propos d'une édition récente, \italique{Bull.} 3° série, t. 27, 1894, pp. 403-421 $=$ [62], pp. 517-534. — Il s'agit de l'éd. J. van Leeuwen publiée à La Haye, en 1893.}%
\biblio{21}{Notes sur les Oiseaux d'Aristophane. \italique{Bull.} 3° série, t. 32, 1806, pp. 603-635 $=$ [63], pp. 254-283.}%
\biblio{22}{Éd.Rahir. Catalogue d'une collection unique de volumes imprimés par les Elzevier et divers typographes hollandais, précédé d'un avant-propos par M. Ferdinand Brunetière et d'une lettre de M. Alphonse Willems. XXIV-493 pp. Paris, D. Morgand, 1896. — La lettre de Willems occupe les pages XV à XVII.}%
\biblio{23}{Rapport sur le mémoire du concours de 1897 : Les croyances et les cultes de l'île de Crète dans l'antiquité. \italique{Bull.}, 3° série, t. 33, 1897, pp. 485-488. — Ce mémoire, non accepté, est resté anonyme.}%
\biblio{24}{Notes sur les Grenouilles. À propos d'une édition récente. \italique{Revue de l'Instruction Publique en Belgique}, t. XL, 1897, pp. 233-239 $=$ [64], pp. 127-134. — Il s'agit de l'éd. J. van Leeuwen publiée à La Haye en 1896.}%
\biblio{25}{Une traduction nouvelle d'Aristophane. \italique{Bull.} 3°série, t. 34, 1897, pp. 970-992 $=$ [64], pp. 361-380. — Il s'agit de la traduction d'Eugène Talbot, Paris, 2 vol, 1897.}%
\biblio{26}{G. Berhman, Supplément à l'ouvrage sur Les Elzevier de M. Alphonse Willems. XVIIL-173 pp., Stockholm, Iduns, 1897. -- La dédicace (p. V) est ainsi rédigée : \italique{À Monsieur Alphonse Willems qui a bien voulu honorer mon travail en y prêtant son concours}.}%
\biblio{27}{Rapport sur un mémoire présenté pour le Prix Stassart : Notice sur un Belge Célèbre (Nicolas Clénard). \italique{Bull.}, 1898, pp. 617-621. — Ce mémoire, renvoyé pour remaniement, a été représenté l'année suivante ; cf. [31].}%
\biblio{28}{Rapport sur H, Delehaye. \italique{Note sur la légende de la lettre du Christ tombée du ciel}. \italique{Bull.}, 1899, p. 136.}%
\biblio{29}{Notes sur les Cavaliers d'Aristophane. À propos d'une édition récente. \italique{Bull.}, 1899, pp. 137-168 $=$ [62], pp. 213-241. — Il s'agit de la \nieme{2} éd. publiée chez Teubner en 1897 par Ad. von Velsen.}%
\biblio{30}{Rapport sur le mémoire du concours de 1897 : \italique{Le temple d'Esculape à Épidaure. Bull.}, 1899, pp. 277-279. — Ce mémoire, non accepté, est resté anonyme.}%
\biblio{31}{Rapport sur le mémoire du concours pour le Prix de Stassart : V. Chauvin et Alph. Roersch, \italique{Nicolas Cleynaerts dit Clenardus. Bull.}, 1899, pp. 300-301.}%
\biblio{32}{Notes sur la Paix d'Aristophane. \italique{Bull.}, 1899, pp. 861-898 $=$ [63], pp. 93-130.}%
\biblio{33}{Note sur un passage des Guêpes, \italique{Bull.}, 1899, pp. 898-900 $=$ [62], pp. 585-587.}%
\biblio{34}{Notes sur deux passages des Oiseaux. \italique{Bull.} 1899, pp. 900-905 $=$ [63], pp. 285-280.}%
\biblio{35}{Rapport sur J. Laminne, Le traité περὶ ἑρμηνείας d'Aristote. \italique{Bull.}, 1901, p. 345.}%
\biblio{36}{Rapport sur le mémoire du concours de 1901 : H. Francotte, \italique{La formation des villes, des États, des confédérations et des ligues dans la Grèce ancienne. Bull.}, 1901, pp. 945-948.}%
\biblio{37}{Le nu dans la comédie ancienne des Grecs, avec un appendice sur la répartition des places au théâtre. \italique{Bull.}, 1901, pp. 1073-1096 $=$ [64], pp. 381-400.}%
\biblio{38}{Notes supplémentaires sur les \italique{Guêpes} et les \italique{Cavaliers} d'Aristophane. \italique{Bull.}, 1901, pp. 1135-1182 $=$ [62], pp. 535-560 et 243-258.}%
\biblio{39}{Notes sur l'Ecclésie des Femmes d'Aristophane. \italique{Album gratulatorium in honorem Henrici Van Herwerden}, Utrecht, 1902, pp. 255-260 $=$ [64], pp. 253-258.}%
\biblio{40}{Du choix des mots chez les Attiques. À propos d'un passage des « Cavaliers » d'Aristophane, \italique{Revue de l'Instruction Publique en Belgique}, t. XLVI, 1903, pp. 1-7 $=$ [64], pp. 401-408.}%
\biblio{41}{Rapport sur le mémoire couronné par le Prix Gantrelle en 1903 : Ém. Boisacq, \italique{Les résultats acquis dans le domaine de l'étymologie grecque depuis la dernière édition des Grundszüge de G. Curtius} (1879). \italique{Bull.}, 1903, pp. 168-273.}%
\biblio{42}{Notes sur les Acharniens d'Aristophane suivies de notes supplémentaires sur les Oiseaux, \italique{Bull.}, 1903, pp. 617-687 $=$ [62], pp. 81-109 et [63], pp. 291-321.}%
\biblio{43}{Les régates à Athènes. À propos d'un passage d'Aristophane, \italique{Revue de l'Instruction Publique en Belgique}, t. XLVII, 1904, pp. 81-87 $=$ [64], pp. 417-424.}%
\biblio{44}{Rapport sur J. Bidez. \italique{Notes sur les lettres de l'empereur Julien. Bull.}, 1904, p. 470.}%
\biblio{45}{Notes sur la Lysistrata d'Aristophane. \italique{Bull.} 1904, pp. 507-645 $=$ [63], pp. 413-454.}%
\biblio{46}{Qu'est-ce que la monodie crétique ? \italique{Mélanges Paul Frédéricq}, Bruxelles, 1904, pp.3-8 $=$ [64], pp. 409-415.}%
\biblio{47}{Les Athéniennes au théâtre, \italique{Revue de l'Instruction Publique en Belgique},t. XLVIII, 1905, pp. 245-252 $=$ [64], pp. 425-433.}%
\biblio{48}{Les Athéniens à l'Ecclésie. \italique{Bull.}, 1905, pp. 809-819 $=$ [64], pp. 435-443.}%
\biblio{49}{Aristophane. Cavaliers, 537-540. \italique{Revue des Grecques}, t. XIX, 1906, pp. 383-388 $=$ [62], pp. 259-264.}%
\biblio{50}{Notes sur les Nuées d'Aristophane. \italique{Bull.} 1906, pp. 635-682 $=$ [62], pp. 375-415.}%
\biblio{51}{Rapport sur le mémoire couronné par le Prix Gantrelle en 1907 : J. Bidez, \italique{Recueil critique des fragments philosophiques de Porphyre le néo-platonicien. Bull.}, 1907, pp. 289-292.}%
\biblio{52}{Aristophane et la démocratie athénienne. \italique{Bull.}, 1907, pp. 338-373 $=$ [64], pp. 445- 473.}%
\biblio{53}{Notes sur les Thesmophories et le Ploutos d'Aristophane. \italique{Bull.}, 1908, pp. 625-693 $=$ [63], pp. 535-579 et [64], pp.345-358.}%
\biblio{54}{Lodewijk Elsevier's geboortejaar. \italique{Tijdschrift voor boek- en bibliotheskwezen}, t. VI, 1908, pp. 72-74.}%
\biblio{55}{Les Guêpes. Comédie d'Aristophane, jouée à Athènes, sous l'archontat d'Aminias, aux Lénéennes (février) de 422. \italique{Revue de l'Université de Bruxelles}, t. XV, 1909-1910, pp. 81-148 $=$ [62], pp. 417-516.}%
\biblio{56}{Veslag over J. V. Van Mirlo, Jr., \italique{Studie over « Zuster Hadewijch»}. \italique{Bull.}, 1910, pp. 147-150. — Non accepté, ce mémoire fut entièrement remanié et couronné en 1912.}%
\biblio{57}{Notes sur l'Ecclésie des Femmes et sur les Grenouilles d'Aristophane. \italique{Bull.}, 1911, pp. 219-272 (à la page 256, le sous-titre porte par erreur « Notes sur les Cavaliers » au lieu de « Notes sur les Grenouilles »  $=$ [64], pp. 215-252 et 109-125.}%
\biblio{58}{Les Nuées, Comédie d'Aristophane, jouée à Athènes, sous l'archontat d'Isarchos, aux Grandes Dionysiaques (mars) de 423. \italique{Revue de l'Université de Bruxelles}, t. XVII, 1911-1912, pp. 81-152 $=$ [62], pp. 267-374.}%
\biblio{59}{La recension des « Guêpes » d'Aristophane par M. de Wilamowitz-Moellendorf. \italique{Bull.}, 1912, pp. 47-70 $=$ [62], pp. 561-584.}%
\biblio{60}{Rapport sur L. Parmentier, \italique{Recherches sur le traité d'Isis et d'Osiris de Plutarque. Bull.}, pp. 361-363.}%
\biblio{61}{Les Cavaliers. Comédie d'Aristophane, jouée à Athènes, sous l'archontat de Stratoclès, aux Lénéennes (février) de 424. \italique{Revue de l'Université de Bruxelles}, t. XVIII, 1912-1913, pp. 81-155 $=$ [62], pp. 113- 212.}%
\biblio{62}{Aristophane. Traduction avec Notes et commentaires critiques. 3 vol, XL-595, 585 et 477 pp. Paris, Hachette — Bruxelles, J. Lebègue, 1919. (Cette édition posthume est l'œuvre de deux anciens élèves de Willems : Ch. Lacomblé et L. Leclère). Tome 1 : Préface (par L. Parmentier, pp. I-XXVII) — Notice biographique (par Ch. Lacomblé, sous le titre \italique{Notes d'un vieil élève}, pp. XXIX-XL). 
Les Acharniens— Les Cavaliers ($=$ [61]) — Les Nuées ($=$ [58]) — Les Guêpes ($=$ [55]). 

Notes sur ces quatre pièces ( $=$ [20, 29, 33, 38, 42, 49, 50, 55,58, 59 et 61]). 

(Signalons, dans la Table des matières de ce vol. I, p. 595, une erreur qui n'a pas été relevée dans les \italique{Errata} : pour les Notes supplémentaires sur les Cavaliers, II, il faut lire pp. 243-258 au lieu de pp. 243-264 et ajouter : Notes sur les Cavaliers, II, Vers 537-540, pp. 259-264).}%
\biblio{63}{Tome II : La Paix — Les Oiseaux — Lysistrata — Les Thesmophories. Notes sur ces quatre pièces ($=$ [21, 32, 34, 42, 45 et 53]).}%
\biblio{64}{Tome III : Les Grenouilles — L'Ecclésie des Femmes — Ploutos. Notes diverses ( $=$ [24, 39, 53 et 57]) et études détachées sur Aristophane ( $=$ [25, 37, 40, 43, 46, 47, 48 et 52]).}%



\Note{\italique{Alphonse Willems. Notes d'un vieil élève.} (par Charles Lacomblé.)}

{\centering \italique{Extrait du t. I, pp. XXIX-XL.}}\\

Quand je pense au maître que nous avons perdu, je retrouve en lui un véritable humaniste. 

Sa haute image se détache, sur le fond terne et affairé de notre milieu, dans la sérénité du culte des bonnes lettres. 

Il lui fut donné de consacrer une longue et saine existence aux études qu'il aimait, en la parant du goût des belles choses avec toutes les curiosités d'un esprit supérieur et désintéressé. 

Nature d'élite, au physique et au moral, il trouva dans l'indépendance de fortune et de position le loisir de développer sa forte originalité. 
Dur à lui-même, travailleur obstiné, actif et l'esprit toujours en éveil, il étudia toutes choses avec méthode et pénétration, réalisant en lui-même la vraie éducation classique. 

Cet esprit complet s'est formé aux pures sources antiques, qui lui donnèrent l'intelligence du monde des idées, et il a vécu dans l'atmosphère de paix que celle élection lui ménageait. 

Simple de mœurs, étranger aux petitesses de nos luttes intestines, il professait de fermes opinions libérales, sans avoir pour la politique cette indifférence qui, chez certains dilettanti, parait être une indifférence à la justice. Il trouva dans le réconfort de son foyer familial le calme et le charme d'une existence vouée tout entière à de chères études et à la réalisation d'aspirations élevées. 

Sa physionomie nous était familière, avec sa haute taille, ses longs cheveux et son clair regard tendu vers ses pensées. 
Il vivait dans un monde à lui, qu'il s'était créé. 

La facilité de son commerce procédait d'une profonde connaissance des hommes et de ce tranquille mépris du vulgaire, qui enseigne au sage l'indulgence envers autrui. 

La nette conscience de sa supériorité s'affirmait hautement dans les domaines qu'il s'était réservés. L'urbanité de ses relations, le charme de sa conversation intarissable de faits et d'aperçus originaux se prodiguaient à tous, par l'effet de la vivacité d'un esprit et de la chaleur d'un cœur que rien d'humain ne laissait indifférent. 
Avec cela, d'un jugement très sûr et d'un sens pratique singulièrement avisé dans les soins du matériel de la vie. 

Cette personnalité puissante et hors de pair ne pouvait être pénétrée et appréciée que dans les sphères d'élite où il fréquentait, à l'étranger et dans le petit cercle d'ici où les choses de l'esprit intéressent encore. 

Dans le domaine multiple où sa vaste intelligence s'exerça, règnent une admirable unité et la réalisation de cette harmonie ordonnée, qui est comme la marque de l'esprit antique sur la vie de notre maître. 

Flamand de race et de cœur, connaissant à fond la Hollande, il se voua à la littérature néerlandaise qui fut l'objet de son premier enseignement. 
Il vint, de longues années, se retremper aux vacances en pays de Flandre, dans sa villa de Knocke, où une œuvre de Lynen lui déroulait l'épopée d'Uylenspiegel. 
Puis ce furent les lettres françaises qu'il pénétra avec cette suite impeccable et cet enthousiasme fécond qu'il savait si bien allier. 
C'est en Hollande qu'il alla préparer son étude sur les Elzéviers, qui servirent comme de trait d'union de la Renaissance et de l'érudition bataves avec l'humanisme français. 
Les \nieme{XVI} et \nieme{XVII} siècles l'occupèrent surtout et lui donnèrent cette exacte et profonde connaissance de la langue dont il se plut à étudier les sources dans les vieux romans de chevalerie. Il en rassembla une collection unique, base de l'incomparable bibliothèque qu'il devait passer sa vie à former. 

Bibliophile érudit, il eut, avec le culte des lettres, l'amour éclairé et averti des beaux livres, des éditions originales de la noble typographie et des prestigieuses reliures de maîtres. 

Voyageur infatigable, épris des belles choses et curieux de toute civilisation, il rapporta de la connaissance des lieux et de l'observation des hommes, le goût des arts, qui en sont comme la fleur. 

Dans sa vieille demeure, d'extérieur modeste, qu'il habita toujours, il passa sa vie de savant au milieu de ses livres et de sa merveilleuse collection de tableaux des grands maîtres flamands et surtout français modernes. 

Il appartient à d'autres d'apprécier avec compétence ces hauts mérites. 

Pour celui qui ne fut que son élève ignoré, il reste à consigner en quelques traits la valeur et la portée de son enseignement et de son œuvre d'helléniste. 
C'est en somme dans l'étude de la calme et éternellement vivante antiquité, pure et maîtresse source de toute culture humaine, qu'Alphonse Willems passa la maturité de sa vie privilégiée. 
Ayant dans sa jeunesse habité Athènes, il revit la Grèce après cinquante années, avec la fidèle compagne de sa vie, et l'on peut dire qu'il vécut cette longue période, de cœur, de pensée et d'œuvres dans l'intime communion de la vieille Hellade, lui souriant toujours de son inaltérable jeunesse. 

Là est sans doute le secret de son détachement des complications de l'existence, de l'affranchissement voulu de tout ce qui pouvait contrarier l'œuvre qu'il poursuivait dans la sérénité de sa pensée. 

C'est à la littérature grecque, la plus belle expression du génie humain, qu'il finit par se consacrer tout entier. 

Il considérait son enseignement comme une mission supérieure. Pendant les vingt-huit années (1881-1909) qu'il professa le grec à l'Université de Bruxelles, il sut instruire ses élèves en les charmant, leur former surtout le jugement et le goût, en les élevant à l'intelligence des nobles pensées exprimées dans la forme parfaite.

Sans s'attacher beaucoup aux questions trop spéciales du métier philologique, ni se livrer aux développements purement littéraires, l'excellence de ses leçons procédait surtout de la saine et exacte compréhension de l'auteur qu'il expliquait. 

Œuvre de conscience, de probité littéraire, de respect du texte, une leçon de Willems était une chose vivante, complète, étudiée et achevée. 

Il préparait laborieusement chaque cours par des notes écrites, fixant le sens, comparant les sources, les interprétations ; toujours il revoyait ses notes avant la leçon. Et alors sa parole claire et animée coulait de source ; son esprit charmant, son érudition variée, sa mémoire infaillible lui fournissaient l'allusion piquante, le trait qui porte, et l'on était sous le charme jusqu'à la fin, ayant recueilli des notions sûres, la confiance dans l'autorité du maître, les révélations de choses élevées et belles. 

Le cours fini, en règle avec son auteur, le cœur satisfait, il se déridait avec ses collègues ; sa sympathie suivait l'élève. 

Pour qui apportait dans ses études autre chose que le souci de l'examen et de la place à conquérir, le maître devenait le guide le plus bienveillant, le plus précieux ; aimable, accueillant, prodiguant les trésors de sa haute culture, s'intéressant aux progrès, s'associant aux travaux de ceux en qui il espérait éveiller quelque vocation. 

Ceux qui lui ont fait honneur, devenus ses collègues ou ses successeurs, savent tout ce qu'ils lui doivent. 

Willems demeurait, par sa nature supérieure, volontiers dans le domaine idéal des conceptions qui lui étaient chères. 
Et il en arrivait à en rechercher la réalisation dans des types successifs à qui il prêtait toutes les promesses de ses propres idées. Celte générosité d'âme ne laissa pas que de lui occasionner quelques désillusions. 

Mais il y avait dans cette riche nature tant de vie, tant d'enthousiasme, une telle foi dans ses nobles buts que rien n'altérait la sérénité de ces sphères d'élection où il avait réfugié ses pensées. 

Sa dominante personnalité se marquait dans sa conversation. Avec l'élève qui l'accompagnait au sortir des cours, il développait ses idées, en abordait tous les aspects, ne tarissait pas en anecdotes amusantes, en aperçus inattendus, et l'on subissait le charme de cette causerie inimitable sans songer à y mettre du sien ; elle était si pressante et impersonnelle qu'on eût craint d'en gâter ou d'en perdre quelque chose. 

Willems inspirait à celui qu'il instruisait le respect de la science, sinon celui de certains savants, la confiance dans ses conseils, le désir de l'entendre, de retenir ses avis et le sentiment intime de tout ce qu'on gagnait en sûreté de jugement, en goût, en culture dans son commerce. 

La distance de sa supériorité à l'ignorance de l'élève se diminuait de sa bienveillance et de l'ardeur de s'instruire, de l'amour qu'il lui donnait pour ses chères études. 
C'était vraiment un maître. Son exemple et son souvenir demeurent, pour qui l'a pratiqué, une force et une consolation. 

Il semblait que ce corps de fer, cet esprit alerte demeuré si plein de jeunesse, de projets, d'intérêt pour ses travaux, dussent le mener, comme les grands penseurs de la Grèce, à l'extrême terme d'une vie privilégiée. Le malheur entra à son foyer, un mal sans merci frappa celle qu'il avait associée à toutes les heures de sa vie. Dès lors, il perdit son admirable foi dans sa destinée. 
Le sentiment de la mort lui apparut, il renonça aux longues espérances, il songea à l'œuvre qui devait lui survivre. 

Pour celui qui fut vingt ans son élève et qu'il honora d'une bienveillance non démentie, le changement fut poignant. 
Lui qui remettait toujours à plus tard la publication de son œuvre, il commença à s'en préoccuper. « Si j'étais sûr de dix-huit mois de vie », répétait-il. Et malgré tout, la divine espérance ne l'abandonna pas. Ne sortant plus depuis six mois, veillant sur sa femme, il rangeait ses livres, revoyait son Aristophane. La veille de sa mort, il regrettait de n'avoir pas devant lui une nouvelle vie, pour la consacrer à Pindare ! 

La mort l'a pris brusquement, dans toute la lucidité de son intelligence, en pleine force, comme il devait finir, debout et à l'œuvre. Elle lui a épargné les angoisses d'assister à la fin de celle qu'il laissait dans la douce inconscience de sa perte.

Et ee fut pour nous tous un vide intellectuel que la disparition du maître vers qui nous orientions notre pensée, qui vivifiait en lui et entretenait en nous le feu sacré des lettres grecques. 

On se trouve si déçu, si dépourvu à l'idée qu'il n'est plus là pour veiller à son œuvre laissée inédite ; et quand on essaie de se ressaisir pour accomplir la tâche que sa confiance nous a léguée, on se sent impuissant par le regret que lui-même n'ait pu la mener à toute la perfection que lui seul pouvait lui donner. 

Mais son œuvre capitale d'helléniste doit lui survivre et prendre la place qui lui revient. 

La traduction d'Aristophane et les mémoires académiques publiés sur chacune des onze pièces, ainsi que sur certaines questions d'ordre plus général se rattachant à l'étude du sujet, ont occupé la meilleure part des vingt dernières années d'Alph. Willems. 

Une pareille intelligence, un esprit aussi meublé, aussi en éveil sur les choses littéraires, appliqués à un tel objet devaient produire une œuvre originale et de marque. 

Il acquit la pleine connaissance de son auteur et vécut en quelque sorte ces années en communion avec la période la plus agitée de la brillante Athènes. Il s'y absorbait, s'y passionnait ; rien ne lui coûtait pour éclaircir une difficulté de texte, pour fixer une interprétation, pour dégager un point d'histoire, la portée d'une allusion, un détail de mœurs, voire une recette de cuisine, un jeu de scène ou un accessoire de costume. 

Et toujours et avant tout, les qualités maîtresses de son esprit et de son enseignement dominèrent son travail : le sérieux, le jugement, le respect du texte, l'étude profonde et consciencieuse, la défiance des autorités, le courage d'aborder les difficultés de front et de source, l'horreur de l'à peu près, du recours présomptueux à la conjecture, du pédantisme serein et satisfait, qui change ce qu'il ne comprend pas et écrit « ce n'est pas grec », parce qu'il ne sait pas dire « \italique{non liquet} ». 

L'Aristophane de Willems nous apparaît — c'est une impression personnelle — moins comme une simple translation d'un texte dans une autre langue, que comme une étude exacte et profonde du génie d'une race et d'un art supérieur, pris sur le vif, dans la plus intense manifestation de son expression littéraire. 

Ce que notre maître nous semble avoir voulu et réalisé, c'est de procurer, dans un type d'élite, la complète compréhension d'une œuvre reflétant un des plus brillants moments de la civilisation grecque. 

Il fallait pour cela tout d'abord fixer le sens et il fallait aussi faire connaître le milieu, l'esprit, les circonstances de l'œuvre. 

C'est cette préoccupation qui s'est traduite dans les notes du texte, les mémoires sur des passages controversés, les études plus générales, fixant et situant l'interprétation. 

Au fond, nous croyons que Willems pensait que qui veut se mêler de connaître Aristophane, ou Pindare ou Théocrite, n'a qu'à savoir du grec assez pour les lire et s'essayer à les comprendre dans l'original. Voici ce que lui-même\note{Note inédite, écrite au crayon et presque effacée au verso d'une convocation imprimée, datée du 3 octobre 1910 à une assemblée générale de la Compagnie du Congo pour le commerce et l'industrie. Nous la donnons telle qu'elle est, parce que c'est la seule relative à sa traduction que nous ayons retrouvée dans les papiers d'Alphonse Willems. Ce n'est évidemment qu'un premier jet, auquel l'auteur n'a pas donné sa forme définitive, peut-être en vue d'une préface à sa traduction. Nous avons cru ne pas devoir en priver le lecteur.} écrivait sur sa traduction : « On ne traduit pas un poète grec en français, et Aristophane moins qu'un autre. Car d'abord, à ne s'en tenir qu'à la langue, nulle traduction ne saurait donner une idée, ni de la richesse du vocabulaire, ni du parfait naturel, du laisser-aller, les négligences voulues qui sont comme la marque de tout ouvrier grec. Ensuite il ne faut pas oublier qu'il n'est rien de plus varié que les mètres grecs, que la forme métrique est toujours assortie à la pensée ; que, dans la comédie surtout, le système se modifie suivant le changement d'idée et de ton, c'est-à-dire avec presque, chaque scène.  

» Tout ce qu'il est permis d'essayer, c'est de rendre avec le plus de fidélité possible la suite des pensées et des images. C'est à quoi nous avons visé, mais sans nous faire d'illusion sur le résultat.

» L'office d'un traducteur est pareil à celui des embaumeurs de l'ancienne Égypte. On sait ce que sont devenues dans leurs mains les plus fières comme les plus délicates beautés du temps des Pharaons : de sèches momies, laissant deviner les lignes principales du corps, mais où toute couleur, toute apparence de vie a disparu. De même que dans la fleur flétrie un botaniste exercé retrouve l'élégance du type et la richesse du coloris qui la caractérisaient autrefois, de même on en vient à pouvoir reconstituer dans son imagination, devant une momie desséchée, l'être vivant qu'elle fut jadis. La comparaison est de M\expo{me} Dacier, et je n'en connais pas de plus juste. 

» Il en est qui prétendent que c'est en vers qu'il faut traduire les poètes grecs. Je ne sais comment il en va dans d'autres langues, mais en français, il faut nettement le déclarer, l'entreprise est impossible. Elle a été déclarée telle par M\expo{me} Dacier et par Leconte de Lisle, un juge dont nul ne, contestera la compétence ; et de fait, à part peut-être quelques épigrammes de l'Anthologie, on ne citera pas vingt vers de suite traduits du grec qui soient autre chose que des ingénieux à peu près ou des paraphrases plus ou moins heureuses. 

» La besogne est déjà assez compliquée de suivre rigoureusement le texte sans rien ajouter, sans rien escamoter, en se rendant un compte exact de la valeur de chaque mot et, chose capitale, en se gardant de toute feinte de langage. Car si Aristophane se trouve être le plus débridé des écrivains attiques, il en est aussi le plus pur, le plus châtié, car à cet égard il l'emporte même sur Platon\note{Sans doute, pour l'impression, Willems eût écrit : « Il ne peut se comparer qu'à Platon ».}. Seul un Racine, le Racine des \italique{Plaideurs}, semblerait digne d'être son interprète. Malheureusement, c'est la méthode contraire qui semble prévaloir aujourd'hui. 

» Je ne sache rien de plus chagrinant pour ma part que de voir celui dont le philosophe disait que les Charités avaient choisi son âme pour leur servir de sanctuaire, rendu dans le style de Labiche ou pis encore dans l'ignoble patois des concierges et des cochers de fiacre. J'avoue que les plus grossiers contresens me paraissent faute vénielle au prix » ... 

Et cependant Willems a consacré à cette œuvre de longues années d'un labeur inlassable. Pourquoi et pour qui ? Nous ne pensons pas qu'il se souciât de vulgariser Aristophane dont le public n'a cure. Il a fait sa traduction pour lui, pour le plaisir de la faire, afin d'acquérir la complète compréhension de son auteur, d'en fixer le sens, d'en noter les subtiles nuances, de le dégager des conjectures et des tares qui en altéraient la pureté native. 

Sa traduction est une justification du texte, du sens et de l'esprit de l'incomparable comique. Elle est la base et la référence d'une grande œuvre critique dont le développement se trouve dans les notes sur les passages controversés, incompris ou faussés, dans les études sur le milieu et les circonstances qui ont produit l'ancienne comédie attique. 

Œuvre exacte, fidèle et complète, d'une probité littéraire singulière, la traduction française, dans sa correction de pureté élégante « se gardant de toute feinte de langage », pourra surprendre le lecteur ordinaire. 

Il arrivera que le texte lui paraisse froid, d'une précision minutieuse, peu coulant, sans laisser-aller, ne faisant grâce de rien, sacrifiant trop peu à l'effet et ne s'adaptant pas assez à notre façon de penser et d'imager l'idée. Et c'est la justement ce qui en fait la haute valeur, la qualité maîtresse. Willems n'a jamais songé à moderniser Aristophane, à accommoder son poète à notre façon ; la rectitude de son jugement et la sûreté de son goût l'ont gardé de fausser ainsi un esprit, un milieu, des mœurs qui ne sont plus les nôtres et dont il faut se pénétrer avec lui pour en apprécier le vivant tableau. 

Pour juger son travail, il faut le lire avec l'original, et c'est alors qu'il se révélera dans toute son admirable conscience, dans la justesse de sens, l'adéquate précision des termes, l'élégance et le bonheur d'expression, et que l'helléniste, charmé de pareille résurrection, se dira : c'est bien cela, maintenant j'ai compris. 

Nous croyons en effet que, mieux que quiconque, notre maître a compris Aristophane et a justifié son interprétation. 

Ses notes critiques, où il n'est plus engagé dans la lutte inégale de rendre en prose française le plus irréductible des poètes, laissaient le champ plus libre à son érudition, à son esprit sagace, à la piquante urbanité de sa verve d'humaniste. Et cette partie de son œuvre, toute personnelle et nouvelle, en est peut être la plus curieuse, la plus sérieusement instructive et la plus caractéristique. 

Ses mémoires académiques sur Aristophane et son temps sont des modèles de jugement, de goût, de critique désintéressée et pénétrante, qui sont appelés à rénover et à fixer l'interprétation du poète. 

Répétons-le, c'est là une œuvre mûrie, fortement élaborée et étayée, qui vaut par elle-même et qui prendra le rang qui lui revient dans l'interprétation et la critique du grand poète d'Athènes. 

Notre maître aimé a bien mérité d'Aristophane et s'en est allé en paix dans cette consolante pensée. \\


Il nous reste à dire quelques mots de l'édition. 

Le manuscrit des onze pièces est de la nette et belle écriture de l'auteur sur de petits papiers à lettres, avec les notes en regard, un certain nombre au crayon. 

Ce texte, d'une correction parfaite, a été scrupuleusement respecté, sans changer, ajouter ou supprimer un mot. \\

\hspace*{1cm}27 novembre 1913.\hspace*{7cm} Ch. L.







\Note{Notes sur les \italique{Thesmophories} d'Alphonse Willems}\note{Extrait des \italique{Bulletins de l'Académie royale de Belgique} (Classe des lettres, etc., n° 11, novembre 1908).}
En publiant ces notes sur les \italique{Thesmophories}, je ne puis que répéter, une fois de plus, que mon intention n'est pas de faire un commentaire perpétuel. 

Je ne touche que certains passages, mal interprétés, à mon avis, et surtout corrigés mal à propos par la généralité des éditeurs. 

Cela ne veut pas dire que je passe condamnation sur le reste et notamment sur les innombrables conjectures introduites dans le texte depuis ces dernières années. 
Bien au contraire. Mais comme elles ont peu de chance, du moins je l'espère, d'être adoptées, ce serait perdre son temps que de les discuter. 

Ainsi, au v. 60, χοανεύειν c'est \italique{jeter en moule}, et pour comprendre la tournure il suffit d'ouvrir un lexique, au mot κατά suivi d'un génitif, \italique{desuper in}, \italique{nieder ein} ; car c'est bien ainsi qu'on emplit un moule. (Cf. Démosth., 403, 4 : ξαίνειν πληγὰς κατὰ τοῦ νώτου.) Quant à θριγκός, mis à dessein au singulier, inutile ici de l'expliquer, tant pis pour ceux qui ne l'entendent pas\note{Voir ma note sur les \italique{Nuées}, v. 676.}. M. van Leeuwen, qui adopte la conjecture de M. Blaydes, εἰς τοὺς θριγκούς , aurait bien dû nous dire en note, non ce qu'elle signifie, je ne demande pas l'impossible, mais ce qu'elle est censée signifier. 

Je voudrais également qu'on expliquât en vertu de quelle raison δικάζειν, v. 79, doit être remplacé par καθίζειν ; τεχνάσμασι,, v. 198, par τεχνήμασι ; δρῶσαν, v. 306, par ὁρῶσαν (cf. \italique{Lysistr}., v. 1047) ; ὑπό του, v. 493, par ὅλην ; χέσειν, v. 570, par χέσαι, contre le témoignage des grammairiens anciens ; οὐ πάρεισιν αἱ, v. 771, par οὐ γὰρ εἰσί μοι ; βέλτιον, v. 774, par βέλτιστον ; pour quoi il y a lieu de mettre οὐδὲ θἔμισυ, v. 452, à la place de οὐδ᾽ εἰς ἥμισυ ; notez que la même distinction existe en français entre : un commerce réduit \italique{de la moitié} ou \italique{de près de moitié}, ce dernier sens étant ici bien mieux approprié\note{On dit εἰς ἥμισυ, sans l'article, comme ὑπὲρ ἥμισυ (cf. Kühner-Gerth, §462 1). Quant à la nuance exprimée par εἰς, cf. Thucyd., VII, 30, 3, εἰς εἴκοσι ἱππέας, une vingtaine de cavaliers; VII, 32, 2, εἰς ὀκτακοσίους , près de huit cents hommes.}. 

Au v. 314, il est d'un intérêt supérieur, j'en conviens, de savoir que M. Blaydes préférerait χαρέντας ἐπιφανῆναι à φανέντας ἐπιχαρῆναι. 
Mais est-ce une raison pour corriger Aristophane ? Sans doute on nous fait valoir que \italique{proprie de divina prœsentia usurpatur vox} ἐπιφανεία (pour moi, je n'en avais jamais douté, et au besoin l'almanach me l'aurait appris). 
Seulement c'est à Aristophane qu'il appartenait de s'en souvenir; et voyez, c'est comme un fait exprès, nulle part il n'emploie le mot propre, alors qu'il use et abuse sans vergogne du simple φάνηθι, φάνητε, etc.\note{Cf. v. 1144 ; \italique{Caval.}, 591 ; \italique{Nuées}, 266 ; \italique{Paix}, 638, etc.}. 

Qu'y a-t-il à redire à δύ᾽ ἢ τρία κακὰ εἶπε δρώσας [ἡμᾶς] μυρία (v. 475) ? Est-il rien de plus régulier que εἰπεῖν τινά τι, comme au v. 552 ? Littéralement : « Sachant deux ou trois de nos méfaits il les a débités, sur nous qui en commettons d'innombrables. » 

Au v. 558, pourquoi, au lieu de τὰ κρέ᾽ ἐξ, faut-il lire κρέ᾽ ἐξ, ce qui rend le vers dur et déplaisant ? Serait-ce, car on s'y perd, pour une raison de métrique ?
Mais l'anapeste du second pied n'est pas plus singulier que, par exemple, dans \italique{Nuées}, 187, et \italique{Ecclésie}, 394. 

Pourquoi, au v. 772, M. van Leeuwen change γένοιτο, leçon du \italique{Ravennas}, en γένοιντο, m'est un problème. 
Avec γίγνεσθαι, le singulier pour le pluriel passait jusqu'ici pour une élégance. On la rencontre plus d'une fois dans Platon (\italique{Banq.}, 188b, \italique{Rép.}, 363a, etc.). 
En français on dirait également : « D'où m'arrivera-t-il des avirons ? » Il en va de même avec εἶναι. Cf. ci-dessus le v. 21 : οἷόν γέ πού 'στιν αἱ σοφαὶ ξυνουσίαι ; \italique{Guêpes}, 58; Soph., \italique{Trach.}, 520; Eur., \italique{Ion}, 1146 ; Plat., \italique{Euthyd}., 302c, etc. 

Je cite sans ordre, au hasard des coups de crayon mis en marge de mon exemplaire. Mais en voilà assez, et je m'arrête. 
Eh bien ! je dis que ce qu'on nous donne à lire n'est plus Aristophane, mais ce qu'autrefois l'on aurait appelé Aristophane travesti. 

Les grands savants du \nieme{XVI} siècle qui, révérence parler, ne le cédaient guère aux nôtres en fait d'érudition et d'esprit critique, se seraient fait un cas de conscience d'introduire dans le texte une leçon qui ne fût pas de tout point justifiée\note{J. Scaliger, parlant de ces insupportables maniaques (car il y en avait de son temps) qui trouvent des fautes partout et corrigent à tort et à travers, disait : \italique{Ego vero nunquam id audendum putaverim, nisi sit error manifestissimus, satiusque judico, ubi aliquid displicet aut ignoro, priorem lectionem retinere (mea tamen sententia annotata) et doctioribus ejus loci correctionem relinquere, qui forte ad id mea admonitione excitabuntur, quam temere vel syllabulam ullam mutare. Et si id tam religiose observassent importuni isti correctores, non haberemus hodie auctores veteres adeo depravatos ; nam hac se ingerente licentia impestive et inconsulto corrigendi quod displicet, totidem varietates locorum prodeunt et prodierunt, quot editiones : et nullos magis corruptos reperias libros, quam qui sœpius excusi sunt et a pluribus perlecti}. (\italique{Scaligerana prima}, Amst., 1740, p. 76.) Qu'il serait à souhaiter que cette déclaration d'un des plus savants hommes qui furent jamais servît de règle à tous les éditeurs !}. Aujourd'hui l'on ne se borne pas à amender tant bien que mal ce qui est ou paraît fautif, on s'applique à dénaturer, à tort et à travers, tout ce qui est susceptible de subir un changement, sans alléguer d'autre raison que l'éternel \italique{displicet} ou \italique{malim} ou \italique{valde arridet}, qu'on a la naïveté d'espérer nous faire passer pour un argument. 

Et comme la plupart du temps c'est le poète même qui est mis en cause, à peine citerez-vous dix vers de suite où il ne reçoive des férules de ceux qui ont péniblement appris de lui à déchiffrer du grec\note{Il y a longtemps que Boissonade disait : \italique{Puto in talibus nonnunquam errores auctorum ipsorum a criticis corrigi, non librariorum}. Sur \italique{Acharn.}, 246.}. Surtout, si l'envie vous prend de relire votre auteur, ou simplement de vérifier un passage, gardez-vous de quitter de l'œil un instant ce qu'on appelle l'annotation critique, sous peine de confondre sans cesse le texte authentique avec ce qui provient de la fantaisie ou, plus souvent encore, de l'impéritie de l'éditeur.

\Note{Note 1 aux \italique{Thesmophories} : vers 16-18}

\citComm{\decalage{ὀφθαλμὸν ἀντί}πρῶτ᾽ ἐμηχανήσατο\\
ὀφθαλμὸν ἀντίμιμον ἡλίου τροχῷ,\\
ἀκοὴν δὲ χοάνης ὦτα διετετρήνατο.\\}


La conjecture de Reiske : δίκην δὲ χοάνης, adoptée par les éditeurs, et celle de Boissonade : ἀκοῇ δὲ χοάνην, donnent un sens plausible, mais sont-elles nécessaires ? Je ne puis admettre pour ma part qu'on ait recours aux conjectures, lorsqu'un texte est susceptible d'interprétation. Or il est inexact de prétendre, comme on l'a fait, que la leçon des manuscrits n'offre aucun sens.

Ἀκοή, dans l'acception d'organe de l'ouïe, \italique{oreille}, se lit maintes fois dans Euripide (\italique{Phén.}, 1480 ; \italique{Iphig. T}., 1496, et passim). 
C'est un terme plus relevé pour οὗς, comme à ὄψις l'est pour ὄμμα, et βάσις pour ποῦς\note{Etym. Magnum : ἀκοή λέγεται καὶ τὸ οὖς ὅλον --- Grammat. Bekkeri Anecd., p. 369, 1 : ἀκοήν, ἀντὶ τοῦ οὖς.}. 
Après ὀφθαλμός, au singulier, ἀκοή se prête mieux au rapprochement que le pluriel ὦτα. 

Χοάνη est défini par Hésychius : τύπος εἰς ὃν μεταχεῖται τὸ χωνευόμενον, « le moule dans lequel on coule la fonte » ; de là le verbe χοανεύειν, qui se lit plus loin (vv. 57 et 62), pour εἰς χοάνην ou κατὰ χοάνης ἐμβάλλειν, « jeter enmoule ». Dès lors la signification de ὦτα ne fait plus doute : ce sont évidemment les \italique{godets}, ou entonnoirs servant à introduire le métal fondu. Le mot me paraît même tout à fait approprié\note{Comparez le français \italique{Orillon}. C'est ainsi qu'on a donné le nom de ὦτα, \italique{oreillettes}, à cause de leur forme, à certains appendices du cœur (lire la définition de Galien).}. 

Le scholiaste, fort sobre de commentaires sur les \italique{Thesmophories}, nous laisse dans l'ignorance à propos de ὦτα et se borne à noter : λείπει ὡς. 
Tenant compte de cette indication, nous construisons : διετετρήνατο δὲ ἀκοὴν ὡς χοάνης ὦτα. « L'Éther fabriqua pour l'homme un œil calqué sur le disque du soleil, et fora une oreille à l'image des godets d'un moule\note{Cf. \italique{Acharn.}, 230, σχοῖνος ἀντεμπαγῶ ; \italique{Nuées}, 178, διαβήτην λαβών ; \italique{Lysist}., 695, κάνθαρός σε μαιεύσομαι ; \italique{Gren.}, 159, ὄνος ἄγων μυστήρια ; \italique{Plout.}, 295, τράγοι δ᾽ ἀκρατιεῖσθε, etc.} ». 

La comparaison est ambitieuse et puérile, mais elle ne manque pas de justesse. C'est à dessein qu'elle est mise dans la bouche d'Euripide cherchant à éblouir son niais de beau-père.
À remarquer d'ailleurs que les Grecs de cette époque étaient très familiarisés avec la sculpture en bronze. Une allusion à la même industrie se lit quelques vers plus loin.

\Note{Note 2 aux \italique{Thesmophories} : vers 21}

\citComm{οἷόν γέ πού ᾽στιν αἱ σοφαὶ ξυνουσίαι\\}

Telle est la leçon du manuscrit, confirmée par le scholiaste. 
Elle est excellente, tandis que les conjectures de Brunck, Fritzsche et M. van Leeuwen ne supportent pas l'examen. 

M. van Leeuwen fait remarquer qu'avec οἷον exclamatif που n'est pas de mise. 
C'est bien possible, mais il s'agit, non de που, mais de γέ που, ce qui est tout différent.
Ainsi associées, ces deux particules ont un sens très précis, sur lequel, il est vrai, on s'est souvent mépris. 
Elles se traduisent par \italique{utique}, \italique{quidquid est} ; en français : en tous cas, de toute façon, de compte fait, somme toute. 
De sorte que οἷόν γέ που, ou οἷόν γέ τοί που, ne diffèrent pas très sensiblement de οἷον ἆρα. 

Dans les tragiques, je ne trouve à citer que le v. 533 de l'\italique{Ajax} de Sophocle : 
\citComm{μὴ σοί γέ που δύστηνος ἀντήσας θάνοι.}


Mais les exemples abondent dans Platon, entre autres \italique{Banq}., 205e, ἐὰν μὴ τυγχάνῃ γὲ που ἀγαθὸν ὄν ; \italique{Théét}., 147bc, πρῶτον μέν γέ που ... ἔπειτά γέ που ; \italique{Polit.}, 259e, παρεχόμενος γέ που ; \italique{Rép.}, V, 478a, ἐπιστήμη γέ που ; \italique{Ibid}., X, 607d, δοῖμεν δέ γέ που ; \italique{Hipp. Maj}., 298a, οἵ τέ γέ που καλοὶ ἄνθρωποι. 

La même locution se lit dans Aristophane, \italique{Guêpes}, 27 : 
\citComm{δεινόν γέ που ᾿στ᾽ ἄνθρωπος ἀποβαλὼν ὅπλα\\}

« C'est une cruelle chose en somme qu'un homme qui a jeté ses armes » ; où les éditeurs, depuis Dindorf, ont eu tort de remplacer cette leçon, qui est celle de tous les manuscrits et du scholiaste, par γε τοὔστ᾽ (τοι ἔστιν). 

Aussi dans les \italique{Acharniens}, 896 : 
\citComm{ἀγορᾶς τέλος ταύτην γέ που δώσεις ἐμοί\\}

« (Cette anguille) tu me la donneras de toute façon pour droit de marché. »

Et dans la \italique{Paix}, 1047 : 3
\citComm{ἀλλ᾽ Ἰεροκλέης οὗτος γέ που ᾽στιν\\}

« Mais c'est bien Hiéroclès après tout... » 

Comme par un fait exprès, le vers qui fait l'objet de cette note se trouve certifié par un vers tout pareil d'Eupolis, cité par Eustathe et confirmé par le scholiaste d'Homère, \italique{Il.}, II, 333 : 

\citComm{οἷόν γὲ πού ᾽στι γλῶσσα κἀνθρώπου λόγος.\\}

« Quelle fameuse chose en somme que la langue et le langage humain ! »%

\Note{Note 3 aux \italique{Thesmophories} : vers 22-24}%

\citComm{\textbf{Εὐρ.} --- πόλλ᾽ ἂν μάθοις τοιαῦτα παρ᾽ ἐμοῦ.}
\citComm{\decalage{πόλλ᾽ ἂν μάθοις τοιαῦτα παρ᾽ ἐμοῦ.}\textbf{Μνη.} --- πῶς ἂν οὖν}
\citComm{\decalage{\textbf{Εὐρ.} --- }πρὸς τοῖς ἀγαθοῖς τούτοισιν ἐξεύροιμ᾽ ὅπως}
\citComm{\decalage{\textbf{Εὐρ.} --- }ἔτι προσμάθοι μὴ χωλὸς εἶναι τὼ σκέλει;\\}

Ces deux derniers vers n'ont pas de sens. Et sans doute les conjectures ne manquent pas, mais la phrase n'en reste pas moins embarrassée. 
Aussi M. van Leeuwen s'est-il décidé à la refaire ; moyen toujours facile d'esquiver une difficulté, mais que par cela même on devrait s'interdire, même dans les cas désespérés. 

Le remède, selon moi, consiste à supprimer le v. 23, qui n'est qu'une glose explicative du suivant, et de lire προσμάθοιμι, au lieu de προσμάθοι μὴ, ce qui peut à peine s'appeler une correction, la confusion de η et de ι s'expliquant par l'iotacisme. L'interpolation me paraît évidente : προσμάθοιμι, c'est-à-dire πρὸς τοῖς ἀγαθοῖς τούτοισιν [ἐξεύροιμι]. 
Lisons donc : 
\trimIamb{\decalage{ἔτι ποοσμάθοιμι χωλὸς }πῶς ἂν οὖν\\
ἔτι ποοσμάθοιμι χωλὸς εἶναι τὼ σχέλη.}\\
EURIPIDE : Avec moi tu apprendras bien des choses du même genre.\\
MNÉSIL. : Ah! que ne puis-je apprendre en outre à être perclus des deux jambe.

Ce que le scholiaste explique fort bien par : « pour n'être pas forcé de suivre ». 
Ce vœu, Mnésiloque peut l'adresser en toute confiance au « faiseur de boiteux », χωλοποιός , qu'était Euripide. ll se justifie d'ailleurs par un passage analogue de la \italique{Paix}. Polémos, qui a entrepris d'écraser dans un mortier les cités grecques, envoie un de ses acolytes à la recherche d'un pilon. Et Trygée, témoin impuissant de cette scène, exprime, lui aussi, le souhait que le messager « se donne une entorse à chaque pied » 

\citComm{\decalage{ἀποστραφῆναι}ἀποστραφῆναι τοῦ μετιόντος τῶ πόδε (v. 279.)}


\clearpage

\Note{Papiri della Società Italiana n° 1194 : vers 145-156 ; 237-245 ; 273-288 ; 594-596 ; 804-809 ; 1231 (et 133-144 \italique{PSI} XIV, p. XV.)}


\begin{figure}[!h]
	\hspace*{-1.5cm}
\centering
\includegraphics[scale=0.55]{PSI1194.jpg}
\caption{\italique{Papiri della Società Italiana} n° 1194 : vers 145-156; 237-245; 273-288; 594-596; 804-809}
\label{POxy1194}
\end{figure}

\Note{Note 4 aux \italique{Thesmophories} : vers 217}

\citComm{ἢ μὴ διδόναι. γ᾽ἐμαυτὸν ὥφελον πότε.}\\


Qu'y a-t-il à redire ? Διδόναι ἑαυτόν τινι est la formule ordinaire, et qui se lit partout : \italique{Cav.}, 740 ; Soph., \italique{Phil.}, 84 ; \italique{Trach.}, 1117 ; Hérod., VI, 108 ; Thucyd., II, 68 ; Démosth., \italique{P. Cor.}, 503d. 
En latin, Térence : \italique{Da te hodie mihi} (\italique{Adelphi}, 5, 3). 
D'autre part, rien n'est mieux justifié que γε. Cf. \italique{Cav}., 413, ἢ μάτην γε... 
Dawes corrige : \\
\citComm{ἢ μὴ ᾿πιδοῦν ἐμαυτὸν ὥφελον πότε,}\\

ce que J. d'Orville, qui était homme de goût en même temps qu'un savant, appelait \italique{mollissimum et novis elegantiis commendabilem versum}\note{\italique{Ad Charitonem}, p. 598.}. Porson, lui, n'est pas de cet avis. Qui a tort ? Je n'en sais rien, et cela m'est égal. Il me suffit de constater que ce vers n'a de raison d'être que parce que le poète, en deux autres endroits (213 et 249), a usé, non du simple, mais du composé ἐπιδιδόναι. Ne disais-je pas naguère que si pour votre malheur vous avez écrit deux fois \italique{autour}, ces grammairiens-là vous interdiront d'écrire désormais \italique{alentour} ? 

Que si vous objectez qu'Aristophane s'est servi ailleurs d'un autre composé, παραδιδόναι\note{\italique{Nuées}, 436 ; \italique{Guêpes}, 1130.}, on vous répondra que c'était son droit dans les \italique{Nuées} et les \italique{Guêpes}, mais non plus dans les \italique{Thesmophories}. Autre exemple :

\Note{Note 5 aux \italique{Thesmophories} : vers 230}

\citComm{ἔχ᾽ ἀτρέμα σαυτὸν κἀνάκυπτε.}\\

C'est la leçon du \italique{Ravennas}. M. van Leeuwen, d'après Dobree, corrige : ἔχ᾽ ἀτρέμα αὐτοῦ, \italique{quia ferri nequit accusativus} ; lequel αὐτοῦ était à sa place \italique{Ois}., 1200, mais pas du tout ici. 
Hérodote a pourtant dit, en parlant de l'armée athénienne : εἶχον ἀτρέμας σφέας αὐτοὺς ἵνα ἐτάχθησαν\note{Hérodote, IX, 54 (aussi III, 79). — De même Platon, \italique{Epist.}, 348 B, écrit : ἐκποδὼν ἑαυτὸν ἔσχεν et Eurip., \italique{Iph. T.}, 1129 : ἐκποδὼν τοῦδ᾽ ἐχειν μιάσματος.}. Hérodote le pouvait, mais non Aristophane, qui ailleurs s'est borné à ἔχ᾽ ἀτρέμα\note{ll est vrai qu'une des quatre fois les bons manuscrits portent ἔχ᾽ ἀτρεμί (\italique{Nuées}, 261). Était-ce assez dégoûtant ? Heureusement que M. van Leeuwen y a mis bon ordre.}. Quoi ! voilà un comique qui, à quatre reprises (car on s'est donné la peine de compter), a écrit : \italique{reste coi}, et la cinquième fois il se permettrait d'écrire : \italique{tiens-toi coi}. À quoi servirait la statistique si l'on tolérait de pareilles licences ?

\Note{Note 6 aux \italique{Thesmophories} : vers 263}

Il faut lire γ᾽ οὐ, et non γοῦν. Mnésiloque, à qui Agathon veut prêter des chaussures, lui demande : « M'iront-elles? car tu n'aimes pas à te chausser large ». Agathon, en effet, nous est donné pour un efféminé, affectant le bel air, et se vantant lui-même d'être recherché dans sa mise (v. 172). 
Il n'aurait eu garde de porter des souliers larges, car c'eût été faire preuve d'inélégance et même de rusticité. (Si telle chose était à démontrer, on pourrait citer là-dessus Théophraste, au chapitre IV de ses \italique{Caractères}.) 
Mnésiloque, en vieillard qui aime ses aises, est donc fondé à craindre que les souliers qu'on lui présente ne soient trop étroits pour lui. Tel est le sens de ce vers. Quoi qu'en dise le magister byzantin, il n'y a pas lieu de raffiner, les termes en grec excluant toute idée d'obscénité. 
Mais c'est le malheur d'Aristophane qu'il ne peut parler ni de figues, ni de myrtes, ni de cochons de lait, ni de grains d'orge, ni d'orne de vignoble (ὄρχος), ni même de moucherons (\italique{Lysist}., 1031) ou de chaussures étroites ou larges, ni surtout du mot ἐραστής suivi d'un génitif masculin\note{Voir la note de M. van Leeuwen sur \italique{Oiseaux}, 324; encore qu'il cite ailleurs (\italique{Cav}., 732) Périclès usant de la même tournure dans sa célèbre oraison funèbre. Mais Périclès n'est pas, comme Aristophane, un pornographe.}, qu'aussitôt certains ne le soupçonnent de gravelure. 

Au surplus, le poète aurait eu l'intention qu'on lui prête, qu'alors surtout la négation devenait indispensable. Si vous ne devinez pas pourquoi, cherchez. Mais ne vous en rapportez pas à la note de M. van Leeuwen. Elle est d'une naïveté qui désarme, et ne peut être comparée qu'à celle qu'il a faite sur un passage de l'\italique{Ecclésie} (v. 991) ayant précisément trait au même sujet.

\Note{Note 7 aux \italique{Thesmophories} : vers 280-281}

\citComm{ὦ Θρᾷττα, θέασαι, καομένων τῶν λαμπάδων\\
ὅσον τὸ χρημ᾽ ἀνέρχεθ᾽ ὑπὸ τῆς λιγνύος.}\\


Les derniers éditeurs ont accepté une soi-disant correction de Reiske : ἀνέρχεται τῆς λιγνύος. Sur quoi se fondent-ils pour dénaturer de la sorte, non plus seulement le texte, mais la pensée même d'Aristophane ? 
Serait-ce sur une faute de langue ? Mais il n'est pas de tournure plus purement grecque que celle qu'on prétend amender. Faut-il citer Photius : οἱ παλαιοὶ τὴν ὑπό ἀντὶ τῆς μετά πολλάκις ἐλάμβανον ? 
Il n'y a qu'à ouvrir le dictionnaire pour trouver des exemples. Bornons-nous à ce vers d'Homère qui rend si bien la même pensée que l'on peut se demander si l'auteur ne l'avait pas présent à l'esprit : \\

\citComm{νύμφας δ᾽ἐκ θαλάμων δαΐδων ὕπο λαμπομενάων\\
ἠγίνεον ἀνὰ ἄστυ}\note{\italique{Iliade}, XVIII, 492. Aussi \italique{Odyssée}, XXIII, 290.}

Voyons maintenant ce que le sens a gagné à cette correction. M. van Leeuwen suppose qu'au moment précis où finit le prologue (v. 279), le fond de la scène s'ouvre, découvrant l'intérieur du Thesmophorion. Là sont assises, tenant chacune une torche à la main, les femmes qui constituent le chœur. Ainsi s'explique, d'après lui, l'exclamation de Mnésiloque : \italique{Incensœ postquam sunt faces, vide quantus splendor inde surgat !}
Malheureusement, même amendé par Reiske, le texte ne peut nullement signifier cela. 
Λιγνύς, comme le dit encore Photius, est une exhalaison fumeuse, ou, si l'on préfère, un dégagement de fumée, καπνώδης ἀναφορά. Je veux bien que, suivant l'intensité du foyer, cette sorte de nuage soit tantôt vivement coloré, στέροψ (\italique{Antig.}, 1127), tantôt purement fuligineux, μέλαινα λιγνύς (\italique{Sept. c. Th}., 494) ; mais, clair ou obscur, jamais il n'a pu évoquer l'idée de splendeur. Il faudrait traduire : « Quelle effroyable fumée s'exhale de ces torches brûlantes ! » Réflexion qui, j'en conviens, ne se justifierait que trop, si la scène se passait en effet comme on l'imagine. 

Songez donc que le chœur représente, non pas deux douzaines de femmes, mais toute la population féminine d'Athènes, censée réunie dans ce sanctuaire, auquel il faudrait nécessairement prêter les dimensions d'une cathédrale. Peuvent-elles faire autrement, ces femmes avec leurs torches, que d'empester tout l'édifice et d'enfumer les trésors, trépieds, vases d'or, ivoires, tapisseries, etc., qui, suivant l'usage, y étaient conservés ? 

Ne confondons pas le théâtre antique avec l'Opéra. S'il est un point bien établi, c'est que jamais, ni chez les comiques, ni non plus dans la tragédie, l'action n'a lieu dans un édifice couvert ; invariablement elle se passe en plein air. De même, le changement à vue consistant à faire s'ouvrir le décor de fond sur les profondeurs d'un palais ou d'un temple est un artifice moderne, dont Racine, par exemple, a tiré un admirable parti dans \italique{Athalie} ; mais il était ignoré des anciens, et même ils en auraient eu l'idée qu'ils n'auraient pu la réaliser. Il suffit, pour s'en convaincre, de jeter un coup d'œil sur le plan de n'importe quel théâtre antique.

Suivant M. van Leeuwen, les femmes, durant le prologue, se sont rendues par un escalier dérobé dans l'enceinte du temple. Lorsque Mnésiloque a fini sa prière, il va s'asseoir parmi elles. De sorte que le long épisode qui va suivre (du v. 279 au v. 659, près du tiers de la pièce) se passe tout entier dans le temple, tandis que la scène et l'orchestre restent vides. Et voilà le plus clair résultat de la malencontreuse conjecture de Reiske. 

Les \italique{Thesmophories} n'ont pas à proprement parler de parodos ou entrée du chœur. À cet égard elles peuvent être comparées à l'\italique{Ecclésie des femmes}. De même que, dans l'\italique{Ecclésie}, les femmes arrivent peu à peu au signal donné par Praxagora au moyen d'une lampe, de même ici elles ont attendu, comme Mnésiloque, le signal hissé sur le Thesmophorion (v. 278). Arrivé le premier dans l'orchestre, Mnésiloque s'écrie envoyant arriver le cortège : « Vois, Thratta, sous la lueur fumeuse des torches ardentes, vois la foule qui monte ». 
Après quoi il fait son offrande et adresse aux déesses son invocation. Quand il a fini, les femmes se trouvent réunies dans l'enceinte sacrée et la prêtresse commence les prières d'usage.

\Note{Note 8 aux \italique{Thesmophories} : vers 289-291}

\citComm{καὶ τὴν θυγατέρα χοῖρον ἀνδρός μοι τυχεῖν\\
πλουτοῦντος, ἄλλως τ᾽ ἠλιθίου κἀβελτέρου,\\
καὶ πρὸς θάληκον νοῦν ἔχειν μοι καὶ φρένας.}\\

Telle est la leçon du \italique{Ravennas}. Mais le premier et le troisième vers ne sont pas grecs, et l'altération est si profonde qu'il est à craindre qu'on ne parvienne jamais à les restituer avec certitude. Heureusement le sens ne paraît pas douteux. Mnésiloque, suivant l'usage, fait des vœux pour ses enfants\note{Les Thesmophories étaient destinées surtout à adorer Démêter comme présidant au mariage et à la production des enfants, assimilés aux fruits de la terre. La déesse y était invoquée comme Calligénie, « à la belle progéniture ».} : il souhaite à sa fille un mari riche, mais benêt, à son fils un esprit sensé et judicieux. 

Il faut commencer par faire justice de l'ancienne conjecture de Biset, πρὸς φάλητα (au lieu de πρὸς θάληκον), reprise sous une autre forme par Fritzsche et M. van Leeuwen : πρὸς τὸ ληκᾶν ou ληκεῖν. 
Tranchons le mot, elles sont plus qu'ineptes, elles sont révoltantes. À qui fera-t-on admettre qu'il ait existé de par le monde un père ou une mère assez dépravés pour souhaiter pareille chose à leur fille ? Notez de plus que Mnésiloque s'adresse à deux déesses, et que dès lors, si inculte, si cynique même qu'on le suppose, il n'aura garde de faire étalage d'obscénités. Car, non moins que Pallas, les déesses d'Éleusis étaient entourées de respect. Jamais poète athénien ne se serait permis à leur endroit, je ne dis pas d'immondes plaisanteries, je dis la moindre irrévérence. 
Le public d'ailleurs ne l'aurait pas souffert. 

Cette considération suffit pour écarter une autre conjecture, très goûtée des anciens éditeurs : τὸν θυγατέρος χοῖρον, ou τὸ θυγατρὸς χοιρίδιον (G. Hermann), une simple saloperie doublée d'un non-sens. Pour qui prend-on Aristophane ? N'insistons pas. Il est de ces fautes de goût tellement exorbitantes, qu'il faut renoncer à les discuter, crainte de s'exposer soi-même à manquer à la mesure. 

Pour revenir à notre texte, il est de la dernière évidence que τὴν θυγατέραν est une glose faisant double emploi avec χοῖρον. 
La preuve c'est que le scholiaste dit expressément que ce dernier mot sert à désigner la fille de Mnésiloque, et en même temps il nous apprend comment il était orthographié, puisqu'il le compare à βοίδιον, χρυσίον et μύρτιον. 
Il faut donc lire τὸ χοιρίον, car l'article ne peut être omis. 
Nous pouvons maintenant restituer le vers presque entier : καὶ τὸ χοιρῖον ἀνδρὸς μοι τυχεῖν. Il ne reste à compléter que le premier mètre. Suppléer par conjecture un texte gâté par une glose est une besogne facile, mais forcément sujette à caution. 
En pareil cas la leçon la plus simple est toujours préférable. 
On pourrait, par exemple, proposer : καὶ συζύγου τὸ χοιρίον... Mais qu'y a-t-il de plus simple que de lire : ἔπειτα καὶ, comme dans \italique{Ploutos}, 796 ; \italique{Œd. Col.}, 1300, etc. ? 

Nous passons au troisième vers. Parmi les corrections proposées, la plus accréditée est celle de G. Dindorf : καὶ ποσθαλίσκον. On a objecté l'absence de l'article, mais je ne pense pas l'objection fondée. Exprimé une première fois devant χοιρίον, l'article peut, à la rigueur, être omis devant ποσθαλίσκον qui suit\note{De même qu'au v. 423 des \italique{Grenouilles}, l'article exprimé devant γνάθους manque devant πρωκτόν, qui précède cf. ma note sur \italique{Lysistr.}, v.151. Si l'article s'omet devant le premier terme, à plus forte raison peut on l'omettre devant le second. Exemples : \italique{Oiseaux}, 653, τὴν ἀλώπεχ᾽ ὡς ἐκοινώνησεν ἀετῷ ; \italique{Eccl.}, 702, τοῖς εὐπρεπέσιν δ᾽ ἀκολουθοῦντες καὶ μειρακίοις ; \italique{Cav.}, 320 ; \italique{Nuées}, 622 ; \italique{Thesm.}, 491 ; \italique{Gren.}, 773 et 1477 ; \italique{Ecclés.}, 51 ; Platon, \italique{Men.}, 244c ; Thucyd., II, 13, 1, etc. Pour moi, il n'est pas douteux qu'il ne faille lire au v. 142 de notre pièce : ποῦ τὸ πέος; comme au vers suivant : ποῦ τὰ τιτθία; (c'est d'ailleurs la leçon de Suidas et du scholiaste des \italique{Guêpes}, v. 1158)}. Seulement il est plus que douteux que ce dernier mot soit grec, les diminutifs se faisant en ισκος, et non en αλισκος\note{On cite, il est vrai, κωραλίσκος, παιδαρίσκος et μειρακίσκος. Ce serait concluant s'il n'y avait pas l'intermédiaire de κωράλιον, παιδάριον et μειράκιον. Si des diminutifs comme ποσθάλιον et ποσθαλίσκος avaient existé, ils n'auraient pu manquer d'être relevés par d'anciens lexicographes. C'est l'unique raison qui m'empêche de proposer σαθίσκον (qui serait parfait paléographiquement) au lieu de σάθωνα.}. 

Il y a peut-être moyen de serrer de plus près la leçon du manuscrit. Je ne crois pas qu'il y ait à reprendre dans καὶ πρός, \italique{et insuper}, formule adverbiale fréquente chez Aristophane : \italique{Cav}., 578 ; \italique{Lysist.}, 1238 ; \italique{Ploutos}, 1001 ; et devant un accusatif, comme ici : \italique{Acharn}., 1229 ; \italique{Lysist}., 628 ; \italique{Ecclés}., 308 ; \italique{Thesmoph.}, 802\note{Je cite ce vers, bien que M. van Leeuwen l'ait encore altéré. Σκεψώμεθα κἀντιτιθῶμεν πρὸς ἕκαστον n'est pas plus difficile à entendre que : δίκην δίδωμι σοι καὶ χάριν πρὸς εἴσομαι (\italique{Guêpes}, 1420), ou bien λαλεῖν γυναῖκας καὶ διαλλάττειν πρὸς ἡμᾶς (\italique{Lysistr.}, 628) et οὗτος δ᾽ ἀκολουθεῖ κἀμε πρὸς βίαζεται (\italique{Ploutos}, 16)}. 
La difficulté gît dans θάληκον, qui n'a pas de sens. En parlant de sa fille, Mnésiloque s'est servi du mot χοιρίον, terme de caresse un peu familier, je veux bien, mais nullement obscène (\italique{Guêpes}, 572; Χοιρίλη est un nom de femme assez fréquent). 
À χοιρίον correspond, en parlant d'un garçon, πόσθων ou σάθων, autres termes de caresse (\italique{Paix}, 1300 ; Photius, p. 496, 5 [σ 21 ed. Theodoridis]). Or il se trouve que ce dernier mot rentre parfaitement dans le vers. On pourrait donc lire le passage comme suit ; je ne réponds pas, bien entendu, que ce passage soit le texte d'Aristophane, mais en tous cas le changement n'est pas grand, et c'est du grec :\\

\citComm{ἔπειτα καὶ τὸ χοιρίον ἀνδρὸς μοι τυχεῖν\\
πλουτοῦντος, ἄλλως τ' ἐλιθίου κἀβελτέρου,\\
καὶ πρὸς σάθωνα νοῦν ἔχειν μοι καὶ φρένας}\note{Le scholiaste n'est pas inconciliable avec la leçon que je propose. Le bonhomme soupçonne que, pris au sens étymologique, χοιρίον et σάθων peuvent passer pour des gravelures, mais il n'en est pas trop sûr. D'abord au sujet de la fillette : ἴσως δὲ τὸ αἰδοῖον τὸ γυναικεῖον αἰνίττεται ; puis à propos du garçon : ἴσως δὲ παρὰ τὴν πόσθην αὐτὸ συνέθηκεν.}\\

G. Hermann trouve que ce dernier trait n'est pas assez piquant. Mais qui nous dit qu'Aristophane ait voulu faire de l'esprit ? Il lui suffit de prêter à son personnage le caquet banal d'une commère athénienne ; c'est sa manière à lui d'être piquant. Le souhait que la soi-disant mère forme pour son fils avait la valeur d'une sorte de dicton. Car on le trouve sous deux formes différentes dans Ménandre :

\citComm{εὐδαιμονία τοῦτ᾽ ἔστιν υἱὸς νοῦν ἔχων.}\note{\italique{Fragm. Com. Gr.}, t. IV, pp. 86 et 261 [$=$ fr. 58(54) ed. Kassel-Austin].}

\Note{Note 9 aux \italique{Thesmophories} : vers 352-354}


\citComm{ξυνευχόμεσθα τέλεα μὲν\\
πόλει, τέλεα δὲ δήμῳ\\
ταδ' εὐγματα γενέσθαι.}\\

Ce dernier vers pèche contre le mètre. M. van Leeuwen adopte la conjecture de von Velsen : ἅπαντα τάδε γενέσθαι. On pourrait en trouver de moins plates. 
Ainsi rien n'empêche de considérer γενέσθαι comme une glose et de lire :

\citComm{ταδ' εὔγματ᾽ ἀποτελεῖσθαι}\\

Cf. Platon, \italique{Rép}., IV, p. 443b : τέλεον ἡμῖν τὸ ἐνύπνιον ἀποτετέλεσται ; \italique{Phœdr.}, p. 249c : τελέους ἀεὶ τελετὰς τελούμενος. Et pour ce qui est du présent, Sophocle, \italique{Ajax}, 686 : εὐχου τελεῖσθαι τοὐμὸν ὧν ἐρᾷ κέαρ.


\Note{Note 10 aux \italique{Thesmophories} : vers 365-367}

\citComm{\decalage{365aaaaa}ὁπόσαι δ᾽}
\citComm{\decalage{365aaaaa}ὅρκους τοὺς νενομισμένους}
\citComm{360\decalage{aaaaa}κερδῶν οὕνεκ᾽ ἐπὶ βλάβῃ...}
\citComm{365\decalage{aaaaa}ἢ Μήδους ἐπάγουσι τῆς}
\citComm{\decalage{365aaaaa}χώρας οὕνεκ᾽ ἐπὶ βλάβῃ}
\citComm{\decalage{365aaaaa}ἀσεβοῦσ᾽ ἀδικοῦσί τε τὴν πόλιν.}\\

Les vers 365 à 367 ont été changés de tant de manières, qu'on serait tenté de les tenir, avec les derniers éditeurs, pour à peu près désespérés. 

L'erreur provient de ce qu'on a attribué à ἐπάγειν le sens de \italique{adducere}, que la moindre réflexion aurait dû faire écarter. 
Les Athéniens savaient trop bien depuis Marathon ce que valait le soldat mède, pour que l'idée pût venir à personne de l'introduire en Attique comme auxiliaire.

Ἐπάγειν signifie ici \italique{incitare}, \italique{stimulare}. L'expression d'Aristo phane se lit dans Hérodote, sous la forme de ἐπῆγον τὸν Πέρσην (IX, 1). 
Et il faut croire qu'elle était d'usage courant, car elle se lit aussi dans une lettre du roi Philippe aux Athéniens : ἐπάγουσι τὸν Πέρσην ἐπὶ τοὺς Ἕλληνας\note{Démosthène, p. 160, 15.}. 

Dès lors la suite s'entend toute seule : ἡ χώρα désigne ἡ χώρα ἡ Ἀττική, les exemples abondent\note{\italique{Guêpes}, 652 ; \italique{Lysistr.}, 524 ; \italique{Grenouilles}, 380, etc.}. Six vers plus haut il était question de ceux qui transgressent leurs serments « dans leur propre intérêt » κερδῶν οὕνεκα [ἐπὶ βλάβῃ]. Ici il s'agit de ceux qui sollicitent les Mèdes « dans l'intérêt du pays » τῆς χώρας οὕνεκα [ἐπὶ βλάβῃ]. 
« Dans l'intérêt du pays » est tout simplement un euphémisme équivalent à « contre les Péloponésiens », ce qu'il fallait éviter de dire. 
On sait, en effet, que les oligarques intriguaient auprès des satrapes d'Asie Mineure pour détacher le grand roi de l'alliance de Sparte et le faire entrer dans celle d'Athènes\note{Thucydide, VIII, 53.}. Mais l'on sait non moins bien que cette politique indignait Aristophane. 
Voilà pourquoi il ajoute ἐπὶ βλάβῃ, formule concise et claire qui se lit également dans Xénophon (\italique{Cyr.}, 2, 2, 12, et \italique{Mémor.}, 2, 3, 19) et dont on peut rapprocher ἐπὶ λύμῃ (Hérod., 3, 14, 1), ἐπὶ λώβᾳ (Soph., \italique{Ant.}, 792) et ἐπὶ δυσκλείᾳ (Id., \italique{Ajax}, 143). 

Que si l'on demande à qui il fait ici le procès (car même quand il généralise, il est rare qu'il n'ait pas en vue quelque particulier), nous pouvons nettement répondre à la question. L'homme qui vise à bouleverser la loi et les décrets, qui révèle à l'ennemi les secrets de l'État et cherche à traiter avec lui dans l'intérêt soi-disant d'Athènes, mais en réalité dans son propre intérêt et celui de sa coterie, cet homme n'est autre que Pisandre. 
La preuve est facile à faire. Deux mois avant les \italique{Thesmophories}, Aristophane avait fait représenter \italique{Lysistrata} ; les deux pièces ont donc été composées en même temps et s'inspirent du même esprit. 
Or dans \italique{Lysistrata} c'est Pisandre qui est désigné, nominativement cette fois, comme le chef de ceux qui conspirent contre l'État, comme l'homme qui ne cesse de susciter du grabuge « pour être en mesure de voler » ; et pendant que l'ex-démagogue converti à l'oligarchie négociait avec Tissapherne\note{Id., VIII, 56.}, le poète patriote dénonçait le Perse comme l'ennemi héréditaire (v. 1133); bien plus, il se complaisait dans le chœur final à évoquer les glorieuses victoires remportées par Athènes et Sparte réunies, à l'Artémision et aux Thermopyles. 

Quant au dernier vers cité ci-dessus, je ne vois pas de motif pour le suspecter. Sans doute le poète aurait pu conclure par un souhait de malheur, quelque chose çomme le ἐξώλεις εἶεν αὐτοὶ καὶ γένος τὸ ἐκείνων, réclamé par M. de Wilamowitz ; mais le poète comique n'est pas astreint à la même ordonnance logique que le prosateur. Il ne s'agit pas ici d'une formule officielle, mais d'une fantaisie ou plutôt d'une parodie, d'ailleurs aisée à entendre. 
Dire que les parjures et les traîtres dont on vient de parler « sont coupables de lèse-majesté divine et de lèse-nation », n'est-ce pas appeler sur eux la colère des dieux et des hommes ? En tous cas le vers n'est pas de ceux qui sentent l'interpolateur : ἀσεβεῖν καὶ ἀδικεῖν sont des termes consacrés, formant une seule locution, et dont on trouve de nombreux exemples dans les auteurs du temps\note{Xénoph., \italique{Hell.}, II, 3, 53 : οὐ μόνον εἰσὶ περὶ ἀνθρώπους ἀδικώτατοι ἀλλὰ καὶ περὶ θεοὺς ἀσεβέστατοι ; \italique{Cyrop}., V, 2, 9 : πολλοὺς μὲν οἶμαι εἶναι ἀνθρώπους, οἱ οὔτε ἀσεβεῖν ἂν θέλοιεν οὔτε ἀδικεῖν ; Démosth., \italique{c. Néer.}, 1349, 6 : περιφανῶς εἰς τοὺς θεοὺς ἀσεβοῦσαν καὶ εἰς τὴν πόλιν ὑβρίζουσαν, etc}.

\Note{Note 11 aux \italique{Thesmophories} : vers 400-401}


\citComm{ \decalage{	γυνὴ στέφανον}ὥστ᾽ ἐάν τις πλέκῃ,\\
	γυνὴ στέφανον ἐρᾶν δοκεῖ.}\\

Il manque une syllabe au v. 400. On y a suppléé de diverses façons, mais on n'a pas songé à la plus simple de toutes, qui serait de lire ξυμπλέκῃ, au lieu de πλέκῃ, comme au v. 61 ξυγγογγυλίσας pour γογγυλίσας, au v.658 διαθρῆσαι pour ἀθρῆσαι, et au v. 1251 des \italique{Guêpes}, συσκεύαζε pour σκεύαζε\note{Lire la note capitale de Brunck sur \italique{Lysistr}., v. 408 ; \italique{Eccl.}, 758 ; \italique{Lysistr.}, 465, 552, 613, 997 ; \italique{Gren.}, 83, 1026 ; \italique{Thesmoph.}, 1129 ; \italique{Guêpes}, 147 ; \italique{Caval.}, 418, 655, 989, 1328, 1363.}. On dit également bien πλέκειν et συμπλέκειν στέφανον\note{Ainsi \italique{Anthol. Pal.}, XII, 257, 6.}.


\Note{Note 12 aux \italique{Thesmophories} : vers 418-420}


\citComm{\decalage{ἃ δ᾽ ἐν ἡμῖν}ἃ δ᾽ ἐν ἡμῖν πρὸ του\\
	αὐταῖς ταμιεύεσθαι προαιρούσας λαβεῖν\\
	ἄλφιτον, ἔλαιον...}\\

Le v. 419 est légèrement altéré. Mais, qu'on lise avec Reiske : ταμιεῦσαι καὶ\note{C'est ainsi qu'on a restitué dans Xénoph. \italique{Hier.}, 3, 9, au lieu de φιλεῖσθαι, la vraie leçon conservée par Stobée : φιλεῖν καὶ}, ou, moins heureusement avec Dobree, ταμιευούσαις, le sens n'en est pas affecté. 4Mais qu'avons-nous à faire de la déplorable conjecture de Scaliger : λαθεῖν (au lieu de λαβεῖν), reprise par les derniers éditeurs ? 
On n'y gagne que de mettre le poète en contradiction avec lui-même. Quoi ! les femmes déclarent qu'autrefois elles administraient elles-mêmes le ménage, et elles ajouteraient que pour toucher aux provisions force leur était de se cacher. Ne voit-on pas que ces deux assertions se détruisent l'une l'autre ? De deux choses l'une, en effet, ou bien elles avaient l'intendance et pouvaient se rendre ouvertement au cellier, ou bien c'était le mari qui disposait des provisions, et dès lors elles n'avaient point l'intendance. 
Impossible de sortir de là. 

Le scholiaste explique προαιρούσας λαβεῖν par ἐκ τοῦ ταμιείου λαβεῖν. En effet, προαιρεῖν c'est \italique{promere}, « tirer une chose de l'endroit où elle était serrée ». 
L'huile, la farine, le vin étaient logés dans le cellier, ταμιεῖον, lequel formait l'un des côtés du portique intérieur de la maison\note{Heindorf sur Platon, \italique{Protag}., VII, p. 314e.}, et tenait lieu de cave. On disait προαιρεῖν ἐκ τοῦ ταμείου, ou simplement προαιρεῖν\note{Thucyd., VIII 90, 5; Phérécrat., ap. Athen., 481a. — Quand il s'agissait d'un lieu autre que le cellier, où l'on ne pénétrait qu'avec de la lumière (\italique{Ecclés.}, 15, on disait ἐξαιρεῖν (\italique{Paix}, 1145 ; \italique{Thesmoph}., 284. La différence est bien marquée dans cette phrase de Thucydide : τὸν σῖτον ἐξαιρεῖσθαι (tirer du navire) καὶ ἐντεῦθεν προαιροῦντας (tirer de l'entrepôt) πωλεῖν VIII, 90, 3.}. Λαβεῖν προαιροῦσας signifie donc « prendre en les tirant du cellier ».

Il paraît par un autre endroit d'Aristophane (\italique{Eccles.}, 26) que les maris qui s'étaient réservé l'accès du cellier n'en étaient pas plus avancés, et qu'en dépit de toutes les précautions les femmes trouvaient moyen d'y pénétrer. De même, on avait eu beau appliquer des cachets, σφραγῖδας (v. 415), sur l'appartement des femmes, celles-ci avaient imaginé de faire graver des cachets tout pareils, σφραγίδια (v. 427), qui leur permettaient d'entrer ou de sortir à volonté. Car M. van Leeuwen se trompe quand il dit que la θύρα du v. 424 désigne la porte du cellier, laquelle se fermait à clef.


\Note{Introduction aux notes sur la \italique{Paix}}\note{Extrait des \italique{Bull. de l'Acad. roy. de Belgique} (Classe des lettres, etc.), n° 12, décembre 1899.}

\italique{Tome second, pp. 93-99.}\\


La \italique{Paix} est la plus endommagée des comédies d'Aristophane. Non que le texte y ait subi de bien graves altérations : nous avons soutenu naguère et nous maintenons, contrairement à l'opinion reçue, que le grand comique est celui des poètes grecs, à part Homère, qui nous a été le plus fidèlement transmis. Mais cela n'empêche pas qu'en le lisant nous ne soyons souvent embarrassés pour des raisons qui ne tiennent pas toutes à notre ignorance. Et c'est le cas surtout pour la \italique{Paix}. 

Pourquoi cette pièce se trouve avoir été moins bien traitée que les autres, nous n'en savons rien, et en vérité il nous importe assez peu de le savoir. C'est heureusement la moins importante de celles qui nous sont parvenues. La cause qu'on y défend était pour les contemporains d'un intérêt actuel et poignant. Mais cet intérêt a dû s'affaiblir assez vite, à mesure que les temps changeaient. Il est à croire que la postérité immédiate en jugeait à peu près comme nous, et que dans l'antiquité même, la \italique{Paix} était moins appréciée, et partant moins lue et étudiée, que les chefs-d'œuvre du poète. 

Il ne faut pas s'exagérer le tort des copistes. Des fautes, ils en ont laissé échapper, et plus ici qu'ailleurs ; mais le plus souvent ils ont péché par excès de conscience. En ne voulant rien omettre, ils ont passé la mesure : car si la pièce a souffert, c'est surtout par les interpolations. 

Je voudrais à ce sujet entrer dans quelques détails. Il arrive que des éditeurs, plus attachés à la lettre qu'à l'esprit, s'obstinent à restituer certains passages qui n'ont et ne peuvent avoir de sens plausible. Sans se préoccuper de la valeur et de la suite des idées, ils estiment avoir besogne faite dès qu'ils ont réussi tant bien que mal à remettre le vers sur pied. D'autres, au contraire, trop confiants dans la sûreté de leur jugement, n'hésitent pas à rejeter sous de futiles prétextes tout ce qui ne cadre pas avec leur propre conception du sujet. Les deux manières de procéder sont également condamnables. En si délicate matière, il faut savoir se tenir en garde et contre le respect superstitieux du texte et contre les raisons dictées par la seule fantaisie. 

Quand on compare entre elles les diverses éditions d'Aristophane, on est frappé de les trouver si peu d'accord sur ce point. Je crois qu'avec un peu d'attention il y aurait moyen de mettre le lecteur hors d'incertitude. Après tout, la difficulté n'est peut être pas si grande qu'on le pense de déterminer à quels signes se reconnaît une interpolation. Il en est en effet des mots intercalés dans un texte comme des retouches dans la peinture : rarement ils sont si bien fondus dans l'ensemble qu'un œil exercé ne puisse les discerner avec assurance. Encore convient-il de distinguer entre les œuvres de poésie et les œuvres de prose. Si dans ces dernières, pour des raisons que chacun devine, la fraude est adroite à se dissimuler, chez les poètes du moins, et notamment chez Aristophane, elle est le plus souvent apparente et se trahit par plus d'un trait. Et même, à les examiner de près et les comparer entre elles, on réussit sans trop d'efforts à classer sous de certains chefs les passages interpolés. Ainsi : 

1°  	 Il arrive que le même vers soit répété, sans raison aucune, en deux endroits de la même pièce : le vers 436 des \italique{Acharniens} a été reproduit, on ne sait pourquoi, d'après le vers 384, et le vers 928 d'après 905 et 931 ; le vers 114 des \italique{Cavaliers}, d'après le vers 96 ; le vers 192 des \italique{Oiseaux}, d'après le vers 1218 ; le vers 1218 de la \italique{Paix}, d'après le vers 1193 ; le vers 281 du \italique{Ploutos}, d'après le vers 260 et 1173 d'après 968.

2° 	 D'autres fois, ce n'est pas le même texte, mais c'est la même pensée exprimée presque dans les mêmes termes en deux vers qui se suivent. Ainsi, dans les \italique{Acharniens}, Dicéopolis demande aux fillettes du Mégarien si elles mangeraient bien des pois chiches. Et il ajoute : « quoi encore, des figues phibalées ? » (vers 802). Un glossateur ajoute dans la marge en guise d'explication :  σῦκα τρώγοις ἄν ; puis un copiste, prenant cela pour un vers, complète à sa façon : σῦκα τρώγοις ἄν αὐτός; -- κοΐ, κοΐ. Aujourd'hui, grâce à G. Dindorf, ce prétendu vers a disparu de toutes les éditions. 

En ce cas, la difficulté consiste à ne pas se méprendre, comme on l'a fait, sur celui des deux qui est l'authentique. Exemple : \italique{Paix}, 896 et 897 ; \italique{Grenouilles}, 1431 et 1432; \italique{Ecclésie}, 799 et 800. Il en est de même des vers 717, 718 des \italique{Acharniens}, qui ne sont autre chose que la traduction en algonquin des deux vers grecs qui précèdent. On peut rapporter encore à cette catégorie levers 1029 des \italique{Guêpes}, emprunté à un passage analogue de la \italique{Paix}, 751. 

3°  	 Parfois aussi un vers interpolé récapitule ce qui précède ou ce qui suit. On dirait que certains manuscrits étaient pourvus de ce qu'en termes d'imprimerie on appelle des manchettes, lesquelles souvent ont passé de la marge dans le texte. Ainsi le vers 176 de la \italique{Lysistrata} résume les trois suivants ; le vers 790 des Grenouilles, les deux qui précèdent et les vers 978, 979 les neuf qui suivent. A cette catégorie se rattache aussi levers 1343 des \italique{Oiseaux}, déjà signalé comme interpolé par le scholiaste. 

4°  	 Mais, le plus souvent, les interpolations consistent en notes explicatives. Parfois simples paraphrases : \italique{Paix}, 420, énumération arbitraire et très inutile des fêtes athéniennes ; \italique{Oiseaux}, 523, où l'intercalation est attestée par le vers antistrophique (611) ; 1343, paraphrase de ὀρνιθομανῶ ; \italique{Thesmophories}, 23, de προσμάθοιμι ; \italique{Grenouilles}, 15, de εἴωθε ποιεῖν.. Cela arrive souvent quand le copiste se trompe sur l'emploi d'un verbe pris absolument : \italique{Guêpes}, 565, est le développement de προστιθέασι ; \italique{Grenouilles}, 168, de μίσθωσαί τινα ; \italique{Ploutos}, 281, sert de régime à φράσαι. 

5° 	Toutefois le copiste se laisse plus volontiers tenter quand la note contient quelque détail précis. Je donnerai comme exemple le vers 627 de l'\italique{Ecclésie}. Aristophane avait écrit 

\citComm{ἀλλὰ φυλάξουσ᾽ οἱ φαυλότεροι τοὺς καλλίους ἀπιόντας\\
	ἀπὸ τοῦ δείπνου κοὐκ ἔξεσται παρὰ τοῖσι καλοῖς καταδαρθεῖν.}\\


Un maître d'école, pour expliquer φυλάξουσι, avait écrit en marge, probablement d'après le vers 694 : καὶ τηρήσουσ᾽ ἐπὶ τοῖσιν δημοσίοισιν : « ils les guetteront dans les rues ». Un copiste a glissé cela dans le texte après les mots ἀπὸ τοῦ δείπνου et, pour compléter le vers suivant, a répété οἱ φαυλότεροι. 

\textbf{Tels sont aussi le vers 500 des \italique{Thesmophories} sur l'anecdote du manteau}, 584 de \italique{Ploutos} sur les jeux olympiques. 

6° 	Ce qui est plus drôle, c'est quand on nous donne comme d'Aristophane les réflexions banales de l'un ou l'autre de ses interprètes. 
Lisez, pour vous en faire une idée, les vers 181-182 des \italique{Oiseaux}, où l'on vous explique d'où le mot πόλος tire son origine. 
Qu'il se soit rencontré un copiste assez borné pour mettre pareilles inepties sur le compte du poète, cela se conçoit à peine, mais ce qui ne se comprend plus du tout, c'est que des éditeurs les aient autorisées. D'autres traits d'esprit, sans être de cette force, trahissent pourtant leur origine. Celui-ci, par exemple, qui se lit dans les \italique{Grenouilles} (935) : « Est-il permis de mettre un coq dans une tragédie ? » ou cet autre dans le \italique{Ploutos} (566) : « N'est-ce pas prouver son respect pour la règle que de se cacher pour voler ? » 
Il arrive même que l'interpolateur soit inintelligent au point de faire dire à un personnage exactement le contraire de ce qu'il lui faudrait dire. Le vers 565, déjà cité, des \italique{Guêpes} et le vers 897 du \italique{Ploutos}, pourraient ici servir d'exemples. Il est rare d'ailleurs qu'une interpolation exprime nettement ce qu'elle est censée dire. 
Presque toujours la fraude se révèle par une entorse, soit au vocabulaire, soit à la grammaire, soit à la prosodie : \italique{Acharniens}, 718, l'article τῷ devant γέροντι et νέῳ ; 803,  σῦκα avec υ long ; \italique{Guêpes}, 565, ἀνιῶν, cheville ; 1029, ἤρξε pour ἤρξατο ; \italique{Paix}, 744, καὶ τούτους οὕνεκα τουδί, simple cheville ; 1135, ἐκπεπρισμένα, cheville et faute de métrique ; \textbf{\italique{Thesmophories}, 500, vers presque tout en chevilles} ; \italique{Grenouilles}, 790, ἐκεῖνος répété à deux vers de distance, pure cacologie, qu'on le rapporte les deux fois à Sophocle, ou la première à Sophocle, la seconde à Eschyle ; \italique{Ecclésie}, 628, οἱ φαυλότεροι, cheville ; \italique{Ploutos}, 387, δικαίους, faute de prosodie. 
J'insiste surtout sur ce dernier vers, fortement interpolé. Il est non seulement redondant, mais en contradiction avec ce qui précède. 
Que si on lit, avec le \italique{Ravennas}, δικαίους, au lieu de δεξιούς, il n'est plus que redondant, c'est vrai, mais il pèche contre la métrique. 
Ce qui caractérise surtout l'interpolation, c'est que, tint-elle directement au sujet, elle fait toujours longueur ; supprimez-la, non seulement le texte s'en trouve allégé, mais il en devient plus clair et généralement mieux suivi. 

Sur ce point, toutefois, il s'agit de bien s'entendre. Tout ce qui n'est pas nécessaire n'est pas pour cela seul explétif. Autre chose est l'ampleur, autre chose la prolixité. Le goût différant chez chacun, l'éditeur risque fort de se tromper, qui prend trop vite parti dans une pure question de goût. Il lui convient de rester dans son rôle et de ne point oublier que ce rôle est passif. 
Où irions-nous si, n'écoutant que ses préférences, chacun s'avisait d'éliminer ce qui lui semble superflu ? C'est pourtant ce qui se voit trop souvent. Je ne citerai qu'un exemple. Pourquoi, dans les \italique{Grenouilles}, M. van Leeuwen écarte-t-il les vers 26 à 29 ? Mais ils sont charmants ces vers, à condition qu'on s'en tienne à la leçon du \italique{Ravennas}. Ce trait si habilement amené : « Le fardeau que tu portes, n'est-ce pas un âne qui le porte? » est de l'Aristophane et du meilleur. Mais ils sont, dites-vous, inutiles. Je reconnais qu'ils ne prouvent pas grand'chose. La belle raison ! Comme si une action dramatique se déroulait à la façon d'un théorème de géométrie ! Relisez plutôt la première scène de \italique{Roméo}. 
Que de détails on y pourrait supprimer sans que l'intérêt de la pièce en soit compromis. Laissons donc ces vers à Aristophane, et croyons bien que sa gloire n'en souffrira pas. 

Une dernière remarque qui ne laisse pas d'avoir son prix. Certaines pièces d'Euripide, l'\italique{Hippolyte} par exemple, sont restées au théâtre ; on n'a cessé durant des siècles de les représenter et partant de les remanier. Les pièces d'Aristophane, nées de circonstances spéciales, n'ont guère été reprises après lui. Les interpolations qu'on y constate sont le fait, non d'acteurs ou de reviseurs, mais de simple copistes. Il en résulte qu'elles sont peu nombreuses et ne dépassent jamais un ou deux vers. 

Soyons indulgents pour les copistes et disons-nous qu'à l'occasion nous aurions fait comme eux. En voulez-vous la preuve ? Dans les \italique{Nuées}, vers 1146, Strepsiade, après avoir salué Socrate, ajoute : « Accepte donc cela »  ἀλλὰ τουτονὶ πρῶτον λαβέ. \italique{Ceci}, c'est un sac de farine. 
Une scholie l'explique par : τὸν θύλακον, ὅς ἔστι μεστὸς ἀλφίτων. Le hasard, qui est quelquefois plaisant, comme dit M\expo{me} de Sévigné, a fait que cette glose ressemble singulièrement à un trimètre iambique\note{πολλοὶ μέτρα ἰαμβικὰ λαλοῦσιν οὐκ εἰδότες.« Bien des gens font des vers sans le savoir. » Démétrius de Phalère, §43.}. G. Hermann, un des maîtres de la philologie en Allemagne, ne manque pas de le constater. Il ne s'en faut, dit-il, que d'une petite lettre, \italique{una litterula}, un simple γ' à ajouter. 

Lisez : 
\citComm{τὸν θύλακον γ᾽ ὅς ἐστι μεστὸς ἀλφίτων,}\\ 
et vous aurez un vers excellent. Il est vrai qu'Hermann s'est fait scrupule d'introduire cet excellent vers dans le texte, ayant considéré sans doute que l'autorité d'une quarantaine de manuscrits compte bien aussi pour quelque chose. Mais on voit qu'il lui en coûte et que si, comme les Anciens, il n'avait eu sous les yeux que l'unique copie contenant la glose, il n'aurait pu résister à la tentation. Eux aussi, les copistes, n'y mettaient point dem alice. Au contraire, ils croyaient servir leur auteur, en lui restituant un vers ou un membre de phrase omis à tort par leurs prédécesseurs. 

Plus, je crois, que les autres comédies, la \italique{Paix} a subi des interpolations. Pour les vers 87 à 89, 98 et 273, il y a lieu tout au moins d'hésiter, et dans le doute la prudence commande de s'abstenir. Le vers 850 est au-dessus de tout soupçon : « Il y en a dans l'Olympe qui vivent de ce métier (de proxénète), » le mot est aussi spirituel que vrai ; de plus, il porte la marque de l'ouvrier, prenant occasion de tout pour se moquer des dieux de l'Olympe. 

Mais le vers 420 est pour sûr apocryphe : Μυστήρι᾽ Ἑρμῇ, Διπόλει, Ἀδώνια, l'énumération est à la fois redondante et incomplète ; en outre elle est inexacte, les Adonies n'étant pas une fête spécialement athénienne. Les vers 744 et 1218 sont dans le même cas. Le premier, incorrect et oiseux, a fait intervertir mal à propos l'ordre des deux précédents. Le second n'est qu'une répétition maladroite du vers 1193. 

Dans les notes qui vont suivre, j'aurai l'occasion de signaler quelques autres passages interpolés, encore que mon principal but soit de défendre la leçon traditionnelle contre des corrections qui me paraissent injustifiées\note{Je me servirai surtout de la grande édition de la \italique{Paix} donnée par M. van Herwerden, professeur à l'Université d'Utrecht : Ἀριστοφάνους Εἰρήνη, \italique{cum scholiorum antiquorum excerptis passim emendatis, recognovit et adnotavit Henricus van Herwerden}, Lugduni Batavorum, A. W. Sijthoff, 1897, 2 vol. in-8°. Elle est faite en conscience, et l'appareil critique y est très au complet. M. van Herwerden sait du grec autant qu'homme d'Europe. Je crains seulement qu'il ne soit un peu trop engagé dans la critique conjecturale.}.



\Note{Note 13 aux \italique{Thesmophories} : vers 499-501}

\citComm{ὡς ἢ γυνὴ δεικνῦσα τἀνδρὶ τούγχυχλον\\
\{ὑπ᾽ αὐγὰς οἷόν ἐστιν, ἐγκεκαλυμμένον\}\\
τὸν μοιχὸν ἐξέπεμψεν, οὐκ εἰρηκέ πω}\\

Il faut supprimer le v. 500, qui tire son origine d'une note marginale introduite mal à propos dans le texte. Tout ce qui décèle l'interpolation s'y trouve : faute de prosodie, chevilles, détail oiseux. 
Ὑπ᾿ αὐγάς pèche contre le mètre; ὑπ᾽ ὄρθρον, que Dindorf y substitue, signifierait non au jour, mais à l'aube, et n'est donc pas de mise ici ; οἷόν ἐστ' est une pure cheville ; ἐγκεκαλυμμένον, au moyen de quoi l'on a complété le vers, est une autre cheville, à moins qu'on ne préfère le traduire par « enveloppé dedans », c'est-à-dire dans le manteau déployé à la lumière, ce qui implique contradiction.

Du vers entier, il ne reste donc que ὑπ᾽ αὐγάς, servant à préciser le sens de δεικνῦσα. Et si l'on veut bien y réfléchir, le récit ne peut que gagner à la suppression de ce détail ; car la femme qui veut faire échapper son amant atteindra plus sûrement son but en ouvrant le manteau devant les yeux de son mari qu'en le déployant à la lumière. 

L'anecdote racontée par Mnésiloque était probablement connue de tout le monde. La reine de Navarre (6° nouv. de l'\italique{Heptaméron}) en raconte une autre à peu près pareille, afin de prouver « combien est prompte et subtille une femme à eschapper d'un dangier ».


\Note{Note 14 aux \italique{Thesmophories} : vers 509-510}


\citComm{ἄπελθ᾽ ἄπελθ᾽, ἤδη γαρ, ὦ ᾽νέρ, μοι δοκῶ\\
τέξειν.}\\

Les Attiques disaient τέξω et τέξομαι, et l'on ne voit pas de raison pour préférer l'une forme à l'autre. 
Les modernes atticistes n'admettent pas cela, non plus qu'ils n'admettent οἴσομαι à côté de οἵσω (\italique{Paix}, 1033), ni ἰδέσθαι à côté de ἰδεῖν (\italique{Thesm.}, 800), ni οὐρήσω à côté de οὐρήσομαι (\italique{Guépes}, 394), ni στύω à côté de στύομαι (\italique{Lysistr.}, 598), ἱκετεύομαι à côté de ἱκετεύω (\italique{Eccl.}, 915), ἔρρευσα à côté de ἐρρύην (\italique{Cav.}, 526), etc., etc. En un mot, ils ne souffrent pas de doubles emplois. Hirschig corrige donc : μοι δοκῶ τίκτειν. 

Vous objecterez que τέξω se lit dans Eschyle et dans Euripide, et même une autre fois dans Aristophane, \italique{Cav.}, 1037. Il paraît que cela ne compte pas. Dans ce dernier cas, il s'agit d'un pseudo-oracle; or Cobet pose en fait que \italique{Sibylla italoquebatur in oraculis et Dii immortales et heroes ; mulierculœ atticœ τέξομαι solebant dicere}. Ce n'est pas impossible, mais ne trouvez-vous pas que pour le prendre sur ce ton, il faudrait avoir beaucoup fréquenté les petites femmes athéniennes ? Je ne sais pourquoi cela me remet en l'esprit l'anecdote si vivement racontée par M\expo{me} de Caylus sur le marquis de Lassay. Un jour qu'il protestait avec insistance de la vertu de M\expo{me} de Maintenon et s'en portait garant, sa femme l'interrompant doucement lui dit : « Comment faites-vous, Monsieur, pour être sûr de ces choses-là ? »

\Note{Note 15 aux \italique{Thesmophories} : vers 531-532}


\citComm{ἀλλ' οὐ γαρ ἐστ' τῶν ἀναισχύντων φύσει γυναικῶν\\
οὐδὲν κάκιον εἰς ἅπαντα πλὴν ἄρ᾽ εἰ γυναῖκες.}\note{M. Blaydes prétend qu'il faut lire κάκιον εἰς τὰ πάντα. Je me borne à citer Thucyd., 8, 76, 3 : ἐς πάντα ποριμωτέρων, et Xénophon, \italique{Cyr.}, 3, 3, 20 : βέλτιον πρὸς πάντα. Rien de plus semblable sous le soleil qu'une correction de M. Blaydes et une autre correction de M. Blaydes.}\\

C'est tout ce qu'il y a de clair : « Les femmes au naturel impudent, il n'est rien de pire à tous égards, si ce n'est des femmes ». 
Oct. Feuillet a dit de même : « Rien de plus semblable sous le soleil qu'une femme — et une femme\note{\italique{Dalila}, partie IV. C'est la traduction presque textuelle d'une sentence de Ménandre, \italique{Monost.}, 109.} ». Et l'on cite de Ménandre : « Il n'est rien de misérable comme un vieillard amoureux, si ce n'est un autre vieillard amoureux\note{\italique{Fragm. Com. Gr.}, t. IV, p. 220} ». Ou encore le mot historique de Talleyrand sur Maret, duc de Bassano : « Je ne connais pas de plus grande bête au monde que Maret, si ce n'est le duc de Bassano ». 

Corriger, comme le font les éditeurs, d'après Fritzsche : πλὴν ἄρ᾽ εἰ γυνή τις, c'est prêter gratuitement au poète un paralogisme. On ne dit pas : « Il n'est rien de pire que les femmes effrontées sauf une femme ». Laquelle ? Au moins faudrait-il ajouter : « une femme plus effrontée », comme on lit encore dans Ménandre : « Rien de plus malheureux qu'un père, si ce n'est un père affligé d'une famille plus nombreuse\note{\italique{Ibid}, p. 261, CX [$=$ fr. 819(596) Kassel-Austin].} », et dans La Bruyère : « Qui est plus esclave qu'un courtisan assidu, si ce n'est un courtisan plus assidu ? » 

Et voilà où mène la manie des conjectures.\\


\Note{Coupe d'Onésime}
\fantome 
\begin{figure}[!h]
\centering
\includegraphics[scale=0.5]{Onesime.png}
\caption{\italique{Coupe d’Onésime.} (Bâle, Antikenmuseum BS 440. ARV 326, 86 \italique{bis. Para} 359. Elle est reproduite par John Boardman, fig. 233 in \italique{Les vases athéniens à figures rouges, La période archaïque} ; première édition sous le titre \italique{Athenian Red Figure Vases, The Archaic Period}, 1975.)}
\label{Onesime}
\end{figure}



\Note{Note 16 aux \italique{Thesmophories} : vers 569}

\citComm{πρόσθες μόνον, κἀγώ σε νὴ τὴν Ἄρτεμιν...}\\

Il se peut que πρόσθες, dans le sens de πρόσφερε τὴν χεῖρα, ne soit pas le mot qui convient. Cobet propose πρόσιθι.

On se rapprocherait davantage, ce me semble, et du sens et de la tradition manuscrite, en lisant : πρόσθιγε. Cf. Soph., \italique{Phil.}, 817 : ἀπὸ μ᾽ ὁλεῖς, ἢν προσθίγῃς; Eur., \italique{Héracl.}, 652 ; \italique{Hipp.}, 1086; et pour une expression équivalente, Esch., \italique{Suppl}., 925 : κλαίοις ἂν εἰ ψαύσειας.

\Note{Note 17 aux \italique{Thesmophories} : vers 575-576}

\citComm{ὅτι μὲν φίλος εἴμ' ὑμῖν, ἐπίδηλος ταῖς γνάθοις·\\
	γυναικομᾶνω γὰρ προξενῶ θ᾽ ὑμῶν ἀεί.}\\

Cobet, suivi de M. van Leeuwen, supprime le premier vers, sous le prétexte que \italique{nisi deletur, γάρ quo referatur non habet}. De la part d'un helléniste si expert, un tel jugement a lieu de surprendre : γάρ est ici confirmatif, plutôt qu'explicatif, et se traduirait en français par \italique{en effet}, \italique{oui}\note{M. van Leeuwen est tombé dans la même erreur, \italique{Oiseaux}, v. 1344.}. Voici une phrase de Platon qu'on dirait calquée sur celle d'Aristophane : ὅτι δ᾽ ἐγὼ τυγχάνω ὤν τοιοῦτος, ἐνθένδε ἂν κατανοήσαιτε· οὐ γὰρ ἀνθρωπίνῳ ἔοικετὸ ἐμὲ τῶν ἐμαυτοῦ ἁπάντων ἠμεληκέναι. (\italique{Apol. Socr}., 31a). Tout un paragraphe de la grammaire de Kühner-Gerth est consacré à cet emploi spécial de γάρ (§545, 3). Et il y a près d'un siècle et demi que Larcher, l'admirable helléniste si méconnu quoique si souvent pillé\note{Exemples : Macrobe nous a conservé le vers suivant de Cratinos : στολὴν δὲ δή τίν᾽ εἶχε τουδομοιφρασον. Dans des \italique{Remarqres critiques sur l'Etymologicum Magnum}, parues en 1791 (t. XLVII des \italique{Mémoires de l'Académie des Inscriptions} ; la date est à noter), Larcher corrige : τοῦτό μοι φράσον. Ouvrez la collection des \italique{Fragments des Comiques}, vous y trouverez cette correction sous le nom de Porson. Et, en effet, le savant anglais la donne comme sienne dans son commentaire sur l'\italique{Hécube}, v. 782. Un autre vers de Cratinos a été incorrectement transcrit par le scholiaste de Pindare sur \italique{Ol.}, VI, 152) : οὗτοι δ᾽εἰσὶν συοβοιωτικὸν πεζοφόρον γένος ἀνδρῶν. Larcher restitue le vrai texte :  συοβοιωτοὶ κρουπεζοφόρον. Porson s'approprie aussi cette heureuse correction, avec le commentaire qui l'accompagne, dans la préface de l'\italique{Hécube,} p. LVII, et depuis c'est encore à lui qu'on en fait honneur. Il m'en coûte, je l'avoue, d'avoir mis en cause un savant que je croyais à l'abri de pareils soupçons. Mais le larcin est indéniable, et facile à vérifier. C'est ainsi que, pour le dire en passant, plus d'une émendation sur Aristophane porte le nom de Meineke, qui est de Boissonade.}, a attiré l'attention sur ce point dans une note de sa traduction de \italique{Chéréas et Callirrhoé} (t. I, p. 204, de l'édit. de 1823).

\Note{Note 18 aux \italique{Thesmophories} : vers 592-593}

\citComm{\decalage{ἠλίθιος, ὅστις τιλ}τίς δ᾽οὕτως ἀνὴρ\\
	ἠλίθιος, ὅστις τιλλόμενος ἠνείχετο;}\\

Brunck lit ἠνείχετ᾽ ἄν, car il était de ces éditeurs — combien nombreux, même aujourd'hui — dont c'est le dada de subtiliser sur la particule ἄν et de la fourrer où elle n'a que faire. Tel est bien le cas ici. Un orateur demandait : τίς οὕτως εὐήθης ἐστὶν ὑμῶν ὅστις ἁγνοεῖ; « Qui de vous est assez simple pour ignorer... ?\note{Démosth., \italique{Ol}., I, §15. — Ce qui équivaut à :οὐκ ἐστιν οὕτως εὐήθης ὅς ἀγνοεῖ. Cf. Soph., \italique{Antig.} ,220} » 
Dès lors on ne voit pas ce qui pouvait empêcher Mnésiloque de dire : τίς οὕτως ἀνὴρ ἠλίθιος ὅστις ἀνέχεται; « Quel homme est dépourvu de sens au point de se laisser faire ? » ou s'il s'agit du passé : ὅστις ἠνείχετο « au point de s'être laissé faire ? » 
G. Hermann, un spécialiste, et M. Sobolewski ont donc eu raison de défendre la vulgate. 

M. van Leeuwen rétablit ἄν. Et d'abord il se fait beau jeu en traduisant : \italique{Quis est vir ille stultissimus qui}, au lieu de : \italique{Ecquis est vir tam stultus qui}. Puis il ajoute : \italique{Non momen viri audire cupit sed negat reperiri posse virum tam stolidum}. 
D'accord, seulement je prie l'éminent professeur de se persuader que quand Mnésiloque dit (je traduis littéralement) : « Quel est l'homme assez sot pour s'être laissé épiler ? » il ne songe pas un instant à demander le nom de l'homme. En grec, comme en français, j'en atteste Démosthène, c'est une façon de nier qu'il puisse se rencontrer un homme de cette espèce. 

Je renverrai volontiers M. van Leeuwen à sa note sur le v. 872, où Brunck avait également faufilé ἄν mal à propos. Citant à l'appui du \italique{Ravennas} le v. 1172 de l'\italique{Œdipe à Colone}, il écarte ἄν par la même considération que je viens de faire valoir : \italique{verba sunt negantis potius quam rogantis}.

\Note{Note 19 aux \italique{Thesmophories} : le chœur compris entre les vv. 655-686}

Ce petit chœur n'a pas été examiné avec assez d'attention. On n'a pas remarqué qu'il s'accompagne d'un jeu de scène unique dans l'œuvre d'Aristophane. Trois morceaux le composent : \\
\ \ \ 1° Les vv. 655 à 658 appartiennent au coryphée. Mnésiloque n'est peut-être pas le seul homme qui ait pénétré parmi les femmes. Il importe de s'en assurer, et, après avoir déposé les manteaux et allumé les torches, de « tourner et virer tout à l'entour du Pnyx, et d'inspecter avec soin les tentes et les passages ».\\
\ \ \ 2° Les quatre suivants (659-662) sont des trochées, et le changement de mètre indique un changement d'interlocuteur. C'est, en effet, la réponse du \italique{parastate} ou chef du demi-chœur. En disant qu'il y avait lieu d'inspecter le Pnyx, les tentes et les passages, la coryphée indiquait clairement qu'il fallait quitter l'orchestre, c'est-à-dire l'enceinte sacrée adjacente au temple\note{On peut conclure de là que le Thesmophorion, qu'on savait situé sur une hauteur, l'était sur la colline du Pnyx et voisin du lieu d'assemblée.}. La parastate annonce donc que le chœur va « s'élancer dehors (ἐξορμᾶν) d'un pied léger\note{Le simple, ὀρμᾶν, qui se lit plus loin, v. 953, et \italique{Gren}., 478, signifierait « prendre sa course ».} ». Puis elle engage celle qu'elle appelle τὴν πρώτην à commencer tout de suite la reconnaissance. Comment faut-il interpréter τὴν πρώτην ? D'y voir, comme on le fait généralement, une locution adverbiale (\latin{quam primum} ou \latin{prima quaque via}) me semble inadmissible. Avec M. Mazon, je crois que ce terme désigne une personne déterminée\note{\italique{Essai sur la composition des comédies d'Aristophane}, Paris (1904), p. 131, n. 4.} ; suivant moi, ce serait τὴν πρώτην ἱέρειαν, appelée aussi τὴν προστάτιν, savoir la première, il y en avait beaucoup\note{Il y en avait deux par dème (Isée, \italique{Succession de Kiron}, 19-20), donc plus de deux cents.}, des prêtresses temporaires, ἀρήτειραι, de Déméter, celle qui, tout à l'heure, avait prononcé les prières d'usage (vv. 295 et suiv.), en un mot, la coryphée. 
Callimaque, parlant d'une de ces prêtresses, ἱερέη Δήμητρος, se sert d'une désignation analogue : πολλῶν προστασίη νέων γυναικῶν\note{\italique{Anthol. Palat}., VII, 728.}. \\
\ \ \ 3° Le morceau qui suit (vv. 663-684) est un \italique{hyporchême}. Une partie du chœur reste en place et chante, tandis que l'autre, en quête de quelque complice de Mnésiloque\note{εἴ τις ἐν τόποις ἑδραῖος. Qu'avons-nous affaire d'ajouter, avec MM. Blaydes et van Leeuwen, τούτοισι τοῖς devant τόποις ? La même locution adverbiale se lit dans un fragment de Nausicratès : [τρίγκλαι] ἃς κλύδων Αἰξωνικὸς ἐν τόποις παιδεύεται (\italique{Fragm. Com. Gr}., t. IV, p. 575 [= Athénée 325e = Nausicratès, fr. 1.6-11 Kassel-Austin]) et la locution inverse se lit dans Sophocle : ἥτις μ᾽ ἀποικίσειεν ἐκ τόπων (\italique{Trach}., 955), où jamais personne n'a trouvé à redire (ἐκ τούτων τῶν τόπων.). Je ne vois pas que ces tournures soient plus bizarres que ἐν ποσί et ἐκ ποδῶν, ou bien que ἐν χώρᾳ et ἐκ χώρας, πρὸ ou παρὰ χειρῶν, etc.}, évolue silencieusement, σιωπῇ (v. 660), va et vient, entre et sort, cherchant partout, dans l'orchestre et dehors, καὶ τὰ τῇδε καὶ τὰ δεῦρο (v. 666), jusqu'à ce qu'il déclare n'avoir rien trouvé (687-688).\\


\Note{Représentation d'un ἔγκυκλον aux franges pourpres ?}

\begin{figure}[!h]
\centering
\includegraphics[scale=0.45]{Encycle.png}
\caption{Vase n° 204020, \italique{Corpus Vasorum Antiquorum}, Beazley Attic Red-Figure Vase-Painters\expo{2} (378.122)}
\label{Encycle}
\end{figure}


\Note{Une crocote que la pudeur réprouve...}

\begin{figure}[!h]
\centering
\includegraphics[scale=0.65]{Crocote transparente.jpg}
\caption{Cratère « St. Agata Antigone », St. Agata dei Goti I, collection Rainone-Mustilli. Inventorié par Arthur D. Trendall , \italique{Phlyax Vases}, « BICS » suppl. LXIX (1967).}
\label{Encycle}
\end{figure}


\Note{Note 20 aux \italique{Thesmophories} : vers 656}

\citComm{ξυζωσαμένας εὖ κἀνδρείως τῶν θ᾽ἱματίων ἀποδύσας}\\

M. van Leeuwen déclare le passage corrompu, vu que, dit-il, ἀποδῦναι \italique{aut absolute ponitur aut accusativum adsciscit}. Cela est-il bien sûr ? Si l'on dit couramment ἱμάτιον ἀποδῦναι, ou simplement ἀποδῦναι, qu'y a-t-il d'incorrect dans ἀποδῦναι τοῦ ἱματίου ? Il existe deux synonymes de ἀποδύεσθαι, savoir γυμνοῦσθαι et ψιλοῦσθαι. L'un et l'autre se prennent aussi absolument. Mais, comme tous les verbes marquant absence ou privation, ils se construisent régulièrement avec le génitif : ainsi l'on trouve γυμνώθη ῥακέων (Hom., \italique{Od.}, 22, 1) et ἐψιλώθην τῶν ἐνδυμάτων (Alciph., 3, 42). 
De même en latin, on dit à la fois \italique{exuo vestem} et \italique{exuor veste} ; et en français \italique{dépouiller son vêtement} ou \italique{se dépouiller d'un vêtement}. 

Que si nous n'avons pas d'autre exemple de ἀποδῦναι τινός, que nous importe ? L'analogie suffit. Il ferait beau voir les éditeurs obligés de citer un texte à l'appui de chaque mot ou de chaque tournure. On oublie trop souvent que sur peut-être trois cents comédies grecques tenues pour classiques dans l'antiquité, nous n'en possédons que onze. 

M. van Leeuwen de son autorité substitue dans le texte ἀμελούσας à ἀποδύσας. C'est bien ici qu'un exemple eût été de rigueur, car on ne dispose pas des mots suivant sa fantaisie. Mais aurions-nous sous la main toute la littérature grecque d'Homère à Quintus de Smyrne, qu'il faudrait désespérer, je le crains, de découvrir un endroit où ἀμελεῖν τοῦ ἱματίου  aurait le sens de \italique{déposer son manteau}.

\Note{Note 21 aux \italique{Thesmophories} : vers 710-712}

\citComm{ἀλλ᾽οὖν ἥχεις ὅθεν ἥχεις,\\
	φαύλως τ᾽ ἀποδρὰς οὐ λέξεις\\
	οἷον δράσας διέδυς ἔργον...}\\

Ce passage a été refait de cinq ou six façons, plus bizarres les unes que les autres. La correction me paraît élémentaire. Les manuscrits concordent, sauf sur le premier ἥκεις (ἥκεις γ᾽ dans Suidas, ἥκεις τ᾽ dans le \italique{Ravennas}, ἥκισθ᾽ dans le manuscrit de Munich). 
Lisez ἥκων γ᾽, le sens devient parfait :

\citComm{ἀλλ᾽ οὖν ἥκων γ᾽ ὅθεν ἥκεις...}\\

Il est à peine besoin d'ajouter que ἥκων est pris dans le sens de ἀνήκων, \italique{redux}, ce qui est des plus fréquents\note{\italique{Paix}, 265, 275 ; \italique{Lysistr.}, 595, 731, 924 ; \italique{Gren.}, 503 ; \italique{Ploutos}, 521, 1179, etc.}. Et quant à ἀλλ᾽ οὖν .. γε, chacun sait qu'il signifie \italique{toujours}, \italique{tout au moins}, \italique{en tous cas}\note{\italique{Nuées}, 985, 1002 ; \italique{Guêpes}, 1129, 1190 ; Esch., \italique{Prom}., 1058 ; Soph., \italique{Él.}, 233 ; \italique{Ant.}, 84, etc}. Traduisez :

\citComm{Toujours tu n'iras pas, revenu au lieu dont tu es parti, après une fuite facile, raconter ton crime et ton évasion...}\\

C'est le mot d'Epicharme, fr. 126 : κἀπῆνθεν ὅθεν ἦνθεν πάλιν; ou de Ménandre (t. IV, p. 211 [$=$ fr. 373(416\expo{a}).3 ed. Kassel-Austin]) : ἀπῆλθεν ὅθεν ἦλθεν ταχύ. Cf. aussi Soph., \italique{Œd. Col}., 1227 : βῆναι κεῖσ᾽ ὁπόθεν περ ἥκει; ; Eur., \italique{Her. F.}, 1247 : θανὼν ὅθενπερ ἤλθον, εἶμι γῆς ὕπο, etc.

\clearpage

\Note{Papyrus d'Oxyrhynque n° 3839 : vers 742-766 (fr. 1) ; vers 941-953 (fr. 2)}


\begin{figure}[!h]
\centering
\includegraphics[scale=0.63]{POxy3839.jpg}
\caption{\italique{The Oxyrhynchus Papyri}, vol. LVI (1989), n°3839. (Ed. Walter E. H. Coockle.)}
\label{POxy3839}
\end{figure}

\pagebreak

\begin{figure}[!h]
\hspace*{-1cm}
\centering
\includegraphics[scale=0.50]{POxy3839fr1Coockle t.png}
\caption{\italique{The Oxyrhynchus Papyri}, vol. LVI (1989), n°3839. (Ed. Walter E. H. Coockle.) Fragment n° 1 vers 742-766 (texte).}
\label{POxy3839fr1t}
\end{figure}

\begin{figure}[!h]
\hspace*{-1cm}
\centering
\includegraphics[scale=0.55]{POxy3839fr1Coockle c.png}
\caption{\italique{The Oxyrhynchus Papyri}, vol. LVI (1989), n°3839. (Ed. Walter E. H. Coockle.) Fragment n° 1 vers 742-766 (commentaire).}
\label{POxy3839fr1c}
\end{figure}

\pagebreak

\begin{figure}[!h]
\hspace*{-1cm}
\centering
\includegraphics[scale=0.50]{POxy3839fr2Coockle t.png}
\caption{\italique{The Oxyrhynchus Papyri}, vol. LVI (1989), n°3839. (Ed. Walter E. H. Coockle.) Fragment n° 2 vers 941-953 (texte).}
\label{POxy3839fr2t}
\end{figure}

\begin{figure}[!h]
\hspace*{-1cm}
\centering
\includegraphics[scale=0.55]{POxy3839fr2Coockle c.png}
\caption{\italique{The Oxyrhynchus Papyri}, vol. LVI (1989), n°3839. (Ed. Walter E. H. Coockle.) Fragment n° 2 vers 941-953 (commentaire).}
\label{POxy3839fr2c}
\end{figure}

\pagebreak


\Note{Note au vers 83 des \italique{Nuées}, ou l'encombrement inutile de la scène.}

\citComm{83 \decalage{aaaaaaaaa}νὴ τὸν Ποσειδῶ τουτονὶ τὸν Ἵππιον.}

Οὗτος ou οὑτοσί (on verra par les exemples qu'il n'y a pas de différence à faire) se prend très souvent dans un sens emphatique, en parlant de personnes ou de choses insignes dans leur genre. Alors il se traduit exactement en latin par \italique{ille}, en français par \italique{ce fameux}, en bonne ou mauvaise part. 

En voici quelques exemples tirés d'Aristophane : \italique{Nuées}, 1403, Σωκράτης οὑτοσί (\italique{vafer ille Sisyphus}, Hor.) ; \italique{Gren}., 429, Καλλίαν τοῦτον τὸν Ἱπποβίνου (\italique{Xenophon Socraticus ille}, Cic.) ; \italique{Guêpes}, 592, ὁ μέγας οὗτος Κολακώνυμος (\italique{magno illi Alexandro simillimus}, Paterculus); \italique{Paix}, 53, τοῖς ὑπερηνορέουσι τούτοις ; \italique{Acharn.}, 427, Βελλεροφόντης ὁ χωλὸς οὑτοσί\note{On trouverait des exemples dans les classiques français. « Comme disait cet historien » (Quinte Curce). Bossuet, \italique{Or. fun. de la reine d'Angleterre}. -- « Il n'y avait qu'une de ces fleurs, a dit cet homme, en Europe » (un maître botaniste). 0. Feuillet, \italique{Rédemption}.}. 

Je me borne à ces échantillons. Si « Bellérophon le glorieux boiteux » est clair, « Poséidon le glorieux hippien » ne l'est pas moins\note{Τοῦτον, ὃν πάντες ἵσασι δηλονότι, dit une glose sur \italique{Gren.}, 429.}. Les deux textes sont parallèles et ne sauraient s'interpréter différemment. Où prendrait-on prétexte de cette distinction ? À moins qu'on ne dénie à la gloire du dieu le lustre qu'on accorde à celle du héros. 
En latin, on traduirait \italique{Neptunus ille equestris}. Et de fait, Plaute a dit, à l'exemple d'Aristophane : \italique{quod ille faciat Jupiter}\note{Fritzsche prétend que \italique{ille Jupiter} signifie « Jupiter qui siège là-haut ». Concédons-lui cela. Mais ne peut-on en dire autant d'Apollon et de Poséidon ? Dans la croyance des Grecs, les dieux ne sont-ils pas présents partout ? Pindare le déclare expressément à propos d'Apollon (\italique{Pyth.}, III, 27). Dans Euripide, Athéna dit, s'adressant à Oreste absent : κλύεις γὰρ αὐδὴν καίπερ οὐ παρὼν θεᾶς (\italique{Iphig. T}., 1447), ce que Plaute traduit textuellement, en le mettant dans la bouche de Jupiter : \italique{Audis quæ dico, tametsi præsens non ades}. (\italique{Amph}., 3, 22.) Vous serez bien avancé quand vous aurez rendu μὰ τὸν Ἀπόλλω τουτονί par « J'en atteste Apollon qui m'entend ».}. 

Quoiqu'il ait été relevé par les grammairiens (entre autres par Kühner-Gerth, § 467, 5), cet emploi spécial de οὗτος n'a été que trop souvent méconnu, et il en est résulté de très singuliers contresens. Ainsi dans notre passage la plupart des interprètes, Lobeck, Teuffel, Kock, Merry, Hickie, Blaydes, van Leeuwen, attribuent à οὗτος une idée d'indication précise et supposent qu'il y a dans un coin ou l'autre de la scène une statue de Poséidon, que l'acteur serait censé montrer du doigt.

Il en va de même d'un autre endroit d'Aristophane, où Apollon est invoqué de la même manière, μὰ τὸν Ἀπόλλω τουτονί (\textbf{\italique{Thesmophories}, v. 748}), non pas cette fois dans la rue, mais dans l'enceinte même du Thesmophorion. Suivant certains de ces éditeurs, ces désignations se rapportent à des idoles installées, non par occasion, mais à demeure, sur la scène. En d'autres termes, ils se figurent celle-ci comme une sorte de plate-forme ornée de colonnes et de statues, quelque chose comme la tribune des Rostres à Rome. Sur ce décor permanent venait s'appliquer le décor spécial à chaque pièce. Comment l'illusion scénique s'accommodait de ce double décor, dont on ne conçoit ni la convenance ni l'utilité, c'est ce dont on ne s'inquiète guère, et ce qu'il fallait pourtant commencer par expliquer. Demandez-vous quel visage devait faire cet Olympe de dieux sur le Caucase, par exemple, en face du calvaire de Prométhée, au bord de la mer, au milieu d'un camp grec (\italique{Hécube}, \italique{Ajax}, \italique{Troyennes}, etc.), ou encore devant la ferme de Trygée (\italique{Paix}), ou parmi les rochers et les halliers servant de cadre au nid de Térée (\italique{Oiseaux}). 

Qu'il y eût des cas où les statues trouvaient leur emploi, qui songe à le nier ? Ainsi quand le décor de fond représentait un palais, souvent le péristyle contenait une ou plusieurs effigies de dieux : Artémis et Cypris, dans l'\italique{Hippolyte} d'Euripide ; Apollon, dans l'\italique{Électre} de Sophocle. De même lorsque la scène figurait une rue, on y mettait un Hermès. Quoi d'étonnant ? Les choses ne se passent guère autrement de nos jours. À Athènes les frais de ce genre incombaient au fermier du théâtre ou plus probablement au chorège. L'auteur de ce bizarre paradoxe sur la décoration scénique, c'est Fritzsche. Sur la foi de témoignages équivoques ou mal compris, il se flatte d'avoir démontré qu'en tout cas une statue d'Apollon faisait partie intégrante de la scène, peut-être aussi, mais il n'en est pas aussi sûr, une autre d'Athéna. De Poséidon il n'a pas trouvé trace dans les textes, mais cela ne l'empêche pas d'affirmer que dans les \italique{Nuées} une statue du dieu, « et splendide », il n'en fait pas doute, se dressait devant la demeure de Strepsiade. 
Un autre monument, un Dinos, faisait pendant devant le logis de Socrate. Ajoutez-y l'Hermès dont il est question au v. 1478 ; de compte fait, voilà au moins cinq statues, sans parler de celles que nous ne connaissons pas. 

Le digne savant est si convaincu de l'excellence de sa théorie, qu'il traite de menteur effronté le scholiaste citant comme une locution courante : μὰ τοῦτον τὸν Ἀσκληπιόν (\italique{quae sunt impudenter mentientis}), et qu'il déclare brouillé de naissance avec la muse comique (\italique{natus musa comicorum irrita}) celui qui s'aviserait d'y redire. 

Lui se contente d'être brouillé avec la logique et le sens commun, car sa théorie ne tient pas debout. Cette dispendieuse idole de Poséidon, qui donc est censé en avoir fait les frais ? Ce n'est pas Strepsiade apparemment, qui depuis son mariage a pris ce dieu en grippe. 
Supposé d'ailleurs que ce soit lui, comment se trouve-t-elle là ? Car Strepsiade, nous le tenons de sa propre bouche, habite à la campagne, bien loin d'Athènes. Il faudrait donc présumer qu'il l'a emportée dans ses bagages, quand la guerre l'a forcé à se réfugier en ville. Accordons encore cela, mais au moins faut-il qu'on nous dise où il l'aurait placée cette statue. Dans la chambre à coucher ? Alors les spectateurs n'auraient pu l'apercevoir. Devant sa maison ? A qui ferez-vous accroire qu'il fût permis d'encombrer la voie publique d'un monument ? Et quand on ne ferait compte d'aucune de ces objections, semble-t-il présumable que le poète ait imposé à son chorège ce surcroît de dépense pour l'unique plaisir de mettre un démonstratif à la suite de la formule très claire par elle-même μὰ τὸν Ποσειδῶ ? 

Plus plaisante encore est l'histoire du Dinos. Socrate, d'après Aristophane, ne croit pas aux dieux. 11 a donc été amené à expliquer à Strepsiade que l'univers est régi, non par Zeus, mais par un tourbillon. De sa part cela s'entend, tourbillon, en grec Δῖνος, étant de ces termes vagues, chers aux idéologues et aux inventeurs de cosmogonies (Anaxagore, Démocrite, Descartes, etc.). Mais δῖνος est aussi le nom d'un vase à boire : il en est question dans une autre comédie\note{\italique{Guêpes}, v. 618.}. Le naïf Strepsiade ne sait pas trop duquel des deux Socrate a voulu parler. Plus tard il avouera sa sottise et regrettera de s'être laissé mettre dedans « à cause de ce maudit Dinos », διὰ τουτονὶ τὸν Δῖνον (v. 1473). Et il ajoute cette apostrophe que l'on a suspectée, j'ignore pourquoi, car elle est aussi claire en grec qu'en français: « Malheureux que j'étais, dire que je te prenais, vile poterie, pour un dieu ! » 

Fritzsche, comme si Socrate était solidaire du quiproquo, a imaginé de caser devant la demeure du philosophe, en guise d'enseigne et de symbole, un Dinos en terre cuite, invention dont il se montre particulièrement glorieux. 

Pour moi, dût la muse comique me renier à jamais, j'avoue que l'erreur de Strepsiade me semble excusable au prix de la sienne, et que des deux interprètes de la doctrine de Socrate, ce n'est pas Strepsiade qui me paraît le moins intelligent. 

M. van Leeuwen n'a eu garde de tomber dans pareille méprise. Mais comme il s'en tient au sens primitif de οὗτος, il n'a trouvé d'autre moyen de se tirer d'affaire que de lire : 

\citComm{\decalage{διὰ τουτονὶ τὸν δῖνον.}ἀλλ᾽ ἐγὼ τοῦτ᾽ ᾠόμην\\
(Subito domum intrat, unde statim redit vas rotundum afferens :)\\
διὰ τουτονὶ τὸν δῖνον.}\\

Voilà en effet ce qui coupe court à toute difficulté. C'est comme qui dirait : « Spectateurs, vous savez comme moi ce que c'est qu'un Dinos ; n'importe, je vais vous chercher l'ustensile en question : il n'est rien de tel que de voir les choses. » 
Seulement soyons logiques jusqu'au bout, et quand, dans \italique{Lysistrata} (v. 1108), nous entendrons le délégué d'Athènes réclamer, en échange de Pylos, la ville d'Echinos, le golfe Maliaque et Mégare, ne manquons pas de mettre :

\citComm{τὸ δεῖνα τοίνυν παράδοθ᾽ ἡμῖν\\
(Subito domum abit, unde non ita multo post redux ingentem tabulam affert :)}
\citComm{\decalage{τὸ δεῖνα τοίνυν παράδοθ᾽ ἡμῖν}τουτονὶ\\
πρώτιστα τὸν Ἐχινοῦντα καὶ τὸν Μηλιᾶ...}\note{J'aurais pu citer aussi bien, quelque six vers plus haut (1162), l'endroit où l'envoyé de Sparte réclamait l'encyclos, c'est-à-dire Pylos. Si je ne l'ai pas fait, c'est que l'intervalle entre les mots τοὔγκυκλον et τοῦτο eût été tel qu'il aurait fallu renvoyer àquatre jours la suite du vers et de la pièce.}

\clearpage

\Note{Le cratère de Würzburg et la scène d'égorgement de l'outre : vers 754-757}

\begin{figure}[!h]
\centering
\includegraphics[scale=0.17]{Wurzburg.jpg}
\caption{Cratère apulien d’environ 370 av. J.-C., Martin von Wagner Museum der Universität Würzburg, Antiken- Sammlung, H5697.}
\label{Wurzburg}
\end{figure}


\Note{Note 22 aux \italique{Thesmophories} : vers 773-774}

\citComm{τί δ᾽ ἄν, εἰ ταδὶ τἀγάλματ᾽ ἀντὶ τῶν πλατῶν\\
	γράφων διαρρίπτοιμι}\\

Faute d'avoir déterminé le sens du mot ἀγάλματα, on s'est cruellement mépris jusqu'ici, non seulement sur les deux vers ci-dessus, mais sur l'épisode tout entier. 

Mnésiloque, dit Fritzsche, se sert en guise de tablettes des bustes et statues de dieux, qui sont en grand nombre dans le Thesmophorion, et les disperse de tous côtés. Pour le scho liaste, ces tablettes ne sont autres que des πινάκια ou petits tableaux votifs suspendus dans le temple. 
Les traducteurs se sont partagés entre ces deux interprétations également insoutenables. \\
\ \ \ 1° \ \ De tableaux votifs, il n'en peut être question, car la scène se passe, non pas à l'intérieur du temple, mais dans le τέμενος, ou enceinte sacrée. C'est le décor de fond qui représente le ἱερόν ou sanctuaire, lequel reste clos durant toute la pièce. \\
\ \ \ 2° \ \ Supposer la présence d'un certain nombre de bustes et statues de dieux n'est guère plus admissible. D'abord, parce que cela ne se concilie point avec la simplicité de la mise en scène chez les Grecs. Songez que les jours de représentation on donnait coup sur coup cinq pièces sans entr'actes. Comment s'y serait-on pris, si pour chaque pièce, et à chaque changement à vue, il avait fallu, comme de nos jours, encombrer la scène d'une quantité d'accessoires ? 
Qui admettra d'ailleurs qu'on écrive sur des statues et des bustes et qu'on les lance de toutes parts ? Au contraire le texte indique clairement qu'il s'agit d'objets légers et d'un maniement facile (vv. 783 et suiv.). \\
\ \ \ 3° \ \ Reste une dernière considération, primant toutes les autres. On a vu dans la scène précédente Mnésiloque, se voyant perdu, se souvenir fort à point d'un épisode du \italique{Télèphe} d'Euripide, qui se résume ainsi. Le héros a pénétré dans le camp grec sous un déguisement. Quand sa ruse est découverte, il se saisit du petit Oreste, et, réfugié sur un autel, menace d'égorger l'enfant avec le coutelas sacré. La situation étant à peu près pareille, Mnésiloque se décide à faire comme Télèphe ; chose d'autant plus facile qu'il a devant lui l'autel, situé, selon l'usage, au pied des degrés qui conduisent au sanctuaire. 
De là il bravera les femmes acharnées à sa perte. En effet, si criminel qu'il fût, celui qui se réfugiait sur un autel se rendait suppliant, ἱκέτης, et se mettait sous la protection des dieux. L'en arracher était un sacrilège. 

Il existait, il est vrai, un moyen d'ôter la vie au suppliant, ou de le faire déguerpir, c'était d'entourer l'autel d'un bûcher et d'y mettre le feu\note{Eurip., \italique{Herc. Fur.}, 243 et 285 ; \italique{Androm.}, 256, et le schol. : « La coutume était d'apporter du feu près des suppliants réfugiés sur l'autel, afin de les contraindre à le lâcher ». Il en était autrement des suppliants réfugiés dans un temple. Ici le feu n'eût servi de rien. Le seul moyen de se débarrasser d'eux était de murer les portes. Cf. Thucyd., I, 134, et III, 81.}. 
Ce moyen est trop classique pour que les femmes n'y aient pas songé. Aussi y songent-elles. Pourquoi elles y renoncent ? C'est, à n'en pas douter, qu'elles se sont dit que cette façon de brusquer le dénouement ne serait pas du goût du public, qui verrait son plaisir sitôt fini, et moins encore de l'acteur ainsi offert en holocauste. 

C'est à ce point de l'action que nous sommes arrivés. Plus tard, nous verrons Mnésiloque commettre l'imprudence de suivre Euripide déguisé en Ménélas, et du coup il sera entouré et livré aux autorités. Mais pour le moment il n'a garde de quitter son asile, dont il ne bougera de longtemps (v. 886). 
Son unique souci est d'envoyer un message à son gendre. Comment faire ? 

Une autre tragédie d'Euripide en fournira le moyen. Palamède ayant été méchamment occis par les Grecs devant Troie, son frère Œax a imaginé d'écrire sur des avirons pour faire savoir la nouvelle à Nauplios, père des deux héros, qui habite en Eubée. Sans doute il est fâcheux que l'autel des deux déesses soit dépourvu d'avirons ; mais qu'à cela ne tienne, on n'est pas des familiers d'Euripide pour rien. Faute d'avirons, on se servira des ἀγάλματα. 

Que sont-ce donc que ces ἀγάλματα ? La définition souvent citée : πᾶν ἐφ᾽ ᾧ τὶς ἀγάλλεται est la meilleure qu'on ait donnée du mot. 
Ainsi, tout ce qu'on offre et consacre aux dieux peut être nommé ἄγαλμα. 
Dans Homère, le cheval de Troie est un ἄγαλμα θεῶν θελκτήριον « une offrande destinée à apaiser les dieux » ; un taureau destiné au sacrifice en est un autre\note{\italique{Odyssée}, VIII, 509 ; III, 438.}. Aristophane, parodiant Euripide, appelle le chien Ἑκάτης ἄγαλμα\note{\italique{Fragm.}, 535, Dind. [$=$ fr. 608(594) Kassel-Austin].}, car c'est la victime que la déesse préfère, d'où son nom de κυνοσφαγής. Un comique nous montre une divinité, probablement Apollon, qu'on cherche à se rendre propice par la promesse « de petites offrandes et d'une branche de laurier chaque premier jour du mois\note{Ἀγαλματίοις ἀγαλοῦμεν αἰεὶ καὶ δάφνῃ. Théopompe, t. II, p. 810, Mein. [$=$ fr. 48(47) Kassel-Austin] À rapprocher de \italique{Paix}, 399 : 6}] ». 
Enfin dans Platon, Socrate, arrivé à l'endroit où il va s'entretenir avec Phèdre, au bord d'une source fraîche ombragée d'un platane, remarque que le lieu « semble consacré à quelques nymphes et à l'Archéloüs, à en juger par les figurines et les offrandes » ἀπὸ τῶν κορῶν τε καὶ ἀγαλμάτων\note{\italique{Phèdre}, 230b.}. 
En un mot, rien n'est plus exact que la définition du lexique de Timée : ἄγαλμα, πᾶν ἀνάθημα. 

Maintenant, prenons attention que l'autel où s'est assis Mnésiloque est couvert d'offrandes. Lui-même en entrant y a déposé τὸ πόπανον « le gâteau d'usage » (v. 285), et il va de soi que les choreutes, sinon tous, du moins la plupart, en ont fait autant. Une autre preuve se tire du passage où, tenant dans ses bras l'enfant, il menace de le frapper ἐνθάδ᾽ ἐπὶ τῶν μηρίων « ici même sur les restes des victimes » (v. 693). 
De victimes, à proprement parler, il n'y en a pas, et il ne peut y en avoir, ce n'est pas le moment\note{Cf. la scholie sur le v. 376. -- \italique{Oiseaux}, v. 1519.}. Mais le public a reconnu de suite que c'est Télèphe qui parle par la bouche de l'acteur, et que les victimes sont représentées par de simples gâteaux.

On voit maintenant ce qu'Aristophane désigne par les ἀγάλματα ; et pour peu qu'on entre dans le sens du poète comique, on conviendra qu'il ne pouvait dans la circonstance faire choix d'un terme plus heureux. Ces prétendues offrandes ne sont en réalité que ce qu'en termes de théâtre on appelle des accessoires. Chez nous les accessoires sont faits de carton ; à Athènes ils étaient de bois. Le comique s'amuse à le faire remarquer : « Elles sont de bois après tout; les rames (d'Œax) elles aussi étaient en bois ». 

Jeu d'esprit qui en rappelle un autre tout pareil dans les \italique{Oiseaux} (v. 901), et sur lequel les interprètes se sont également mépris. Cette fois c'est un bouc qu'on apporte pour le sacrifice, mais le bouc est postiche. Et le chœur de dire : « Il faut que j'invoque les dieux, — un d'eux sans plus..., car en fait de victime, il n'y a guère là que de la barbe et des cornes ». De même dans la Paix, le mouton sera immolé dans la coulisse et l'on n'apportera sur la scène que les cuisses, « afin que l'animal demeure au chorège » (v. 1022)\note{Lire là-dessus Dindorf, qui n'y a rien compris. [Drôle de conseil !..]}.


\Note{Note 23 aux \italique{Thesmophories} : vers 808-809}

\citComm{ἀλλ᾽ Εὐβούλης τῶν πέρυσίν τις βουλευτής ἐστιν ἀμείνων,\\
	παραδοὺς ἑτέρῳ τὴν βουλείαν; οὐδ᾽ αὐτὸς τοῦτό γε φήσει.}\\

Passage très controversé, bien qu'il ne paraisse guère difficile à entendre. On se souvient des réformes constitutionnelles qui avaient marqué l'hiver de l'année d'avant (πέρυσίν), à la nouvelle du désastre de Sicile. La Boulè n'offrant plus assez de garanties de prudence, « on avait décidé d'élire un Conseil de vieillards (les Probules) qui, suivant les circonstances, donneraient un avis préalable sur toutes les affaires\note{Thucydide, VIII, 1, 4.} ». 
C'était délivrer au Conseil des Cinq-Cents un certificat d'incapacité, le droit de formuler un avis préalable étant précisément l'essentiel de sa prérogative. 
Aussi, quand il sortit de charge en skirophorion (juin) 412, n'existait-il plus que de nom. 
Et si, ce que nous ignorons, l'usage prévalait déjà, comme au temps de Démosthène\note{Deuxième hypothèse du discours de Démosthène, \italique{C. Androtion}, 22, p. 590.}, de décerner à la Boulè démissionnaire une couronne d'or pour prix de ses services, il est bien sûr que cette fois nul ne se leva à l'Ecclésie pour proposer pareil décret. 

M. M. Croiset a fort bien vu que c'est à cette transmission annuelle de pouvoirs (παράδοσις, c'est le terme consacré), par laquelle chaque bouleute sortant cédait sa place à son successeur, que le chœur fait allusion\note{\italique{Aristophane et les partis à Athènes}, p. 238.}. 

J'ajouterai que τις n'est pas l'interrogatif, mais l'enclitique, se rapportant à βουλευτής\note{La phrase correspond exactement à cette autre d'Euripide : ἐν δ' εὐπροσηγόροισιν ἐστί τις χάρις; (\italique{Hipp.}, 95).}. Et c'est une erreur de prétendre, avec von Velsen et M. van Leeuwen, que φήσει ne peut se rapporter à l'adjectif indéfini τις. Je me bornerai à citer l'exemple suivant tiré de l'\italique{Ecclésie des femmes}, 899 : οὐδέ τις στέργειν ἂν ἐθέλοι μᾶλλον ἢ ᾿γὼ τὸν φίλον... ἀλλ᾽ ἐφ ἕτερον ἂν πέτοιτο.. Il serait facile de multiplier les exemples. Traduisons donc : 

\citComm{Et Euboulé, est-il un des bouleutes de l'an dernier, ayant résigné à un autre sa fonction, qui l'emporte sur elle ? Lui-même s'en défendra.}

\Note{Note 24 aux \italique{Thesmophories} : vers 811-812}

\citComm{οὐδ᾽ ἂν κλέψασα γυνὴ ζεύγει κατὰ πεντήκοντα τάλαντα\\
	ἐς πόλιν ἔλθοι τῶν δημοσίων.}\\

M. van Leeuwen prétend sans ombre ni apparence que ce trait est dirigé contre les Probules. Rapprochez-le, dit-il, du v. 421 de \italique{Lysistrata}, et concluez hardiment, \italique{tute}, qu'aux yeux d'Aristophane ce collège se composait de voleurs. 

Nous venons devoir dans quelles circonstances les Probules avaient été créés. Mais il importe de citer la phrase où Thucydide nous dit « qu'il fut décidé en même temps qu'on équiperait une flotte, en rassemblant des bois et de l'argent » (VIII, 1). 
Or c'est précisément ce que le Probule déclare dans \italique{Lysistrata} : « J'ai à fournir des bois de marine, et maintenant qu'il faut de l'argent, les femmes me ferment au nez les portes de l'Acropole » (où était le trésor). 
Comment a-t-il pu venir à l'idée de personne de voir dans ce passage, où le poète ne fait que confirmer l'historien, une insinuation contre la probité du magistrat ? Celui-ci joue dans la pièce un piteux personnage, dénué de tout prestige, de toute autorité, mais rien ne tend à le faire soupçonner, ni aucun de ses collègues, de malversation. 
M. van Leeuwen est bien le premier qui ait représenté ces vieillards, chargés de réprimer les abus de la démocratie, comme cherchant à capter la faveur populaire. 

Les voleurs dont parle Aristophane, nous savons où les trouver. Thucydide raconte qu'en apprenant la défection de Chios, les Athéniens jugèrent le péril si pressant qu'ils décidèrent de disposer des mille talents du trésor mis en réserve pour les cas d'urgence (VIII, 15). 
Cette mesure donna lieu à de graves détournements. Une dénonciation fut portée devant la Boulè contre un certain Philinos et trois autres, parmi lesquels un sous-greffier des Thesmothètes. Le procès eut lieu dans l'été de 412, peu avant la représentation. Ce fut Antiphon qui composa le plaidoyer contre les accusés, lesquels furent condamnés\note{Antiphon, \nieme{6} discours, 35 et 50.}. M. de Wilamowitz a fort ingénieusement mis ce procès en rapport avec notre passage, et sa démonstration parait tout à fait convaincante\note{\italique{Aristoteles und Athen.} Berlin, 1893, t. II, p. 346.}. 

M. van Leeuwen lit κλέψουσα (au lieu de κλέψασα) et rapporte à ce verbe le mot ζεύγει. Même il reprend pour son compte la remarque du Byzantin, expliquant que le voleur, pour emporter son butin, s'est servi de chars, dont chacun contenait cinquante talents. Mais alors pourquoi le singulier ζεύγει ? 
Et puis, que signifierait ce détail, alors qu'on nous laisse ignorer l'importance du vol ? Qu'on ait usé de camions, de chevaux ou de mulets, qu'importe au public ? Il n'est que trop clair qu'à une époque où la monnaie était presque exclusivement d'argent, il eût été difficile d'emporter sous son bras quelque trois cent mille francs. 

On nous avertit, il est vrai, que si nous repoussons la correction, le texte devient intraduisible ; car il ne saurait s'entendre, comme dans les \italique{Nuées} (v. 69), d'un vainqueur aux grands jeux montant à l'Acropole en tunique de pourpre. Qui songe à le contester ? Mais il est d'autres circonstances où l'Athénien pouvait se rendre en char à l'Acropole. 
Ainsi, chacun sait qu'à l'occasion des Panathénées, une merveilleuse procession partait du Pompéion, traversait l'Agora et montait au temple d'Athêna Poliade. 
On sait aussi que la partie du cortège qui attirait surtout les regards était la file brillante des chars. Participer à ce défilé était un luxe que seuls se permettaient les opulents. 
De même, lors des fêtes d'Éleusis, Démosthène nous montre la femme de Midias se rendant aux Mystères sur un char attelé de deux chevaux blancs\note{\italique{In Midiam}, p. 565, 27 ; aussi \italique{Ploutos}, v. 1013.}. Dès lors je ne vois pas ce qu'il y a à reprendre dans notre vers. Ne dirait-on pas, encore aujourd'hui, qu'on a vu tel financier véreux se faire conduire au grand prix de Paris dans un attelage à quatre ?

Quant à la construction, elle n'est guère plus enchevêtrée que \italique{Guêpes}, 717 (déjà cité par Dindorf), et \italique{Ibid.}, 604, \italique{Ois.}, 1150, \italique{Eccl.}, 1069, et vingt autres passages. Rien n'empêche d'ailleurs de lire avec M. van Herwerden : ζεύγει κλέψασα γυνή.


\Note{Note 25 aux \italique{Thesmophories} : vers 1019-1020}

\citComm{κλύεις ὦ προσαιδουσσαι τὰς ἐν ἄντροις\\
	κατάνευσον...}\\

Encore un vers qui a donné lieu à quantité de conjectures. Les principaux éditeurs, Dindorf, Bergk, Meineke, Fr. G. Wagner (\italique{Euripid. fragm}.) ont adopté celle de Seidler : πρὸς Αἰδοῦς σὲ τὰν. Elle a le mérite de se rapprocher du texte, mais le sens n'est pas satisfaisant : \italique{Pudori quis hic locus} ? dit avec raison M. van Leeuwen. Après quoi l'on s'étonne de voir l'éminent éditeur faire sienne une autre correction : πρὸς Ἅιδου σ᾽ ἀϋτὰς.... ἀπόπαυσον, laquelle s'écarte beaucoup plus du manuscrit et sou lève la même objection. 
À notre tour nous demanderons : \italique{Orco quis hic locus} ? 

Mieux valait reprendre, comme M. Blaydes, l'ancienne conjecture de G. Hermann : προσαυδῶ σὲ τὰν. C'est jusqu'ici la plus plausible, et l'on pourrait s'y tenir, si elle ne laissait à désirer au point de vue de la clarté : « toi qui habites dans des antres » est trop vague et ne désigne pas suffisamment l'Écho. 
Il y avait d'autres déesses et nymphes habitant dans des antres : par exemple Thétis, invoquée de la sorte dans l'\italique{Andromaque} : ὦ κατ᾽ ἄντρα μύχια Νηρέως κόρη (v. 1224), et de même les nymphes de montagnes, appelées ἀντροχαρεῖς (\italique{Oiseaux}, v. 1097). 
Or la précision est ici essentielle. Car c'est cet appel de Mnésiloque qui détermine Euripide, quoique travesti en Persée, à jouer dans la coulisse le rôle d'Écho. 

La leçon primitive se laisse dégager sans trop de peine. Thiersch et, après lui, Dobree remarquent fort à propos que προσαιδεῖν décèle nettement προσᾴδειν, vu que dans le \italique{Ravennas} l'iota est régulièrement adscrit (cf. 102, 271, 554, 986, etc.). Lisez :

\citComm{κλύεις, ὦ προσᾴδουσα τἄμ᾽ ἐν ἄντροις....}\\

C'est un bacchiaque, tout pareil au v. 1295 de l'\italique{Oreste}. Προσᾴδειν signifie « accompagner de la voix » comme προσαυλεῖν, « accompagner de la flûte ». 
Après τἄμα, sous-entendez ἄσματα ou μέλη. 
Il n'est pas d'ellipse plus fréquente après ᾄδειν : \italique{Guêpes}, 269 et 1225, ᾄδων Φρυνίχου ou Ἁρμοδίου ; \italique{Lysistr}., 1243, ἀείσω καλόν (inutile d'ajouter ἄεισμα, avec Meineke) ; Théognis, 242, καλά τε καὶ λιγέα ᾄσονται ; Platon, \italique{Timée}, 21b, τὰ Σόλωνος πολλοὶ τῶν παίδων ᾔσαμεν, etc.


\clearpage

\Note{Papyrus d'Oxyrhynque n° 4935 : vers 1043-1051 (fr. 1) ; vers 1201-1210 (fr. 2)}


\begin{figure}[!h]
\centering
\includegraphics[scale=1.83]{POxy4935.jpg}
\caption{\italique{The Oxyrhynchus Papyri}, LXXIII (2009), n° 4935. (Ed. A. Benaissa.)}
\label{POxy4935}
\end{figure}

\pagebreak

\begin{figure}[!h]
\hspace*{-1cm}
\centering
\includegraphics[scale=0.50]{POxy4935fr t.png}
\caption{\italique{The Oxyrhynchus Papyri}, LXXIII (2009), n° 4935. (Ed. A. Benaissa.) Fragment 1 : vers 1043-1051 \& Fragment 2 : vers 1201-1210.}
\label{POxy4935fr1t}
\end{figure}

\begin{figure}[!h]
\hspace*{-0cm}
\centering
\includegraphics[scale=0.35]{POxy4935fr1 c.png}
\caption{\italique{The Oxyrhynchus Papyri}, LXXIII (2009), n° 4935. (Ed. A. Benaissa.) Fragment 1 : vers 1043-1051.}
\label{POxy4935fr1c}
\end{figure}

\begin{figure}[!h]
\hspace*{-0cm}
\centering
\includegraphics[scale=0.35]{POxy4935fr2 c.png}
\caption{\italique{The Oxyrhynchus Papyri}, LXXIII (2009), n° 4935. (Ed. A. Benaissa.) Fragment 2 : vers 1201-1210.}
\label{POxy4935fr2c}
\end{figure}

%\pagebreak

\Note{Note 26 aux \italique{Thesmophories} : vers 1050-1051}

\citComm{εἶθε μὲ πυρφόρος αἰθέρος ἀστὴρ\\
	τὸν δύσμορον ἐξολέσειεν.}\\

« C'est peut-être le soleil, mais c'est plutôt la foudre que Mnésiloque appelle ἀστήρ. » Ainsi parle le scholiaste; et Kuster, en homme de sens qu'il est, d'ajouter : « Voilà pour désigner la foudre une singulière périphrase ». Singulière, en effet, au point qu'on a peine à concevoir que tant d'éditeurs, avant et après Kuster, se soient accommodés d'une telle explication. 
M. van Herwerden ne s'y est pas trompé ; seulement au lieu de se prendre à l'ignare Byzantin, auteur de la scholie, il a eu tort de s'attaquer au texte du poète, en proposant, très ingénieusement sans doute, mais très inutilement, de remplacer αἰθέρος ἀστήρ par ἀστεροπητής  « le dieu qui lance l'éclair ». 

Il n'y a rien à changer : le πυρφόρος αἰθέρος ἀστήρ ne désigne ni le soleil ni la foudre, mais tout simplement l'étoile de la Canicule, ou, pour parler grec, le Chien ou Sirius, car les deux termes sont synonymes et se prennent indifféremment l'un pour l'autre. 

Ce passage d'Aristophane touche à l'un des points les plus curieux des croyances antiques, et fournirait matière à une ample dissertation. Ce n'en est pas la place ici. Mais le sujet étant nouveau, au moins que je sache, on me permettra d'en donner un exposé sommaire.

L'histoire du mythe remonte jusqu'à Homère. L'\italique{Iliade} nous montre Achille « resplendissant comme l'astre qui s'élève durant la canicule. Ses rayons éclatants brillent dans la nuit sombre, entre la foule des étoiles. On le désigne sous le nom de chien d'Orion ». Et le poète ajoute : « C'est le plus éclatant des astres, mais il constitue un signe funeste et apporte force fièvre aux chétifs mortels\note{\italique{Iliade}, XXII, 29.} ». Ce que Virgile a traduit : 

\citComm{Ille [Sirius] sitim morbosque ferens mortalibus aegris\\
Nascitur et lævo contristat lumine cælum.}\note{\italique{Énéide}, X, 274.}\\ 

Cette étoile, Homère l'appelle ailleurs οὔλιος « la pernicieuse », et d'Homère à Nonnos, c'est-à-dire durant l'antiquité tout entière, elle gardera ce caractère de divinité malfaisante, ou plutôt elle sera le symbole même de la malfaisance\note{\italique{Iliade}, XI, 62. Oppien : οὐλόμενος κυνὸς ἀστήρ (\italique{Cyn.}, III, 325). Nonnos : λοίγιον ἀστέρα Μαίρης.}. 
Dès qu'elle se lève sur l'horizon, au plus fort de l'été, elle amène avec elle un cortège de maux. Elle stérilise les champs, brûle les plantes, prive de nourriture les germes dans les entrailles de la terre\note{\italique{Énéide}, III, 141.}, lance sur les troupeaux des misères sans nombre\note{Apollonius, \italique{Argon.}, III, 959.}. Le moindre mal qu'elle fasse aux hommes, c'est de leur dessécher la tête et les genoux. L'expression est d'Hésiode, et lui a été reprise par Alcée\note{Hésiode, \italique{Trav. et J.}, 585 ; Alcée, 45, Bergk, p. 141, Farnell [$=$ fr. 111 Reinach-Puech $=$ 347a Lobel-Page. Ces deux passages figurent dans la très belle \italique{Anthologie de la poésie grecque} de Robert Brasillach].}. Et déjà au temps d'Alcée, son ascendant est si bien établi, que seule, entre toutes les étoiles, on l'appelle Ἄστρον « l'Étoile ».

Nulle puissance n'est plus redoutable, pas même les Érinyes, qui, après tout, ne pourchassent que les coupables et dont les châtiments sont des expiations. On serait tenté plutôt de rapprocher Sirius et Até, vu l'action également funeste qu'ils exercent, celle-ci sur l'esprit des hommes, l'autre sur le corps. Encore Até ne fait-elle qu'accomplir les ordres de Zeus, dont elle est la fille. Si elle produit l'égarement et la folie, du moins est elle suivie de loin par les prières boiteuses, et l'on peut espérer la fléchir. Tandis que Sirius extermine au hasard, à la façon d'une force aveugle, et n'épargne pas plus les innocents que les criminels. 

Tous les maux soudains qui frappent la vie à sa racine ou la tarissent dans sa source lui sont attribués. Une vigne, une plante quelconque périt-elle, subitement dépouillée de son feuillage et de ses fleurs, on la dit ἀστροπλὴξ ou ἀστρόβλητος « frappée de l'astre\note{Lire là-dessus Théophraste, \italique{de Causis Plant.}, V, 9.} ». Car ce que nos médecins nomment héliose (insolation), les anciens l'appellent siriase. Un enfant au berceau est-il pris d'une de ces fièvres de l'encéphale si fréquentes à cet âge et dont la terminaison est le plus souvent fatale, le mal auquel il succombe se nomme ἀστροβλησία, en latin \italique{sideratio}\note{Ἀστροβοληθῆναι, ἐπὶ τῶν παιδίων (Gr. Étymol.). — Sans être médecin, je ne crains pas d'affirmer que l'ἀστροβλησία  ou σειρίασις τῶν παιδίων est bien la méningite. Pline, qui en parle par trois fois, traduit tantôt par \italique{adustio}, tantôt par \italique{destillatio infantium} (XXX, 47 ; XXXII, 48 ; XXII, 29). Sur quoi je ferai remarquer : \ \ \ 1° \ que le mot \italique{infantium} (παιδίων, non \italique{puerorum}) marque une affection spéciale aux enfants en bas âge ; \ \ \ 2° \  qu'\italique{adustio} signifie inflammation, tandis que \italique{deslillatio} répond exactement à ce que l'on nomme aujourd'hui exsudat ; \ \ \ 3° \ que Pline recommande comme topiques une éponge humectée souvent avec de l'eau froide, ou une grenouille appliquée par le dos sur la tête, ou bien encore les feuilles de l'héliotrope, « bonnes aussi, ajoute-t-il, dans les convulsions, même épileptiques » ; c'est-à-dire qu'on traitait le mal, comme on n'a cessé de le faire depuis, par des antiphlogistiques. Tous ces traits ne peuvent se rapporter qu'à la méningite. Les dictionnaires, tant grecs que latins, doivent être corrigés sur ce point. [Cette correction n'a toujours pas été effectuée...]}. Ce dernier mot, sous sa forme française, se lit encore dans les dictionnaires de médecine. L'homme fait se voit-il au milieu d'un état de santé florissant enlevé aux siens par une affection foudroyante, fièvre aiguë, paralysie, coup de sang, c'est encore une victime de l'astre. 

Et ne croyez pas, d'après les dictionnaires, à l'influence maligne d'un astre quelconque. C'est bien de Sirius ou du Chien qu'il est question. Hésychius le dit expressément\note{Ἀστροβολήτους· τοὺς ὑπὸ τοῦ κυνὸς βαλλομένους.}, et ce qui achève la preuve, c'est qu'à côté d'ἀστρόβλητος on trouve κυνόπληκτος, et que l'ἀστροβλησία s'appelle aussi σειρίασις. 

Contre l'astre maudit, il n'y a abri ni recours. Ceux qu'il frappe sont le plus souvent blessés à mort. Dans le chœur célèbre d'\italique{Hippolyte} où Euripide dépeint sous de si vives couleurs la puissance de l'Amour, c'est Sirius non moins que Zeus qui fournira le terme de comparaison : rien, y est-il dit, n'est supérieur au trait que lance Éros, οὔτε πυρὸς οὔτ᾽ ἄστρων βέλος « ni le trait du feu céleste, ni celui de l'étoile pernicieuse\note{\italique{Hippolyte}, 531.} » ; car c'est ainsi qu'il convient d'interpréter ce passage, qui a donné lieu à tant de méprises\note{Ἄστρων est un pluriel poétique, indiquant probablement la constellation entière. Cf. le schol. sur Eurip., \italique{Héc.}, 1104 : ὅλον τὸ ἄστρον κύων λέγεται, ὅπερ ἐκ πολλῶν ἀστέρων σύγκειται. Il se peut toutefois que le poète fasse allusion en même temps au trait de Sélènè, car nous savons par le scholiaste des \italique{Nuées}, 398, qu'on disait σεληνόπληκτος ou σεληνόβλητος pour un lunatique.}. 

Et on peut en rapprocher un autre, tiré de l'\italique{Agamemnon} d'Eschyle. Clytemnestre, qui accueille par de perfides démonstrations de respect son époux revenu de Troie, le compare à un arbre « dont le feuillage étend sur la maison son ombre protectrice de l'étoile du Chien\note{\italique{Agamemnon}, 958.} ». Suprême et adroite flatterie, donnant à entendre qu'une fois le roi rentré dans ses foyers, nul fléau, fût-il de ceux qui déjouent toute résistance, n'est plus à redouter. 

Après cela, faut-il s'étonner que les Grecs aient rendu à Sirius des honneurs divins ? Sans entrer à ce propos dans toutes les discussions auxquelles la mythologie astronomique a donné lieu, je me bornerai à constater que le Chien a pris place de bonne heure parmi les divinités, et que son culte se perd dans la nuit des temps. Car c'est le fils bienfaisant d'Apollon, celui qui présidait à la fécondité des troupeaux, à l'élève des abeilles, aux travaux des champs, c'est Aristée qui l'institua. La légende est racontée au long par Apollonius de Rhodes. Les habitants des Cyclades étaient livrés sans défense aux mortelles ardeurs de l'astre. Aristée, sur l'ordre de son père, écarta d'eux le fléau : il éleva dans l'ile de Céos un grand autel à Zeus Ikmaios (qui répand la pluie) et célébra sur les montagnes des sacrifices en l'honneur de Sirius à la fois et de Zeus : « Et maintenant encore à Céos, des prêtres immolent des victimes un peu avant le lever du Chien\note{\italique{Argonautiques}, II, v. 526. — Une allusion à ce culte se lit aussi dans un fragment des \italique{Aitia} de Callimaque, retrouvé en Égypte, sur le mythe d'Acantios et Cydippé, vv 34 et suiv.} ». Ce fait si curieux d'un culte permanent rendu à Sirius nous est confirmé par des médailles\note{\italique{Catal. of the Greek coins of the British Museum. Creta and the Ægean islands}, by W. Wroth, 1886. L'île de Céos comprenait quatre villes. Le Chien entouré de rayons se trouve sur les monnaies de Céos et de Carthæa. L'Astre, sur celles de Coresia, Iulis et Poiessa.}, ainsi que par Théophraste, Diodore et Cicéron\note{Théophraste, \italique{de Ventis}, 14 ; Diodore, IV, 82 ; Cicéron, \italique{de Divin}., I, 57.}. 

C'est donc bien l'astre promu à la dignité de dieu qu'invoque Mnésiloque. Il ne nous reste plus qu'à déterminer le sens de πυρφόρος. On aurait tort de tenir ce mot pour un équivalent de φλογοείς « embrasé ». Sans doute, rien ne s'oppose à ce qu'on dise, et on a dit en effet, φλογοεὶς κύων, par rapport aux ardeurs de la Canicule\note{Straton, dans l'\italique{Anthol. Pal}., XII, 225.}. Mais πυρφόρος  a un sens bien plus précis : il veut dire « armé d'un feu dévorant, d'un feu meurtrier » par allusion à la fièvre (πυρετός, aussi nommée πῦρ). C'est aussi l'épithète dont use Sophocle, au début de l'\italique{Œdipe Roi}, pour caractériser le fléau qui désole Thèbes, ὁ πυρφόρος θεὸς. Λοιμὸς ἔχθιστος (v. 27). On peut y voir, si l'on veut, une réminiscence du vers de l'\italique{Iliade} que nous citions en commençant. 

Si, maintenant pour finir, nous voulions traduire en langage de comédie l'imprécation de Mnésiloque, il ne faudrait pas la rendre par : Que la foudre m'écrase ! Ce souhait de malheur est le premier peut-être qui vienne à l'esprit ; c'est pourquoi le scholiaste s'y est arrêté, mais à contre-sens. Molière nous offre le choix entre deux autres formules : Que la fièvre me serre ! Que la peste m'étouffe\note{Dans une nouvelle d'Ém. Souvestre, La male mort me vienne ! \italique{Revue des Deux Mondes}, 1853, t. IV, p. 965.} ! C'est bien ainsi qu'il faut interpréter. Car si la foudre n'obéit qu'à Zeus, la peste et la fièvre sont aux ordres de l'Astre.



\Note{Note 27 aux \italique{Thesmophories} : vers 1052-1055}


\citComm{οὐ γὰρ ἔτ᾽ ἀθανάταν φλόγα λεύσσειν\\
	ἐστὶν ἐμοὶ φῦλον, ὡς ἐκρεμάσθην\\
	λαιμότμητ᾽ ἄχη δαιμόνων αἰόλαν\\
	νέκυσιν ἔπι πορείαν.}\\

C'est abuser de la syntaxe que de prétendre avec Fritzsche que les mots ἄχη et πορείαν sont des accusatifs d'apposition à la phrase entière. Pareils accusatifs ne sont en somme que des tournures elliptiques toujours faciles à ramener à une construction régulière, comme dans ce vers d'Euripide : Ἑλένην κτάνωμεν, Μενέλεῳ λύπην πικκράν. Ici l'on ne saurait exposer clairement en vertu de quelle ellipse les mots en question se rattacheraient à la proposition principale. 

D'autre part δαιμόνων est dénué de sens ; et c'est évidemment là, ainsi que l'a vu Bentley, qu'est le nœud à résoudre. Je n'hésiterais pas à lire :

\citComm{\decalage{ἐστὶν ἐμοὶ φῦλον, }ὡς ἐχρεμάσθην,\\
	λαιμότμητ᾽ ἄχη δία μένων αἰόλαν....}\\

De la sorte tout devient clair. Construisez : διὰ λαιμότμητα ἄχη μένων αἰόλαν πορείαν ἐπὶ νέκυσιν. -- Διά avec l'accusatif signifie \italique{propter} ; μένω, c'est \italique{manere}, \italique{expectare} ; l'adjectif αἰόλος, interprété de si étrange façon, est pris dans son acception première de ταχύς.

Mnésiloque n'a plus goût à la vie ; l'excès de ses souffrances fait qu'il n'aspire qu'à un prompt trépas. C'est ce qu'il exprime à sa manière : « Car me voilà suspendu, n'attendant en raison d'un supplice qui me jugule\note{Λαιμότμητα fait allusion au carcan qui lui serre le cou (v. 1104). N'allons pas traduire avec M. van Leeuwen : un supplice à se couper la gorge. Le savant professeur suppose que Mnésiloque a toujours en main le coutelas qu'il a pris sur l'autel et qui lui a servi à graver son message. Ce serait, je l'avoue un plaisant spectacle, mais qui donnerait une singulière idée de la vigilance du Scythe qui a attaché le patient au poteau.} qu'un prompt trajet vers les morts ». 

Après cela, je ne fais pas difficulté de reconnaître que l'anastrophe est moins fréquente avec διά qu'avec d'autres prépositions. Mettons, bien qu'on la trouve encore dans Hermésianax, qu'au temps d'Aristophane elle pouvait passer pour un archaïsme\note{Hésiode, \italique{Tr. et J.}, 3 : ὅντε δία ; Tyrtée (Bergk, 5) : ὃν δία ; Eschyle, \italique{Sept.}, 206 : πηδαλίων δία ; \italique{Agam.}, 590 et 1133. Cf. Kühner-Blass, §86, 1}. Ce serait une raison de plus de l'admettre dans une parodie bouffonne où le poète, pour se moquer d'Euripide, s'est complu à prodiguer les termes affectés et les singularités de langage. Et c'est aussi pourquoi il faut préférer le datif νέκυσιν ἔπι, pris au sens prégnant, à l'accusatif νέκυς (Kühner-Gerth, §447, A, 2).

\Note{Note 28 aux \italique{Thesmophories} : vers 1181}

\citComm{φέρε θοἰμάτιον ἄνωθεν, ὦ τέκνον, τοδί.}\\

Ce vers est d'une clarté telle qu'il est à peine concevable qu'on ait pu s'y méprendre. Les éditeurs y ont substitué : φέρε νυν, κατάθου θοἰμάτιον...  ou φέρε, θοἰμάτιον ἄνες μὲν... « Or çà, dépose ton manteau. » 
Ils oublient qu'il s'agit d'une danseuse, et que, pas plus que les joueuses de flûte, les danseuses ne s'embarrassaient de manteaux. Ἱμάτιον, chacun le sait, se dit aussi d'un vêtement quelconque ; ici il désigne le chiton : preuve, le démonstratif τοδί, qui serait superflu s'il s'agissait du manteau. 

Φέρε signifie \italique{trade} ou \italique{porrige mihi}, comme aux vv. 1115, 1196, \italique{Caval.}, 118, \italique{Nuées}, 1297, \italique{Paix}, 15, et vingt fois ailleurs ; ἄνωθεν veut dire \italique{par en haut}, c'est-à-dire en le passant par-dessus la tête, ce qui est, en effet, la manière laplus usitée d'ôter une tunique ou une chemise. Euripide, qui a son plan, dit à la danseuse : « Tends-moi ton vêtement, ma fille, en le tirant par en haut. Puis assieds-toi sur les genoux du Scythe, et avance les pieds que je les déchausse ». 

Lisez ce qui précède, vous verrez que quand elle reçoit cet ordre il y a déjà tout un temps que la jeune femme dansait (sans manteau, est-il besoin de le dire ?). Lisez après cela ce qui suit, vous verrez clairement qu'une fois qu'on lui aura tiré son vêtement et ses chaussures, elle demeurera nue comme la main. Et c'est bien sur ce moyen qu'Euripide compte pour mener le Scythe à ses fins. 

Il ne restera plus l'ombre d'un doute, je suppose, quand on aura lu le texte suivant de Platon, qui pourrait servir de scholie à notre passage : Ἐκτείνας τὴν χεῖρα καὶ λαβόμενος τοῦ ἱματίου ἄνωθεν αὐτοῦ παρὰ τὸν ὦμον\note{\italique{Républ}., V, p. 449b.}. « Polémarque avança la main et saisit le manteau par le haut de l'endroit de l'épaule. »

%\pagebreak

\Note{Papyrus d'Oxyrhynque n° 3840 : vers 1185-1193}


\begin{figure}[!h]
\centering
\includegraphics[scale=0.33]{POxy3840 t.png}
\caption{\italique{The Oxyrhynchus Papyri}, LVI (1989), n°3840. (Ed. Peter J. Parsons.)}
\label{POxy3840}
\end{figure}
\Note{Le nu dans la Comédie ancienne des Grecs}\note{Extrait des \italique{Bull. de l'Acad. roy. de Belgique} (Classe des beaux-arts), n° 9-10, (sept.-oct.), 1901. L'article était suivi d'un appendice sur la répartition des places au théâtre (reproduit au tiers tome de l'\italique{Aristophane} posthume de Willems).}
La question du nu dans l'art prête à des controverses sans fin, dans lesquelles il ne me convient pas d'entrer. Je ne discute point, je constate. On s'est beaucoup occupé dans ces derniers temps de ce qu'on appelle l'âme des peuples. Peut-être trouvera-t-on que la présente étude jette une lumière inattendue sur l'âme grecque. Des spectacles que la police proscrirait aujourd'hui, non seulement l'usage les autorisait, mais je ne sache pas qu'aucun moraliste ancien ait songé à les censurer. Au contraire, Socrate me sera garant qu'en introduisant le nu sur leur scène comme dans leurs arts plastiques, les Grecs, loin de viser à une excitation des sens, ne faisaient que poursuivre leur éternel idéal de la beauté.\\

Quand on examine sans parti pris la question du nu dans la comédie antique, on ne laisse pas d'être surpris des conclusions où l'on aboutit, car elles sont en complète contradiction avec toutes les idées reçues.

Le problème se présente sous une double face, suivant qu'il a trait à l'un ou l'autre sexe. À première vue, il semblerait que chez un peuple habitué aux exercices de la palestre, le nu masculin ne dût choquer personne ; qu'en fait de femmes le nu au contraire n'eût pas même l'occasion de se produire au théâtre, par la raison que les rôles féminins, nous dit-on, étaient exclusivement réservés aux hommes. Cette manière de raisonner paraît logique, malheureusement les faits la démentent. Pour rester dans le vrai, il faut prendre tout le contre-pied. 

D'abord, quant au nu masculin, on peut affirmer catégoriquement que les Grecs l'ignoraient en matière scénique. Acteurs et choreutes portaient une sorte de maillot collant, au besoin rembourré, qui couvrait le corps entier et qu'on appelait σωμάτιον. Comment en eût-il été autrement ? C'est précisément parce que des acteurs étaient chargés des rôles de femmes que le maillot s'imposait. Il est inutile, je suppose, d'entrer dans des explications. Or, à part les \italique{Cavaliers}, il n'est pas une comédie d'où l'élément féminin soit exclu. Souvent même une femme tient le premier rôle, et le chœur dans l'orchestre se compose uniquement de femmes. Il arrive aussi que le chœur soit divisé en deux parties, l'une de femmes, l'autre d'hommes. Dans tous ces cas, on reconnaîtra que le maillot est indispensable, aussi pour ces derniers. Car on ne se figure ni une troupe d'acteurs ni un corps de ballet, jouant, dansant et s'entre mêlant, moitié avec, moitié sans maillot. La disparate eût été trop sensible. 

Grâce à cet uniforme de dessous, hommes et femmes pouvaient sans inconvénient déposer le manteau et rester en simple tunique : tels les choreutes au début des parabases\note{\italique{Acharniens}, v.627.}, ou quand ils se disposent à danser\note{Femmes en simple tunique, \italique{Lysistr.}, 686 ; \italique{Thesmoph.}, 656.}. De la tunique même, on pouvait au besoin se passer. Ainsi dans les \italique{Oiseaux}\note{\italique{Oiseaux}, vv. 993-947.}, au milieu d'une cérémonie religieuse, l'acolyte du prêtre est invité à se dépouiller de ses vêtements au profit d'un poète dépenaillé, et il ôte successivement sa pelisse et sa tunique, sans qu'il en résulte le moindre effet comique, le σωμάτιον parant à tout. 

C'est donc, si l'on veut, par un sentiment de décence (nous ne disons pas de pudeur), mais c'est surtout par des considérations d'ordre esthétique que s'explique l'emploi du maillot. Mais la gaieté n'y perdait rien. L'attribut viril que le costume dissimulait, il ne tenait qu'à l'auteur comique de le rétablir artificiellement au moyen du phallus. Ne perdons pas de vue que le théâtre grec procédait en droite ligne des fêtes du dieu Dionysos, et rappelons la définition si nette d'Aristote : la tragédie tire son origine du dithyrambe ; la comédie, des chants phalliques\note{Poétique, chap. IV.}. Ces chants, on s'en fait à peine une idée d'après celui qui se lit dans les \italique{Acharniens}\note{\italique{Acharniens}, vv. 263-279.}.
Ils devaient être autrement licencieux et folâtres quand ils accompagnaient le comos bachique qui marquait la fin des vendanges. D'après leur nom, on juge de la place qu'y occupait le phallus, symbole de la force productrice et de la fécondité de la nature. Pour mal renseignés que nous soyons sur la comédie primitive, telle qu'elle subsistait à l'époque de Périclès, il est sans doute que le symbole en question y avait un rôle prépondérant. 

Aristophane s'est-il conformé sur ce point à la tradition ? Le problème est des plus intéressants, car il touche au caractère même de son art. On l'a débattu souvent et tranché en sens divers. Nous croyons qu'en pesant attentivement et en interprétant comme il convient les données fournies par son théâtre, il n'est pas impossible d'arriver à une solution définitive.

À cette question : Aristophane a-t-il, comme ses devanciers, usé régulièrement du phallus ? nous n'hésitons pas à répondre non. S'il en a usé, ce n'est que dans des cas très exceptionnels et pour ainsi dire à son corps défendant. Nous avons épluché scène par scène ce qui reste de son œuvre, et voici ce que nous y avons relevé : dans les \italique{Acharniens}, le corps des Odomantes ; dans les \italique{Thesmophories}, Mnésiloque et l'archer scythe. La première scène des \italique{Acharniens} a lieu dans le Pnyx. Théoros, de retour d'une ambassade auprès du roi thrace Sitalcès, promet monts et merveilles aux Athéniens, et à l'appui de ses dires produit un prétendu corps auxiliaire d'Odomantes. Ceux-ci s'avancent accoutrés d'une manière burlesque, à demi nus et portant le phallus. Or c'est ce dernier détail qui sert à déceler la fraude, ces soi-disant barbares n'étant que des débauchés pris dans la lie de la population athénienne. L'épisode est fort court et n'a pour but que de convaincre les gobe-mouches de l'ecclésie de se laisser prendre à d'immondes supercheries. 

Quant aux \italique{Thesmophories}, il suffit de rappeler le sujet en deux mots. Les femmes d'Athènes ont pris occasion de la fête qui les réunit pour conspirer contre Euripide. Afin de détourner le coup, son beau-père Mnésiloque s'est déguisé en femme. Mais il est reconnu et condamné lui-même à périr. On voit que le point de constater le vrai sexe du coupable est le principal ressort du drame. De là, le phallus. Le dénouement est à l'unisson avec le reste. Après avoir essayé de divers moyens pour délivrer son parent, Euripide se décide à recourir à de grossières séductions et à prendre par les sens l'agent de police qui a la garde du prisonnier. Le moyen, sur lequel nous aurons à revenir, lui succède, et c'est par là que finit la pièce. Les \italique{Thesmophories} n'ont rien de commun avec la politique ; le poète n'a songé qu'à amuser le public en ridiculisant Euripide sous le masque de son beau-père. C'est la seule de ses comédies qui soit exclusivement bouffonne.

Les Odomantes, Mnésiloque et l'archer scythe, voilà à peu près tout ce que nous trouvons à citer. Ajoutez Dicéopolis, mais seulement tout à la fin des \italique{Acharniens}, et Philocléon, dans une scène unique des \italique{Guêpes}. Mais dans les deux cas le détail est simplement emprunté à la réalité. L'un et l'autre personnage montent sur la scène ivres et débraillés, en compagnie de courtisanes. Le premier a remporté le prix de boire à la fête des Choés, l'autre s'en revient d'une orgie. Les peintures de vases font foi qu'en pareille circonstance les Athéniens n'observaient pas mieux qu'eux la décence, et n'étaient guère plus gênés de donner aux passants le spectacle de leur nudité. 

D'autres témoignages qu'on serait tenté de citer ne prouvent rien ou prouvent le contraire. Quand Strepsiade couché sur son grabat se livre à ses méditations, et qu'à cette demande de Socrate : « Tiens-tu quelque chose ? », il répond : « Rien que mon cas de la main droite »\note{\italique{Nuées}, v.734.} ; quand Trygée, sur le point de monter au ciel sur son escarbot, rassure ses filles en leur disant : « J'ai pour me diriger un gouvernail dont j'userai »\note{\italique{Paix}, v.142.}, ce sont là de simples saillies d'esprit. S'ils avaient eu l'accessoire qu'on suppose, ils auraient employé le démonstratif τόδε ou τοδί, et à coup sûr il y serait fait allusion plus d'une fois dans le courant de la pièce. Autant peut-on dire de la boutade contre Lamachos, dans les \italique{Acharniens} : εὔοπλος γὰρ εἶ (\italique{tu es bien outillé}, mais aussi \italique{tu es bien armé}) : pur jeu de mots suggéré par la lance, le casque à triple panache et le bouclier à tête de Gorgone, dont Lamachos, en fanfaron qu'il est ou qu'on veut le faire paraître, s'est ridiculement affublé\note{\italique{Acharniens}, v.592.}. Voilà donc, sauf une exception dont nous parlerons tout à l'heure, voilà, tout compte fait, cinq fois que l'acteur se montre muni du phallus ; encore Mnésiloque est-il seul à le porter d'un bout à l'autre de la pièce.

Quoi d'étonnant ? La nature même des sujets traités n'excluait elle pas toute obscénité dans le costume ? L'acteur aurait eu un joli succès qui, armé d'un priape, serait venu flétrir les excès de la démagogie, se poser en réformateur des abus, en avocat de la paix, en champion de l'austère éducation d'autrefois, en défenseur des hautes traditions de l'art. Autre est la farce, autre la comédie, surtout la comédie politique. Celle-ci peut tenir le langage le plus libre et le plus cru, ne reculer ni devant le détail scabreux ni devant le mot propre, faire arme de tout pour bafouer l'adversaire, le fond n'en est pas moins sérieux, parfois triste, et elle n'aura garde de provoquer le rire aux dépens des causes qu'elle défend. Joignez à cela qu'Aristophane est par excellence un esprit aristocratique. Il se vante hautement d'avoir dignifié et ennobli sa profession\note{\italique{Paix}, vv. 739-753.}, et nous devons l'en croire. Jamais il ne manque l'occasion de protester contre les inventions triviales, contre « la scurrilité et les basses bouffonneries », et précisément, parlant des \italique{Nuées}, il déclare dans la parabase : « Ma comédie est venue sans avoir cousu sur elle un phallus de cuir, pendant, épais, rouge par le bout, pour donner à rire aux gamins\note{\italique{Nuées}, v. 538.} ». 

Comment faut-il prendre cette déclaration ? Faut-il y voir l'expression d'un système arrêté, et se donner, à l'exemple de nombreux interprètes, le facile plaisir de mettre le poète en contradiction avec lui-même ? Pas le moins du monde. Tout obscène qu'il était, l'emblème générateur ne choquait pas les Grecs au même point que nous. Il faisait partie intégrante du costume des choreutes dans le drame satyrique, trouvait son emploi dans certains rites religieux et, dans toutes les rues d'Athènes, il s'étalait impudemment sur la gaine des Hermès. Il est à croire que l'objet en soi laissait le poète assez indifférent. Ce qu'il n'admettait point, c'est qu'on l'employât sans motif, qu'on en fît un rouage essentiel tenant lieu d'esprit et d'intrigue, et ne servant qu'à déguiser la pauvreté d'invention. Il en va de cela comme des licences de la cordace, de certains mots orduriers, des plaisanteries rebattues, du pont aux ânes de la voracité d'Héraclès\note{Comparez à cet égard : \italique{Nuées}, 555 (sur la cordace), avec la cordace dansée à la fin des \italique{Guêpes} ; \italique{Grenouilles}, 8 (sur les plaisanteries ordurières), avec \italique{Nuées}, 295 ; \italique{Guêpes}, 60, et \italique{Paix}, 741 (sur Héraclès), avec \italique{Oiseaux}, 1585 et suiv., et \italique{Grenouilles}, 549 et suiv.}, autant de moyens dont il se raille, mais dont il se réserve d'user en temps et lieu, quand la situation le comportera ; et sur ce point nous pouvons nous fier à lui, car jamais auteur dramatique ne sut mieux son jeu. 

En somme, hormis les cas cités, il n'a guère eu recours à cet expédient qu'il jugeait vulgaire. Même dans les épisodes où son caprice s'amuse et se joue, ses personnages portent le costume de tous les jours. Telle scène où l'emploi de l'outil de cuir eût paru justifié est conduite de manière à ne pas laisser de place au doute. Prenons, pour nous en tenir à un exemple, celle de l'\italique{Ecclésie des Femmes} où un jeune homme est cruellement tiraillé par trois vieilles en amour ; il faudrait n'avoir aucune expérience scénique pour ne pas voir que les choses se seraient passées différemment et que le dialogue aurait pris un tout autre tour, si l'insuffisance du costume y eût donné prétexte. 

Il en est de même partout. Une seule pièce fait exception, et le cas a paru si rare que le scholiaste le fait remarquer\note{Schol. des \italique{Nuées}, v. 542 : τοὺς μὲν γὰρ φάλητας εἰσήγαγεν ἐν τῇ Λυσιστράτῃ.}. Chacun connaît le sujet de \italique{Lysistrata}. Les femmes grecques, pour forcer leurs maris à conclure la paix, ont juré de leur refuser le devoir conjugal, puis elles s'emparent de l'acropole. La conspiration réussit, et les résultats d'une continence forcée se trahissent, tant chez les Athéniens que chez les Spartiates, de façon la plus visible. Il s'ensuit que la pièce dégénère un moment en une véritable phallophorie. Ah ! n'accusez pas le poète, car jamais il n'a fait plus noble besogne. Sa patrie se trouvait dans la plus grave des situations. Après le désastre de Sicile, après la défection des alliés, la paix restait pour Athènes l'unique chance de salut. Pour y amener les esprits, il ne s'agissait de rien moins que de braver le préjugé populaire, plaider dans un certain sens la cause de l'ennemi en montrant que les torts n'étaient pas tous de son côté. 
Si cette fois il a recouru à des moyens qu'il réprouvait, c'est, nous le montrerons ailleurs, qu'il n'avait pas d'autre choix. Le rire, j'entends le rire immodéré, le rire fou, il fallait le provoquer à tout prix. C'était l'unique ressource dont il disposât pour atteindre son but, et ce but était sacré. 

Nous passons maintenant à l'autre face du problème du nu, celle qui concerne les femmes. Car on va voir que la règle interdisant à celles-ci l'accès de la scène souffrait d'assez nombreuses exceptions. Si nous avons pu affirmer que jamais l'homme au théâtre ne se montre complètement nu, il n'en va pas de même pour la femme. La question soulevée ici est neuve. Les historiens littéraires, qui ont relevé jusqu'aux moindres particularités du théâtre attique, n'ont pas touché à celle-là, de peur évidemment d'être taxés de paradoxe. Rien en effet ne s'éloigne davantage de nos idées et de nos mœurs. Et pourtant, pour qui sait lire autrement que des yeux, le doute n'est pas possible, et l'on peut même dire que l'ancienne comédie a marché dans cette voie aussi loin qu'il y a moyen d'aller. 

L'étonnement cessera, si l'on veut bien se mettre au diapason des anciens et entrer dans leur manière de voir et de sentir. Les Grecs, ne nous lassons pas de le répéter, n'avaient pas la même mesure de morale et de pudeur que nous. Le nu, pourvu qu'il eût le caractère de la beauté, éveillait chez eux un sentiment de pure admiration. Loin de le blâmer, les meilleurs en recherchaient le spectacle. Qui en douterait n'a qu'à lire le \italique{Banquet} de Xénophon. Le maître de la maison, pour charmer ses convives, a imaginé de faire venir un Syracusain, lequel entre, suivi de deux beaux adolescents : un jeune garçon qui danse et joue de la cithare, et une joueuse de flûte qui est en même temps une acrobate accomplie. 
Quand ils ont suffisamment fait montre de leurs talents, un des conviés, séduit par la beauté du jeune couple, s'adresse au Syracusain, et après avoir exprimé sa répugnance pour les jeux dangereux : « S'ils dansaient dans le costume sous lequel on nous dépeint les Charites, les Heures et les Nymphes, ils s'en tireraient plus aisément et le banquet nous paraîtrait beaucoup plus agréable\note{Xénophon, \italique{Banquet}, VIII, 5.} ». 
Et le Syracusain de déférer avec empressement à ce désir. Or savez-vous bien qui parle ainsi ? C'est Socrate, l'homme que n'effleura jamais la pensée même du mal, et le trait nous est rapporté dans un écrit fait à sa louange par le plus zélé et le plus austère de ses disciples. Qui osera prétendre après cela que les Athéniens, au théâtre, outrageaient de parti pris la pudeur, ou leur faire un grief de ce que le plus sage d'entre eux jugeait irrépréhensible ? 

Quoi qu'il en soit, Aristophane non seulement n'évite pas, mais il cherche manifestement l'occasion de pareilles exhibitions. 
La moitié de ses pièces nous donnent le spectacle d'une femme nue. Les \italique{Guêpes} d'abord, où le vieux Philocléon, à peu près nu lui-même, s'avance tenant par la main une joueuse de flûte entièrement nue\note{\italique{Guêpes}, vv. 1342 et suiv.}. Et si je cite cette scène avant les autres, c'est qu'elle nous dispose de recourir aux inductions. Le dialogue est suffisamment clair. Certains détails, trop vifs pour être cités, mettraient en défaut quiconque essaierait de soutenir que la courtisane enlevée par Philocléon a conservé le moindre vêtement, fût-ce le σωμάτιον.

De même Théoria dans la \italique{Paix}. L'auteur aurait pu, s'il l'avait voulu, faire jouer à ce personnage allégorique un des rôles principaux. Il a préféré le réserver pour la scène où Trygée en fait remise à la Boulê. Théoria est arrivée brillamment parée de tous ses attributs. Trygée lui ordonne de quitter ses vêtements, sans en excepter aucun : lui-même, quelques vers plus loin, a soin de le souligner\note{\italique{Paix}, vv. 886 et suiv}. Après quoi il descend dans l'orchestre, et, on en pensera ce qu'on voudra, mais le fait est indéniable, il traverse toute la place\note{Voir l'appendice.} et vient se poser, avec la jeune femme dans l'état de pure nature, devant la section occupée par le Conseil des Cinq-Cents : « Allons, prytanes, recevez Théoria ! » 
Les prytanes sont les bouleutes temporairement de service, qui à ce titre occupent les places de devant. Et cette fine mouche qu'est Aristophane a deviné d'avance la petite scène qui va se jouer entre les magistrats ainsi mis en demeure et la figurante : « Voyez l'accueil empressé que lui fait ce prytane. Quelle différence s'il avait fallu introduire une affaire gratis ! Tu n'aurais pas manqué de te retrancher sur les vacances ». 
Attrapez, prytanes! Vous faites payer cher votre concours ; tant pis pour celui qui refuse, il aura à attendre par fois une année entière avant d'être autorisé par vous à présenter sa requête\note{Xénophon, \italique{Républ. Athén.}, 3, 1.}. Le poète, qui n'oublie jamais son rôle de justicier, a trouvé cette occasion de vous le dire au nez, et voilà ce qui s'appelle faire d'une pierre deux coups. 

La scène de l'archer scythe, dans les \italique{Thesmophories}, n'est pas moins significative. Elle forme le pendant de celle de la \italique{Paix}. Seulement Trygée est remplacé par Euripide déguisé en vieille entremetteuse, et Théoria par une danseuse. La vieille fait passer à celle-ci son vêtement par-dessus la tête, puis l'ayant fait asseoir sur les genoux du Scythe, elle lui ôte jusqu'à ses chaussures\note{\italique{Thesmophories}, vv. 1181-1183.}. Ici mon analyse s'arrête : le reste est à lire dans l'original. À la fin des \italique{Acharniens}, Dicéopolis, vainqueur aux Choés, revient à demi ivre entre deux courtisanes. Cette fois, il est vrai, le contexte n'est pas aussi explicite. Mais si l'on raisonne par analogie et qu'on fasse compte des propos tenus par l'acteur principal, on se convaincra que les courtisanes ne sont guère vêtues, pas plus d'ailleurs qu'il ne l'est lui-même\note{\italique{Acharniens}, vv. 1198 et suiv.}. 
Il est aisé de se faire une idée du groupe par la vignette insérée dans le Dictionnaire de Daremberg et Saglio, laquelle représente au naturel un comos revenant du banquet des Choés\note{Art. \italique{Dionysia}, t. III, p. 237.}. 
J'allais oublier Διαλλαγή, la \italique{Conciliation} et le rôle qui lui revient dans la comédie de \italique{Lysistrata}. Déjà dans les \italique{Acharniens}, Aristophane l'avait invoquée, et nullement comme une abstraction, car il lui prête un beau visage, et fait d'elle la compagne d'enfance de Cypris et des Charites. 
Cette fois il la fait apparaître en personne à l'appel de Lysistrata. D'où vient-elle? Il n'y a pas de doute, de même que Théoria et Opora dans la \italique{Paix}, elle descend directement du ciel, où la fureur de la guerre l'a tenue reléguée. 
Sous quels traits faut-il se la représenter ? Il n'y a pas davantage à hésiter. Elle n'a point, comme ses émules, des attributs qui la fassent reconnaître. A quoi bon d'ailleurs ? Son unique rôle est de se faire aimer et désirer. Aussi se montre-t-elle, comme Cypris et les Charites, « ses compagnes inséparables », sans voile aucun, et n'ayant pour tout moyen de séduction que sa beauté. 

Lysistrata a su ce qu'elle faisait en l'appelant auprès d'elle. En femme avisée qu'elle est, elle s'est rendu compte qu'à elle seule l'éloquence la plus persuasive n'aurait pas de prise sur des rivaux uniquement butés à se nuire réciproquement. Si de beaux discours suffisaient, il y a bien longtemps que la paix serait faite. Voilà pourquoi elle compte surtout sur la présence réelle, sur les attraits visibles et palpables de la Conciliation. Et les choses se passent comme elle l'avait prévu. Ses auditeurs, pendant qu'elle pérore, ne lui prêtent qu'une attention distraite. Tous les yeux sont fixés sur les appas de la jeune femme, et chacun, suivant sa complexion propre, traduit son ardeur en des termes qu'autorise la liberté du théâtre attique, mais que le plus osé des interprètes se refuserait à traduire textuellement.
Le résultat ne se fait pas attendre. Mordus par le désir, les délégués de Sparte et d'Athènes souscrivent à toutes les conditions qu'on leur propose et s'en vont joyeusement célébrer la paix par un banquet. 

Enfin, dans les \italique{Grenouilles}, quand apparaît la muse d'Euripide sous les traits d'une joueuse de castagnettes, il n'y a pas de doute qu'elle soit nue. La réflexion de Dionysos le donne clairement à entendre et les peintures de vases le confirment\note{\italique{Grenouilles}, v. 1308. — Cf. A.-S. Murray, \italique{Designs from greek vases in the British Museum}. Londres, 1894, pl. VI, n° 23.}. 

On pourra s'étonner que je n'aie pas cité dès l'abord la fameuse scène de Myrrhine et Cinésias dans \italique{Lysistrata}. Elle passe en effet pour la plus libre du théâtre d'Aristophane, et je n'entends pas y contredire. Mais si Myrrhine y perd toute retenue, ce n'est qu'en paroles ; en somme elle reste, sinon drapée, du moins couverte jusqu'à la fin. A. de Musset a dit quelque part : « A voir cette rusée commère, les spectateurs eux-mêmes devaient partager le tourment de Cinésias, pour peu que la scène fût bien rendue\note{\italique{Lettres de Dupuis et Cotonet}, première lettre.} ». Le mot est joli, mais il porte à faux. Il ne tient pas compte d'une circonstance essentielle, à savoir que Myrrhine, les spectateurs ne l'ignoraient pas, était jouée par un homme. 

Ici nous touchons à un des points les plus curieux, et j'ose dire les moins compris, de la pratique théâtrale chez les Grecs. Les personnages de femmes étaient remplis par des hommes, c'est entendu. Mais qu'on ne se figure pas que le respect du sexe et la morale y fussent pour quelque chose\note{Méprise où est tombé après bien d'autres un écrivain que je me plais à citer, parce qu'il joint le savoir au goût : « Qu'eussent dit les Athéniens de Périclès s'ils avaient vu une femme se montrer sur la scène ? Il est indécent qu'une femme paraisse en public. » Anatole France, \italique{Thais}, fin de la première partie.}. 
Cette exclusion s'explique suffisamment par des raisons toutes matérielles et physiques. Il n'était pas de profession plus fatigante que celle de l'acteur, ni qui impliquât des aptitudes plus diverses. 
Suivez à ce point de vue le Philocléon des \italique{Guêpes}. Vous verrez qu'aux dons, naturels ou acquis, nécessaires à tout homme qui se produit devant un public, la beauté des formes, la distinction des gestes et des attitudes, la sûreté de mémoire (il n'y avait pas de souffleur), il fallait en ajouter d'autres plus spéciaux. 
Le comédien devait être doublé d'un chanteur et d'un danseur, et compliqué aussi d'un gymnaste rompu à tous les exercices du corps. 
Et ce n'est pas tout. Il devait en outre posséder un organe d'une sonorité dont on ne peut plus se faire une idée. Souvenons-nous que les théâtres étaient à ciel ouvert et faits pour contenir toute la population, sans compter les étrangers : à Athènes, 17,000 spectateurs\note{C'est l'évaluation de M. Dörpfeld. Platon exagère évidemment quand il parle de plus de 30,000 spectateurs (\italique{Banquet}, chap. III), non moins qu'Aristophane avec ses « innombrables myriades » (\italique{Guêpes}, v. 1010)}.
La voix devait porter sans difficulté jusqu'aux derniers rangs des gradins. Les anciens nous ont transmis une foule de termes désignant cette puissance vocale : βομβεῖν, ληκυθίζειν, λαρυγγίζειν, φαρυγγίζειν. 
Aussi les acteurs exerçaient leur voix comme les athlètes leurs membres, et l'on ne réussissait à se faire entendre qu'au prix d'un long entraînement. 

Prétendre, comme on l'a fait encore récemment, qu'Aristophane, faute d'un acteur qui osât s'en charger, joua lui-même le rôle du protagoniste dans les \italique{Cavaliers}, provoque le sourire. On soutiendrait plus plausiblement que Meyerbeer s'offrait au besoin à chanter la partie de Raoul de Nangis dans les \italique{Huguenots}. 
L'auteur n'était que le διδάσκαλος, c'est-à-dire l'instructeur ou le répétiteur ; encore se fit-il de bonne heure suppléer par un spécialiste dans cette ingrate fonction. Comment les choses se passaient à l'époque de Thespis, nul ne le sait. Il se peut qu'à l'origine l'auteur jouât lui-même\note{Aristote l'affirme positivement, \italique{Rhetor.}, III, 1, 3}. Mais cela devint impossible quand le théâtre eut atteint son plein développement. 
C'est là une question, non de textes, mais de bon sens. On citera Cratès, mais le cas est différent : Cratès est un acteur qui, comme Molière, de comédien devint auteur comique. Il en est de même probablement de Carkinos et de son fils Xénoclès, de Phérécratès, de Philémon\note{Meineke, \italique{Hist. crit. com. gr.}, t. I, pp. 59, 66 et 369.}, peut-être aussi de Mélanthios. 
On nous citera encore Sophocle émerveillant les spectateurs dans \italique{Nausicaa} par son habileté à jouer de la balle. Qu'y a-t-il d'extraordinaire ? Sophocle passait pour le plus beau des Athéniens. N'est-il pas tout simple qu'il ait séduit par sa grâce dans un rôle muet et purement épisodique ? 

Nous ne sommes pas si loin qu'il semblerait de notre sujet. Sophocle nous y ramène naturellement et va nous fournir notre conclusion. Ce que nous disions de lui peut se dire également des femmes. Si des motifs d'ordre physique fermaient à celles-ci la carrière dramatique, ces motifs n'existaient plus dès qu'il s'agissait de personnages muets. 

L'artiste qui joue Procné dans les \italique{Oiseaux} est bien une femme, une joueuse de flûte qui devait être d'une rare beauté. 
Le dirai je ? J'ai toujours soupçonné que le poète nourrissait pour elle quelque prédilection. Car la strophe dont il la salue à son entrée est plus qu'un hommage rendu aux attraits d'une musicienne aimée du public, elle trahit une sollicitude pénétrée et presque émue dont on chercherait vainement trace en aucun autre endroit de son œuvre\note{\italique{Oiseaux}, vv. 676-684.}. 

La Trêve dans les \italique{Cavaliers} est un personnage du même genre, ainsi que Théoria et Opora dans la \italique{Paix}, Basiléia dans les \italique{Oiseaux}, et les autres que nous citons ci-dessus. 
Le scholiaste a soin de nous avertir la plupart des fois que le rôle est rempli par une hétaïre\note{Schol. sur \italique{Oiseaux}, 667 : ἑταιρίδιον πρόσεισι ; \italique{Cavaliers}, 1388 : πόρνας εἰσφέρει ; \italique{Ibid.}, 1390 : ἑταῖραι ὡραῖαι ; \italique{Paix}, 706 : ἦσαν δὲ ἑταῖραι, etc. Le poète le fait entendre lui-même au v. 849 de la \italique{Paix}.}, nécessairement choisie parmi les plus intelligentes et les plus belles. Car l'hétaïre grecque par son éducation est souvent tout autre chose qu'une vulgaire courtisane, et alors elle valait bien l'acteur. Cette Théodote, par exemple, que Socrate va visiter chez elle et à qui il donne de si piquants conseils\note{Xénophon, \italique{Memorab.}, III, 11}, devait être un sujet accompli. Maintenant faut-il s'étonner que chez ce peuple artiste avant tout et épris des pures lignes de la statuaire, les poètes se soient plu à faire paraître ces créatures de charme comme de vivantes statues, sans voile ni déguisement, dans toute la grâce harmonieuse et correcte de leurs formes ?


\Note{Note au vers 1341 des \italique{Guêpes}}

\citComm{ἀνάβαινε δεῦρο χρυσομηλολόνθιον.}\\


Ἀνάβαινεe, suivant M. de Wilamowitz, signifierait tout simplement : Avance. Il n'y a, dit-il, qu'à comparer ce passage de l'\italique{Apologie de Socrate} de Platon : ἡνίκ᾽ ἀνέβαινον ἐπὶ τὸ δικαστήριον (40b). 
Je ne saisis pas la pertinence de celte comparaison. On n'a jamais dit autrement que ἀναβαίνειν ἐπὶ δικαστήριον ; c'est le terme consacré pour « aller en justice » (\italique{Apol. Socr.}, 17d, \italique{Euthyd}., 35c, \italique{Gorgias}, 486b, Andocide, I, 4, 18, etc.), de même que ἀναβιβάζειν τινα est le terme courant pour « appeler en justice » (Démosthène, 440, 18, Isée, 78, 4, etc.). Pris à la lettre, ces deux verbes signifient : monter ou faire monter sur l'estrade du tribunal en qualité d'accusé ou de témoin. De même quand on se rendait au Pnvx, on était dit ἀναβαίνειν (ss-ent. εἰς τὴν ἐκκλησίαν) , et quand on en revenait, καταβαίνειν, le Pnyx étant situé sur une colline ; comme on disait ἀναβαίνειν lorsqu'on allait au Thesmophorion (\italique{Thesmoph.} , 281, 585), car ce temple était voisin du Pnyx ; et quantité de tournures analogues : ἀναβαίνειν (avec ou sans ναῦν), monter à bord, — (avec ou sans ἐφ᾿ ἵππον), monter à cheval, -- ἐπὶ τὸν ὀκρίβαντα ou ἐπὶ τὴν σκηνήν monter sur la scène (Platon, \italique{Banq.}, 19b, Pollux, IV, 127). 

Δεῦρ᾽ ἀνάβαινε, en dépit de toutes les subtilités, ne signifie et ne peut signifier que « monte ici ». C'est le mot du serviteur de Démos au charcutier se dirigeant vers l'agora: « Monte ici », tu y apercevras les files de peuple et les îles et le reste (\italique{Cav.}, 1 49) ; celui du paysan de Mégare à ses fillettes : ἄμβατε ποττὰν μᾶδδαν « Montez ici », on y a installé un marché, et peut-être y trouverez-vous de quoi vous emplir le ventre (\italique{Acharn.}, 732) ; et enfin le mot de Philocléon à sa flûtiste : « Monte ici, mon petit hanneton d'or », et entre chez moi, où tu n'auras rien à redouter des passants. 
Tout cela ne fait que confirmer la définition si claire d'Aristote, quand il oppose τὰ τοῦ χοροῦ, c'est-à-dire l'orchestre, et τὰ ἐπὶ τῆς σκηνῆς καὶ τῶν ὑποκριτῶν μέρος, c'est-à-dire la scène, que nos dictionnaires définissent encore aujourd'hui, comme au temps d'Aristote : partie du théâtre où jouent les acteurs. 
On traduit ἐπὶ τῆς σκηνῆς par « dans le voisinage de la scène », ce qui est un indéniable contresens : il faudrait au moins πρὸς τῇ σκηνῇ ou ἐγγὺς τῆς σκηνῆς. 

L'inverse de ἀναβαίνειν, au théâtre comme ailleurs, c'est καταβαίνειν. À la fin des \italique{Guêpes} (v. 1514), Philocléon ivre sort en dansant de chez lui el déclare : καταβατέον ἐπ᾿ αὐτούς μοι « Il me faut descendre vers eux [les Carkinites] ». On traduit : « 11 faut me mesurer avec eux », ce qui à la rigueur serait plausible, si le cas était isolé ; mais ce mot doit être rapproché de celui de la servante de Blépyros s'adressant à son maître, à la fin de l'\italique{Ecclésie des femmes} : ἐν ὅσῳ δὲ καταβαίνεις « pendant que tu descendras dans l'orchestre » (v. 1152). Car il est si peu question pour Blépyros de se mesurer avec les amies de sa femme, qu'au contraire il va les inviter à le suivre au banquet, ἐπὶ τὸ δεῖπνον ὑπανακινεῖν.

Ces verbes précédés de ἀνά ou κατά jouent aux éditeurs d'Aristophane les plus méchants tours. Lorsque dans les \italique{Acharniens} la fille de Dicéopolis, parée de ses bijoux et portant sur la main la tarte consacrée, dépose à terre sa ciste et dit à sa mère : ἀνάδος δεῦρο τὴν ἐτνήρυσιν « passe-moi la cuiller» (v. 245), quel est le fantaisiste qui s'écrie d'un ton triomphant : Vous voyez bien que δός et ἀνάδος c'est absolument la même chose ? Pas si vite, cher maître ; pas plus que βαίνειν et ἀναβαίνειν, ni que ἐλθεῖν et περιελθεῖν. 
Ne voyez-vous pas que la mère s'est accroupie devant la ciste et tendra de là à sa fille l'objet qu'elle réclame ? 

Pour nous en tenir au seul verbe ἀναβαίνειν, il se retrouve partout, exprimé ou sous-entendu, dans le théâtre grec. Quand ils n'entrent pas en scène par une des portes du fond, les acteurs y montent par l'orchestre\note{Pollux. IV, 127 : εἰσελθόντες δὲ κατὰ τὴν ὀρχήστραν ἐπὶ τὴν σκηνὴν ἀναβαίνουσιν διὰ κλιμάκων.}. 
Tel, dans l'\italique{Électre} d'Euripide, le vieillard qui éleva la fille d'Agamemnon, devenue la femme d'un simple laboureur ; parvenu au pied de la rampe qui mène à la chaumière, il regrette que le chemin soit si escarpé et si pénible, πρόσβασις ὀρθία (v. 489). 
Dans l'\italique{Ion}, le pédagogue, arrivé d'Athènes, trouve bien ardue la route qui mène au sanctuaire de Delphes, αἰπεινά τοι μαντεῖα (v. 739). 
Ailleurs Tirésias aveugle est obligé de s'appuyer sur un enfant pour gravir jusqu'au palais de Créon (\italique{Phénic}., v. 846). 
Dans \italique{Lysistrata}, les vétérans qui viennent assiéger les conjurées dans l'Acropole se vexent d'avoir à escalader sans l'aide d'un âne le raidillon qui s'élève jusqu'à l'entrée, τὸ σιμὸν τὸ πρὸς πόλιν (v. 288). 

S'il n'y avait point de scène, me fera-t-on la grâce de me dire ce que signifient ces divers témoignages, et pourquoi, par exemple, Euripide s'est avisé de placer la cabane d'Électre au sommet d'une colline, alors qu'il était si simple de la mettre en plein champ, comme aussi de situer le palais de Créon à front d'une place publique ? Et de pareils témoignages il serait facile d'en réunir tout un dossier.

Je sais que M. G. Perrot, dans une série d'articles du \italique{Journal des Savants}, répond à cela qu'en réalité les acteurs ne montaient pas du tout, mais faisaient semblant, et que le public, prêtant son concours au poète, voyait la pente en imagination. 
Ainsi, à la même époque où les Athéniens bâtissaient les Propylées et le Parthénon, c'est à ce piteux résultat qu'aurait abouti l'effort combiné de leurs poètes, de leurs architectes et de leurs ingénieurs, incapables, paraît-il, d'imaginer telle chose qu'une estrade du haut de laquelle les artistes dramatiques se seraient fait voir et entendre. 
Encore M. Perrot omet-il d'expliquer pour quel motif l'acteur, qui s'obstine à gravir sans nécessité aucune cette pente imaginaire, ne la descend jamais. Car enfin si Tirésias, le devin aveugle, cheminant en plein orchestre, s'échine et s'essouffle à atteindre un palais qui se trouve à six pas devant lui, si les vétérans de \italique{Lysistrata} suent sang et eau pour gagner les Propylées dont rien ne les sépare, pour quoi Héraclès, pourquoi Thésée et son ami Pirithoos, pourquoi Orphée ne font-ils pas les mêmes grimaces et contorsions en feignant de s'enfoncer dans l'Hadès ? 
Voilà du moins qui eût valu la peine, et du coup nos théâtres de marionnettes n'auraient eu plus rien à envier à la scène attique. 

Et voici tantôt vingt ans que la moitié de l'Europe savante s'alambique la cervelle, sans réussir à nous faire comprendre dans quel dessein le poète s'amusait à proposer à tous coups de captieuses énigmes grammaticales qui, comme la présente note le prouve abondamment, tiennent en échec les plus autorisés parmi les interprètes, mais que le public d'alors était tenu de saisir et résoudre à la volée.

\Note{Un commentaire du commentaire et son droit de réponse\\(ou la \italique{Passion de Philologus})}

D'antan, j'avais soumis quelques \italique{Notes sur les }Thesmophories au λύχνος, alias \italique{Connaissance hellénique}, revue fondée par Jean-Victor Vernhes et dirigée alors par Christian Boudignon.
J'y reçus un accueil si chaleureux que je retentai l'expérience quelques mois après ; la sagacité des lecteurs sollicitée ne fut pas sans retour ! Je reçus notamment une réponse insolite qui vaut d'être rapportée.\\
Ces \italique{impressions} -- comme la réponse s'intitulait --, à la lecture rapicolante par le style et la vaste culture, – pas même ma trop rapide incidente historique n'avait échappé à l’œil implacable de leur auteur –, ces impressions n'en restaient pas moins à mes yeux l'archétype des nuées embrouillaminant le texte d'Aristophane depuis que ce dernier eut la mauvaise idée de mourir... 
Et ces nuées, je n'aurai de cesse d'essayer de les dissiper par tous les moyens ! \\
La gifle étant alors fort opportunément revenue à la mode --  le cocaïnomane qui nous servait, paraît-il, de président, en fit la douloureuse expérience -- il m'avait semblé à la fois sain et juste de ne pas me priver de ce remède de grand-mère qui en un tournemain bien claqué peut remettre les idées à l'endroit et faire redescendre tous les philologus présents et à venir ! 

Je répondis donc à ces \italique{impressions} par une saynète vengeresse. (Pour se bien entendre, la courte pièce qui suit doit se lire équipé d'une puissante sono' pour les parties chorales. Il s'agit d'extraits de la \italique{Passion selon saint Matthieu} de Bach. Helmuth Rilling – évidemment ! – à la direction.)\\
Les voici toutes, \italique{impressions} et saynète, reproduites ci-dessous.

\subsection*{Impressions sur les Nouvelles notes thesmophoriennes parues dans le n°158 d’ὁ λύχνος (mars 2021)}

Disons-le tout de go, ces \italique{impressions} sont bonnes et, critère qui ne trompe pas, elles sont confirmées par la relecture attentive de l’ensemble, notes comprises, car les notes comprennent elles-mêmes des notes, ce dont il ne faut pas s’effaroucher. Non sans quelques émendations, comme nous le verrons... 

C’est un plaisir qui se mérite, pour le non-spécialiste que je suis, d’être confronté à un travail de fond portant sur la langue grecque, que j’apprécie sans la connaître autrement que par quelques notions éparses héritées de mes années de philosophie en Sorbonne. 

Travail d’érudition d’autant plus méritoire qu’il tape dans le dur, en s’attaquant à une comédie d’Aristophane qui s’épanouit dans un incroyable florilège d’allusions, allant d’Homère aux tragédies primées d’Euripide en passant par les gaudrioles de l’argot attique… 

Les déesses thesmophores, qui m’interdiraient l’accès à leur fête annuelle -- étant un homme -- me permettront-elles d’entrer dans le temple du savoir, moi qui ne suis même pas helléniste ? 

Mais avant d’endosser l’habit de Candide, je prendrai un instant la plume d’oie de M. Prudhomme pour … une petite remarque concernant le XIXe siècle : \\
\italique{« puissions-nous ne plus manger des herbes sauvages ni de scandix »… ce qui paraît tout-de-même avoir été des conditions plus enviables que celle des Parisiens lors de la Commune, réduits à manger des chats et des rats.} [note n°28] \\
Que la situation alimentaire ne fut guère brillante sous la Commune, elle était bien pire quelques mois plus tôt, sous le siège de Paris (17 septembre 1870 - 28 janvier 1871), qui a vu la population contrainte effectivement de manger des chats et des rats, sans omettre les éléphants du Jardin des Plantes. 
À noter que les singes ont été préservés du fait de leur trop grande ressemblance avec l’homme… 

Il s’agit là d’une simple remarque \italique{adventice} comme le \italique{scandix}, ou peigne de Vénus, la plante qui va nous occuper, étant au milieu du chemin de ces Nouvelles notes :\\
« \italique{Intéressons-nous maintenant plus particulièrement au scandix ; parmi tout l’étal que pouvaient proposer ces marchandes de verdure, seul le scandix est raillé par les comiques.} » \\
Reste à savoir pourquoi, et la réponse apportée s’avère peu convaincante : Le scandix servirait de condiment bon marché dont on pourrait assaisonner les plats. Mais, car il y a un ``mais"\\
« \italique{… le vers suivant : `ἀλλ᾽ εὑρέ τιν᾽ ἀπόκινον ἀπὸ τοῦ δεσπότου' laisse cependant circonspect : on ne comprend pas bien le rapport entre le scandix et l’ἀπόκινος, défini par le scholiaste comme une danse licencieuse et qui fait entendre par son étymologie le verbe ἀπο-κινέω,} se retirer. »\\

Nous devons donc rechercher, par nos propres moyens, quel est cet étrange privilège du scandix… Heureusement, Internet, dieu des temps modernes, vole à notre secours avec un 1\expo{er}  indice :\\
\italique{Le nom du peigne de Vénus (scandix pecten veneris) provient de la forme de ses fruits ridés avec de très longs becs qui, regroupés en ombelles, évoquent un peigne}.\\
(François Couplan : \underline{Les Plantes et leurs noms - histoires insolites}).\\
Va pour le peigne, mais pourquoi Vénus ? Les photos disponibles sur le net montrent que les dents du peigne, à la floraison, se couvrent de petites fleurs blanches qui pourraient peut-être évoquer la chevelure luxuriante de la déesse, lorsqu’elle se prend pour Flore…\\
Explication de bisounours addict à Bonne nuit les petits ! La vérité est autrement plus crue : la plante avec ses multiples petites fleurs blanches évoque la foufoune de Vénus au petit matin, après les ébats de la nuit.\\
Le scandix \underline{est donc l’évocation directe de débordements sexuels}, et c’est là l’unique raison qui fait son succès auprès des comiques. 
Nous comprenons maintenant le rapport auparavant mystérieux entre le scandix et l’apokinos, danse lascive. De même pour le retrait suivi de l’explosion florale… Nous remarquons que l’auteur des Nouvelles notes tourne à loisir autour du pot et semble souvent bien près de la vérité sans jamais l’atteindre… 

Ainsi, lorsqu’il commente la discussion des deux généraux qui débute Les Cavaliers d’Aristophane :\\
« \italique{On peut ainsi proposer la traduction pour le vers 19 des Cavaliers :\\
Δημοσθένης - μὴ ‘μοί γε, μὴ ‘μοί, μὴ διασκανδικίσῃς·\\
Démosthène - Arrête-moi ça : ne me saupoudre pas ça de scandix !} »\\
De même, un peu plus loin :\\
« \italique{Je crois avoir établi dans cette note - note qui tend à l’obésité, mais si c’est à la manière d’Eschyle, cela me convient ! - que si Aristophane parle si souvent du σκάνδιξ d’Euripide, c’est parce que le grec usait lui-même d’un verbe correspondant peu ou prou à notre ``enfariner" : διασκανδικίζειν.} » \\
Nous y voyons une bonne intuition, même si le verbe proposé demeure un peu faible pour restituer pleinement le sens. Nous préférerons la traduction plus crue \italique{empapaouter}. 
Nous ne résistons pas cependant à la tentation de citer un autre passage des Notes, qui ne manque pas de sel : \\
« \italique{Le verbe διασκανδικίζω n’est pas un hapax d’Aristophane ; il se trouve aussi dans un fragment de Téléclide, un autre poète comique contemporain d’Euripide. Ce fragment nous a été conservé par Athénée :\\
Ξυγγενέσθαι διὰ χρόνου †λιπαρείτω† με δρυπεπέσι, μάζαις καὶ διασκανδικίσαι.\\
Tout ce que l’on peut inférer du contexte chez Athénée, c’est que δρυπεπής, bien mûr, lorsqu’il ne se rapporte pas à un nom comme c’est le cas ici, désigne une olive. Il me semble bien avoir trouvé un sens possible à διασκανδικίζω qui rendrait compte des deux passages : non pas comme il est indiqué dans les dictionnaires ``se nourrir de σκάνδιξ" mais plutôt : ``assaisonner à l’aide de σκάνδιξ" et, dans un sens imagé, ``agrémenter", ``enjoliver" ; le σκάνδιξ jouant un rôle analogue à notre piment, et les multiples sens du verbe ``pimenter". 
Je construis donc : λιπαρείτω (?) με διὰ χρόνου συγγενέσθαι δρυπεπέσι καὶ διασκανδικίσαι μάζαις, ``Qu’il m’importune depuis longtemps pour ajouter des olives ainsi qu’agrémenter de scandix des galettes [!]".} »\\
Tout le monde comprenant qu’il ne s’agit pas d’une dispute portant sur une recette de cuisine, le sens - à peine caché - de la citation de Téléclide est le suivant :\\
« Ca fait longtemps qu’il veut me coller ses couilles au cul et me napper les fesses ! »\\
Le sens brut(al) étant restitué, il nous faut, pour traduire, trouver un équivalent du genre :\\
« On peut dire de moi que j’ai des c… au cul, mais encore faut-il que ce soit les miennes ! »

Ayant constaté que l’auteur des Nouvelles notes donne parfois des verges pour se faire battre, il ne manque pas, en outre, de fustiger le vil troupeau des pourceaux d’Épicure qui mettent leur groin dans la traduction. 
Prenons pour exemple la seconde femme intervenant à charge au procès d’Euripide : \\
« \italique{Je laisse à d’autres commentateurs le soin d’expliquer que la vendeuse de couronnes avoue à mi-mot être une petite vertu rentrant de travailler ; songez donc ! évoquer à la fois Chypre, le marché aux myrtes, la tresse de couronnes ! c’en était bien assez pour que nos fouines dégainent mécaniquement leur lecture }sensu obsceno\italique{ ! On repassera par contre pour obtenir de leur part un embryon d’explication :\\
1° en quoi l’impiété d’Euripide nuirait-elle effectivement à la prostitution ?\\
2° comment une mère de cinq enfants pourrait-elle se permettre de faire le tapin sans provoquer l’hilarité générale ?} » \\
Notre auteur manquerait-il le sens du passage en question ? Il s’agit, techniquement parlant, d’un intermède burlesque, après le trop long plaidoyer de la première intervenante qui, d’ailleurs, cible non seulement Euripide, mais aussi le machisme ambiant de la société athénienne. 
Pour dérider son public, Aristophane met en scène une matrone de petite vertu bien contente de paraître devant l’auguste assemblée des femmes pour recouvrer une légitimité qu’elle croit trouver suffisante dans l’activité de tresseuse de couronnes qui lui tient lieu de couverture. Bien entendu, elle s’empêtre entre ses deux activités, à la grande joie des spectateurs :\\
1° l’impiété d’Euripide, qui se communique au peuple, nuit sans doute à son travail de couverture, mais favorise, bien au contraire, son activité réelle : \\
« Mais je m’en vas à l’agora.\\
Car j’ai vingt couronnes à tresser à prix convenu. » \\
Bien entendu, ce glissement de sens est applaudi par les spectateurs ! À noter le bonheur de traduction trouvé par Louis Hauseux, inspiré de la fameuse phrase de Vaugelas sur son lit de mort : « Je m’en vais ou je m’en vas ; l’un et l’autre se dit ou se disent ».\\
2° l’hilarité générale est justement alimentée par ce second décalage entre l’apparence respectable de cette matrone athénienne mère de cinq enfants et ses appétits démesurés ; que son mari soit mort à Chypre signifie qu’elle l’a simplement tué au déduit… \\
Hypocrite lecteur, mon semblable, mon frère, tu as naturellement compris que je me plaçais du côté de ceux qui lisent Aristophane \italique{sensu obsceno}.

Encore faut-il le faire avec quelque talent, en suivant sa pente naturelle. Car des traducteurs fort sérieux, comprenant qu’il fallait parfois envoyer pour restituer la verve aristophanesque, s’y attellent-ils, mais en manquant totalement leur objectif. 
Prenons pour exemple les vers 286-291 :\\
« \italique{δέσποινα πολυτίμητε Δήμητερ φίλη\\
καὶ Φερρέφαττα, πολλὰ πολλάκις μέ σοι\\
θύειν ἔχουσαν, εἰ δὲ μή, ἀλλὰ νῦν λαθεῖν\\
καὶ τὴν θυγατέρα χοιρίον ἀνδρός μοι τυχεῖν\\
πλουτοῦντος, ἄλλως τ᾽ ἠλιθίου κἀβελτέρου,\\
καὶ πρὸς θάληκον νοῦν ἔχειν μοι καὶ φρένας. \\
-- Maîtresse très honorée, chère Déméter,\\
et toi Perséphone, puissé-je maintes et maintes fois\\
être à même de te sacrifier (à voix basse) et à tout le moins rester caché aujourd’hui.\\
(haut) Et que ma petite chatte de fille me trouve un mari\\
riche, d’ailleurs benêt et stupide …\\
et m’avoir l’esprit et l’intelligence appliqués à la cassette.} »

Louis Hauseux corrige le \italique{pros thalékon} du manuscrit de Ravenne en \italique{pros thalika}, tirant un trait sur le fils supposé de Mnésiloque, lequel fait de louables efforts pour se mettre dans la peau de la matrone athénienne moyenne ayant une fille à marier et un fils pas très dégourdi… \\
Mais prenons la traduction de ce passage par le très-honorable Hilaire van Daele, doyen de la Faculté des Lettres de l’Université de Besançon, pour la Collection Budé :\\
« \italique{-- Maîtresse tant honorée, chère Déméter,\\
et toi, Perséphoné, puissé-je échapper bien des fois [?]\\
à condition de te sacrifier aussi souvent, ou du moins de n’être pas reconnu aujourd’hui.\\
Puisse ma fille à la belle ``truie" trouver un mari\\
riche, par ailleurs sot et imbécile,\\
que de plus ma vergette de gars ait de l’esprit et du jugement.} »\\
Ma fille ``à la belle truie" ne passe pas en français, tandis que ma ``vergette de gars" est simplement ridicule. Comme quoi n’est pas obscène qui veut, et pour les vers litigieux, je proposerai la traduction suivante :\\
« Puisse ma gourgandine de fille se trouver un riche mari, si possible pas trop fute-fute, et que mon couillon de fils prenne un peu de plomb dans la cervelle ! »\\

Mais mes notes en réponse à des notes prennent déjà trop d’ampleur, et il faut savoir sacrifier pour ne pas lasser le lecteur… \\
Je finirai donc par le passage où Euripide, qui est entré lui-même dans l’action pour tenter de libérer son beau-père, en vient à parodier sa propre tragédie, La nouvelle Hélène. \\
Au vers 910, Mnésiloque-Hélène reconnaît en Ménélas-Euripide, vêtu en naufragé, son vieux mari, ὅσα γ᾽ἐκ τῶν ἀφύων, « d’après ton menu fretin » selon le manuscrit de Ravenne, le seul à comporter le texte intégral des Thesmophories et sans doute aussi le plus fiable. 

Bien entendu, ce passage passant pour obscur a fait l’objet de nombreuses interpolations divergentes que Louis Hauseux énumère avec un courage certain, la plus courante et non pas la moins ridicule étant qu’Hélène reconnaisse son mari ``à ses herbes" sinon ``à ses lavandes". \\
Passons, mais notre auteur, dans sa note n°64, fustige une de ces interpolations fautives qui cherche la vérité dans des oignons marins fort peu crédibles, et qui désigneraient le sexe masculin ! \\
D’où une sainte colère : \\
« \italique{Pour réussir à faire perdre son sang-froid au pourtant toujours modéré Colin Austin, il fallait effectivement que ce fut une belle salauderie… Nous aurions volontiers laissé la misérable hypo-foutaise ``qui se réfute d’elle-même" de Gag*** dans un oubli bien mérité, malheureusement, il nous a été infligé de la voir exhumée et prise au sérieux dans plusieurs articles récents… je crois donc de santé publique de l’autopsier une bonne fois pour toutes et d’enterrer - définitivement, espérons-le ! - cette charogne ! 
Ne nous attardons pas dans cette opération et allons droit au but : la saloperie de Gag*** présente à la fois le quadruple exploit : \\
1° de ne pas rendre compte de la leçon ‘ἀφύων’ du témoin principal de la pièce : le manuscrit de Ravenne ; \\
2° d’élucubrer sur un soi-disant sens obscène de τίφυον qui n’est attesté nulle part ; \\
3° d’être intrinsèquement absurde : si Ménélas est réputé avoir bien peu d’ardeur (raison du départ de l’insatisfaite madame ?), comment se fait-il qu’il soit en pleine érection et que ce soit à cela qu’Hélène le reconnaisse ?\\
4° last but not least, d’être anti-comique à souhait : un cocu impuissant n’est drôle que parce qu’il est impuissant : représentez-le en pleine érection et vous gâterez à coup sûr la farce !} » \\

Une fois de plus, n’est pas obscène qui veut : il ne saurait être question d’érection chez Ménélas naufragé sur cette plage d’Égypte et c’est justement ce que nous dit le manuscrit de Ravenne avec son \italique{menu fretin} : Hélène reconnaît son mari, qu’elle n’a plus revu depuis dix ans, \underline{à son petit zizi}. 

C’est là simple évidence, mais semble-t-il pas pour tout le monde, et c’est un peu étonnant de voir notre brillant helléniste, auteur des Nouvelles notes, sombrer dans des recherches inopérantes tendant à prouver qu’Hélène reconnaît Ménélas à ses hardes, alors qu’au contraire elle l’a toujours connu revêtu de riches atours… 

Mais n’oublions pas que ce qui donne de la force à cette parodie, c’est la référence, pour le coup incontournable, au chant VI de l’\italique{Odyssée}, avec la rencontre d’Ulysse et de Nausicaa, souverainement narrée par Homère et qui a fait fantasmer bien des auditeurs, jeunes et moins jeunes. 

Si la fille d’Alcinoos reste seule face à Ulysse, alors que ses servantes effrayées se sont égayées hors de la plage, c’est parce qu’Athéna a mis de la force dans son cœur, la force d’âme d’une future reine. 

Sans doute, les esprits portés à la mauvaise pente ont-ils de tout temps gauchi \italique{sensu obsceno} cette scène magnifique, en suggérant que la branche dont Ulysse naufragé dissimulait sa virilité avait légèrement glissé et suscité par-là même l’intérêt de Nausicaa, donnant, par contraste, sa force comique à la parodie aristophanesque aux dépens pêle-mêle de Ménélas, d’Hélène et d’Euripide… 

Ceci pour dire, en conclusion, que la beauté chez les Grecs n’est jamais loin, que ce soit dans le secret d’un temple ou, sous le soleil, sur la plage des Phéaciens face au rire éternel de la mer. 

Elle est toujours présente, même en arrière-plan, dans les textes qu’Aristophane éclabousse de son génie, scandix ou pas-scandix. 

\begin{flushright}Michel Devezeaux de Lavergne - avril 2021\end{flushright}

\subsection*{Droit de réponse aux \italique{Impressions}, ou \italique{La Passion de Philologus}}


Quatre personnages sur scène, et un chœur dans l'orchestre : \\
– Philologus, vieillard par l'âge et plus encore par l'aspect, une longue barbe sale et hirsute. Sorte de Michel Simon dont les yeux, embusqués derrière de petites lunettes rondes, exhalent le vice ; \\
– Bibi ; \\
– Pluton-Oui-Pluton-Non, le dieu des infernales tergiversations ; \\
– le fantôme d'Aristophane, bâillonné ; \\
– le chœur bacchien. \\

[Philologus toque à la porte de Pluton.]\\

\perso{Philologus} — Ouvre-moi ta porte, Pluton ! \\

\perso{Pluton} — Qu'on cesse ce raffut ! Τίς ἐσθ’ ὁ κόψας τὴν θύραν; \\

\perso{Philologus} — C'est moi : Philologus de Cancrebridge. Je cherche Aristophane : j'ai imaginé comment améliorer en maints endroits ses \italique{Thesmophories} et je tiens à lui en faire part ! \\

\perso{Pluton} — Aristophane ? Il est bien ici... et sous bonne garde ! cet empêcheur-de-mourir-en-rond ne cesse de vouloir faire la tombe buissonnière depuis qu'il est mon pensionnaire. Aristophane, viens par là ! \\

[Arrive le fantôme d'Aristophane, bâillonné et fortement entravé, gardé par une cohorte de rachitiques à la mine patibulaire ; d'aucuns prétendent qu'Aristophane fut ainsi enchaîné et mis sous bonne garde par ses propres commentateurs... Une voix tonnante et pleine d'une violente colère se fait entendre au loin.] \\

\perso{Philologus} — Qu'est-ce là ? j'entends le tonnerre... mais où sont les éclairs ? \\

\perso{Chœur} — \allemand{Sind Blitze, sind Donner\\
In Wolken verschwunden ?}\note{Le texte allemand de la \italique{Passion selon saint Matthieu} de Bach, la numérotation ainsi que les traductions françaises sont tirés du Cantagrel : \italique{J.-S. Bach, Passions Messes et Motets}.\\
« Les éclairs, le tonnerre\\
ont-ils disparu dans les nuages ? » 27b.}\\

\perso{Pluton} — Le tonnerre ? détrompe-toi : c'est Bibi de Thunder-ten-Trunk ! On entend sa voix tonner ainsi dans le désert universitaire depuis quelque temps. Le voilà qui vient... \\

[Bibi entre sur scène tonitruant. Les éclairs de son regard et sa crinière ébouriffée lui donnent des faux airs de Gorgone. Il aperçoit Aristophane pieds et poings liés.]\\

\perso{Bibi} — Ne touchez plus à Aristophane ! éloignez-vous de lui ! \\

\perso{Bibi}\textbf{, accompagné du chœur} — \allemand{Lasst ihn, haltet, bindet nicht !}\note{« Lâchez-le ! Arrêtez ! Ne l’enchaînez pas ! » 27a.}\\

[La cohorte rachitique fait bloc autour du pauvre Aristophane que Bibi tente désespérément de délivrer...] \\

\perso{Chœur} — \allemand{Sie führen ihn, er ist gebunden.}\note{« Ils l’emmènent. Il est lié. » 27a.}\\

\perso{Bibi + chœur} — \allemand{Lasst ihn, haltet, bindet nicht !} \\

\perso{Philologus} (regrettant tout d'un coup de s'être mis dans un tel pétrin) — Peut-être serait-il prudent pour moi de prendre mon courage à deux jambes... \\

[Bibi lui bloque toute retraite.] \\

\perso{Bibi} — Race de vipère ! qui t'a appris à fuir la colère qui vient ? C'est à Aristophane que tu dois tous tes titres universitaires ; est-ce ainsi que tu le remercies : en le mordant de ton venin, vil serpent ! \\

\perso{Chœur} — \allemand{Ach, ein Kind, das du erzogen, \\
Das an deiner Brust gesogen, \\
Droht den Pfleger zu ermorden, \\
Denn es ist zur Schlange worden.}\note{« Hélas ! un enfant que tu as mis au monde, \\
que tu as nourri de ton sein, \\
menace la vie de son bienfaiteur, \\
car il est devenu serpent. » 8.} \\

\perso{Bibi} — Je me suis affligé la lecture de tes \italique{Impressions} et, disons-le tout de go, elles sont très mauvaises : \italique{ça} m'en est tombé des mains, j'en suis resté littéralement médusé... Maintenant est venu le temps de la fureur ! il te faut payer pour ton crime de lèse-Aristophane ! \\

\perso{Bibi avec le chœur} — \allemand{Eröffne den feurigen Abgrund, \\
O Hölle, \\
Zertrümmre, verderbe, \\
Verschlinge, zerschelle \\
Mit plötzlicher Wut \\
Den Falschen Verräter, \\
Das mördrische Blut !}\note{« Ouvre ton abîme de feu, \\
ô enfer. \\
Pulvérise, ruine, \\
engloutis, fracasse \\
d’une soudaine fureur \\
le faux traître, \\
le meurtrier sanguinaire ! » 27b.} \\

\perso{Pluton} — Allons, Bibi !.. on entend les gens avant de se fâcher... \\

\perso{Bibi} — Hors qu’un commandement exprès d'Aristophane me vienne de trouver bonnes ces impressions, je soutiendrai toujours qu'elles sont viciées et qu’un homme est pendable après les avoir faites. \\

\perso{Bibi avec le chœur} — \allemand{Lass ihn kreuzigen !}\note{« Qu’il soit crucifié ! » 45b.}\\

\perso{Pluton} — Qu'a-t-il fait donc de si grave ? \\

\perso{Bibi} — Rien que la mort n'est capable d'expier son forfait ! \\

\perso{Bibi + chœur} — \allemand{Er ist des Todes schuldig !}\note{« Il est passible de mort ! » 36b.}\\

\perso{Pluton} — Cette attitude ne me plaît guère ! je te trouve assez plaisant de vouloir user d'un tel empire. On n'est jamais trop lent pour faire périr un homme ; la Camarde est assez vigilante, elle n'a pas besoin qu'on lui tienne la faux. Dois-je vous rappeler que vous êtes ici aux portes de l'Enfer et que j'en suis le seul maître et dieu ! 
J'arbitrerai donc votre contentieux mais avant cela, je vous saurai gré de me rendre les hommages dus, ainsi que l'exigent l'usage et les bienséances ! \\

\perso{Bibi} — Pardonne-moi, ô Pluton : je m'emporte et manque effectivement à tous mes devoirs. Il me sera bien facile de te rendre hommage, ô toi le dieu des morts, en même temps qu'à mes chères Muses. \\
Mes maîtresses les Muses et la Mort sont intimement liées : si la Mort nous renvoie à notre finitude de mortels, c'est bien les Muses qui nous en délivrent par leur promesse d'immortalité. C'est donc toi et ton auxiliaire la Mort qui êtes le cœur de toute création artistique ; les dieux immortels n'ont pas besoin de ce succédané d'éternité : ils sont uniquement tournés vers les plaisirs... mais c'est bien là une activité incompatible avec notre condition d'homme ! \\
Ô Muses, mes éternels secours, restez toujours à mon chevet ; \\
Car vous seules savez me consoler. \\
Après les coups, les ruines, séparations ; \\
Lorsque mes larmes humaines couleront,\\
Vous seules saurez me chanter comment le bel art incorporé à d'étincelantes matières \\
Fait des promesses de salut définitif à la beauté comme à la vertu. \\
Vous seules savez que les substances sacrées ne peuvent s'altérer \\
Mais font retour au cœur du monde Où chacun les recouvrera... \\
Pour peu que le monde ait un cœur... \\

\perso{Chœur} — \allemand{Wenn ich einmal soll scheiden, \\
So scheide nicht von mir ! \\
Wenn ich den Tod soll leiden, \\
So tritt du denn herfür ! \\
Wenn mir am allerbängsten \\
Wird um das Herze sein, \\
So reiss mich aus den Ängsten \\
Kraft deiner Angst und Pein !}\note{« Quand, un jour, je devrais m’en aller, \\
ne te sépare pas de moi ! \\
Quand il me faudra souffrir la mort, \\
montre-toi à moi ! \\
Quand mon cœur sera pris \\
des plus terribles angoisses, \\
arrache-moi alors à mes angoisses \\
par ta propre angoisse et ta souffrance ! » 62.} \\

\perso{Philologus} — Peuh... la sempiternelle ritournelle sur la beauté qui sauvera le monde... et d'abord, où as-tu vu que le monde avait un cœur ? Il n'y a que des guerres et des angoisses et des horreurs ; de cœur, je n'en ai encore jamais perçu la moindre trace. \\

\perso{Bibi} — Tu crois donc qu'il n'y a pas de justice ? que Dieu –  \\

\perso{Philologus} — Dieu ?! c'est l'Imagination de l'homme ! Mais laisse-moi plutôt répondre à tes pleurs d'enfant inconsolable devant les ténèbres par une autre prière : \\
Que valent donc amours et larmes des aimables Muses ?\\
Elles ne font que retenir la Forme Immortelle – par des liens certes tendres ! –\\
Mais qui ne sont que des abîmes sur les plaines ravagées de la Raison :\\
No-man's-lands où règne l'abstraction philosophique menant sa méchante guerre contre l'\italique{Imagination}\\
(laquelle est la seule puissance capable d'immortaliser l'homme ;\\
de conférer à ce phoque rampant et se tordant dans sa fange \\
Une destinée divine ; de hisser à l'égal de Dieu\\
Celui qui serait le maître de cette Puissance de l'Esprit.)\\
Nuées étourdissantes, Éther vénérable,\\
Saisissez-vous de moi : emportez-moi dans un de vos tourbillons\\
Où Sensations et Esprit se trouvent liés \\
– Liés entre eux et non aux contingences matérielles –, \\
Déchaînez mon Imagination engourdie : \\
Vous ferez ainsi mes quatre voluptés. \\

\perso{Pluton} — Les prières d'usage étant faites, passons à la confrontation ; vous connaissez le principe : j'ai apporté la balance qui vous départagera : il suffit que chacun de vous y dépose ses conjectures et la balance penchera du côté de la meilleure interprétation. \\

\perso{Bibi} — Mais que ne demande-t-on directement à Aristophane ici présent ? \\

\perso{Pluton} — L'adage le dit : on ne fait pas parler les morts. \\

\perso{Bibi} — C'est vrai : les morts sont d'or pour cela !.. et cela me donne bien des idées pour faire cesser les radotages de Philologus !!... N'aie crainte, Aristophane, ce que tu ne peux plus faire par toi-même, je m'en chargerai ! Cède-moi un peu de ta verve vengeresse, que je puisse faire justice de la tourbe dont tous les Philogugus du monde t'ont maculé. Ô Ciel ! de mes transports puis-je être ici le maître ? \\

[Aristophane acquiesce.] \\

\perso{Bibi} — Et puisque \italique{nous} sommes les offensés, à nous le choix des armes ! commençons par remonter les bretelles de ce bretteur-pornographe ; Philologus, prends garde à ta scandaleuse scandhypofoutaise ! \\

\perso{Philologus} — Je maintiens que le scandix est l'évocation directe de débordements sexuels, débordements suggérés par la multiplicité des petites fleurs blanches tels la foufoune de Vénus à l'aube qui suit les nuits les plus agitées. \\

\perso{Bibi} — Ô le galant homme qui affuble gracieusement le plus bel apanage de Vénus de fleurs blanches... mais as-tu seulement pris garde que ta chimère reposait sur le nom latin de la plante, \italique{Pecten Veneris} (si tant est qu'elle soit correctement identifiée)... et qu'il n'y a rien de tel dans le grec σκάνδιξ ! Si Aristophane avait voulu suggérer à son auditoire quelque idée de stupre, la Flore attique surabondait de telles fleurs érotiques : παρθένιον, βλήχων, μύρτον, et j'en passe et des plus suggestives... 
Quant à continuer de rapprocher le scandix de l'apokinos, c'est persévérer dans ses turlutaines quand toute l'intelligence du passage est restituée par l'apok\textbf{y}nos (avec un `i' \italique{grec} et non latin, encore une fois !) \\

\perso{Pluton} — Le contentieux ayant été mis à plat, passons à la pesée ! \\

\perso{Philologus} — Mais je n'ai même pas pu parler... \\

\perso{Bibi} — Tu ne l'as déjà que trop fait ; la pesée ! \\

[Les propositions sont pesées ; la balance penche du côté de Bibi.] \\

\perso{Philologus} — Comment ?! mais l'explication sexuelle est l'unique raison pouvant expliquer le succès du scandix auprès des Comiques !.. \\

\perso{Pluton} — Regarde sur la balance ; qu'est-il marqué ? \\

\perso{Philologus} — Τά... λαν... τον : τάλαντον ! \\

\perso{Pluton} — Oui-da : τάλαντον, la balance, et non \italique{bilanx} ; elle juge du grec et non du latin ! \\

\perso{Bibi} — Dans la même veine, parlons de la tresseuse de couronne et de feu son mari, mort à Chypre ! \\

\perso{Philologus} — Cette fois, je suis inexpugnable : j'ai résolu la contradiction apparente entre la plainte de la tresseuse de couronne dont les affaires ne vont pas fort à cause de l'impiété d'Euripide, et les commandes qui affluent pourtant par vingtaines : elle s'emmêle les fils entre ceux de sa couverture (tresseuse de couronnes, métier honorable s'il en est !) et ceux de sa véritable couverture, dont elle a besoin pour son activité de petite vertu. \\

\perso{Bibi} (rongeant son frein depuis tout-à-l'heure) — Patience !.. patience... je te réglerai ton compte d'un grand coup de force et de rage... \\
\allemand{Geduld ! \\
Wenn mich Falsche Zungen stechen ! \\
Leid ich wider meine Schuld \\
Schimpf und Spott, \\
Ei, so mag der liebe Gott \\
Meines Herzens Unschuld rächen.}\note{« Patience ! \\
Quand me piquent les mauvaises langues, \\
quand par injustice je souffre \\
outrages et railleries, \\
courage ! Le Dieu aimé \\
vengera l’innocence de mon cœur. » 35.} \\

\perso{Philologus} — C'est aussi l'occasion pour Aristophane de dérider ses spectateurs avec un intermède burlesque. \\

\perso{Bibi} — Un intermède burlesque ?! quelle idée originale !.. mais j'y pense, on y a déjà eu droit précédemment avec la scène d'épilation... un autre va bientôt suivre avec le déshabillage de Mnésiloque... pour faire patienter le spectateur entre ces deux intermèdes burlesques, quoi de mieux qu'un intermède burlesque ?.. c'est donc ainsi que tu te représentes tout l'art d'Aristophane : nous écrire une longue suite ininterrompue d'intermèdes burlesques !.. \\

\perso{Philologus} — Je maintiens : l'hilarité générale est justement amenée par le contraste entre la respectabilité que se dessine cette matrone athénienne mère de cinq enfants et ses appétits démesurés ; que son mari soit mort à Chypre signifie qu'elle l'a simplement tué au déduit. \\

\perso{Bibi} — Mourir de sa belle mort, comme l'heureux \italique{Félix Faurenicator} ; crois-tu que je n'y ai pas songé ? L'hypothèse est plaisante mais gratuite et donc à rejeter ! tous les lexiques consultés ne soufflent mot d'une telle expression qu'on aurait nécessairement relevée. D'autres se sont essayés à prendre ce vers littéralement : là encore, en se heurtant à un mur : les dernière expéditions sur l'île de Chypre feraient de notre procurette une septuagénaire au moins !.. et comme il ne semble pas y avoir de jeux de mots, il ne reste qu'une seule voie : une erreur de transcription du texte. C'est d'autant plus vraisemblable que l'expression « ἀπέθανεν ἐν Κύπρῳ », si elle s'était véritablement trouvée dans le texte original, n'aurait pas manqué d'intriguer les commentateurs qui, même sans y rien connaître, auraient ouvert les vannes de leur verve intarissable... au contraire : ce vers n'a étrangement pas laissé la moindre scholie derrière lui. Contrairement à ce que tu soutiens, cette intervention n'a rien de burlesque : elle essaie au contraire d'introduire du pathétique. Or, qu'y a-t-il de plus banal pour apitoyer son auditoire à peu de frais que de sortir le lieu-commun de la femme qui se retrouve veuve alors qu'elle est enceinte ? \\
ἐμοὶ γὰρ ἁνὴρ ἀπέθανεν μὲν \textbf{ἐγκύμων} \\
Mon mari mourut alors que j'étais enceinte \\
Me laissant cinq enfants à charge etc. etc. \\

\perso{Philologus} — Ah ! voilà notre champion de la tradition et du respect filial du texte transmis soudainement mué en tripatouilleur de vers !!! \\

\perso{Bibi} — La véritable tradition est critique ! Et la retouche a ses lois : elle ne doit être qu'un pis-aller, le dernier recours après avoir épuisé toutes les autres pistes, ce qui me semble être le cas ici. De toute façon, notre intervenante ne peut en aucun cas être la nigaude que tu imagines : si le chœur vante son esprit quelques vers plus loin, il faut l'en croire ! \\

\perso{Pluton} — Passons à la pesée : à vos remarques ! prêts ? pesez ! \\

\perso{Bibi} — J'y dépose mon émendation en tremblant, espérant qu'elle n'amaigrisse pas ma lecture \italique{sensu pathetico}... \\

\perso{Philologus} — \italique{Pathice} ! Et moi j'y balance sans trembler ma lecture \italique{sensu obsceno}... Et en prime, j'y balance mon porc d'Épicure ; il pèsera de tout son poids en ma faveur ! \\

[Entre une troupe de Bacchantes en furie qui se jette sur le porc de Philologus et le met en pièce avec la délicatesse qu'on connaît aux féministes.] \\

\perso{Philologus} — Ces charmantes demoiselles ont tranché net : cette fois, je l'emporte ! \\

\perso{Pluton} — Au contraire, comme dit le poète : \\
« Malheureusement, si les filles \\
Aux jeux de l'amour valent leur prix, \\
On sait qu'en revanche elles ne brillent \\
Ni par le goût, ni par l'esprit. » \\
Ces bacchantes sont nos boussoles indiquant le sud ! À Bibi la victoire ! Point suivant. \\

\perso{Philologus} — Cette fois, je prends la main (et je sais m'en servir !). Laissez-moi vous déniaiser un tant soit peu : n'est pas docteur \italique{ès scato' \& porno'} qui veut, sots puceaux que vous êtes ! Si, au vers 910, Hélène reconnaît le pseudo-Ménélas à ses ἀφύων, son \italique{menu fretin}, c'est l'évidence qu'elle pointe son \italique{petit zizi} (à défaut de l'inverse). Et toute la force de cette parodie, c'est la référence, pour le coup incontournable, au chant VI de l'\italique{Odyssée}, avec la rencontre de Nausicaa, souverainement narrée par Homère et qui a fait fantasmer bien des auditeurs, jeunes et moins jeunes. Si la fille d’Alcinoos reste seule face à Ulysse, alors que ses servantes effrayées se sont égayées hors de la plage, c’est parce qu’Athéna a mis de la force dans son cœur, la force d’âme d’une future reine. Sans doute, les esprits portés à la mauvaise pente ont-ils de tout temps gauchi \italique{sensu obsceno} cette scène magnifique, en suggérant que la branche dont Ulysse naufragé dissimulait sa virilité avait légèrement glissé et suscité par-là même l’intérêt de Nausicaa, donnant, par contraste, sa force comique à la parodie aristophanesque aux dépens pêle-mêle de Ménélas, d’Hélène et d’Euripide... \\

\perso{Bibi} — Voilà du nouveau ! d'émoustillantes équivoques chez Homère ?.. un tel flair devrait te permettre d'aller débusquer de l'ordure jusques dans les évangiles. Mais ici, ton imagination résonne à vide : tu t'appuies sur un texte qui ne respecte même pas le mètre du vers ! Encore un fois, l'obscénité t'est suggérée par la traduction française de « menu fretin », fais le choix de traduire \italique{aphyes} par « anchois » et ton hypothèse s'effondre ! \\

\perso{Philologus} — Qu'importe le vers, pourvu qu'on ait l'ivresse de l'Imagination ! \\

\perso{Bibi} — Pourvu qu'on ait lu vrai ! Avec toi, l'on a l'ivraie, c'est vrai !.. mais jamais le bon grain... \\

[Aristophane profite de l'inattention de ses geôliers, ébahis par les savants rapprochements de leur maître Philologus, pour se défaire de ses liens et les assommer un à un. Il se débâillonne et se lance dans une tirade inspirée de ses lectures d'outre-tombe :] \\

\perso{Aristophane} — Ô philologues ! philo-logues ! PHILOLOGUES ! créatures faibles et décevantes ; nul animal ne peut manquer à son instinct, le tien est-il donc de te tromper ? Pourvu que je ne parle en mes écrits ni de figue, ni d'olive, ni de galette, ni même de légumes dont vous allez aussitôt chercher les vertus prétendûment aphrodisiaques, de poissons dont vous farfouillerez les vieux grimoires pour dénicher une traduction portant à équivoque ; ni bien-sûr d'individus du beau-sexe, fussent-elles fanées par une vie de labeurs et de quinquenfantements, ni de tout mot qui, de loin plutôt que de près, ait pu avoir dans l'histoire de l'humanité une acception déshonnête, je puis tout écrire librement sous l'inspection de vos groins fureteurs sans que vous y détectiez grivoiserie. Ô... que vous me donnez de tourments !.. \\
Mais le verbiage est inutile face à cette racaille de la pensée ! Qu'on m'apporte du feu ! Je ferai de sorte que chacun d'eux, si fanfarons qu'ils soient, reçoive aujourd'hui son châtiment ! Au bûcher tous ces marmiteux, ces calamiteux de la recherche ! \\

\perso{Philologus} (en train d'imaginer assez aisément le sort qui lui est réservé) — Ô ma chère Imagination, cesse de me faire entr'apercevoir ce qui m'attend... Je crois qu'il est venu le temps de me mettre en règle... comment cela commence-t-il déjà ?.. ah ! oui : \\
Notre Père qui êtes dans le Sur-Moi, \\
Donnez-nous aujourd'hui notre libido quotidienne, \\
Pardonnez-nous nos offenses, c'est-à-dire nos agressions sadiques, \\
Comme nous les pardonnons masochistiquement à ceux qui nous ont offensés. \\
Mais ne nous laissez pas succomber à la névrose. \\
Ainsi soit-il. \\

\perso{Pluton} — Halte-là ! ils doivent payer par là où ils ont péché : je les condamne plutôt à lire éternellement du Euripide, les œuvres écrites, les œuvres perdues, et celles que ce boulimique à l'écriture compulsive ne cesse d'écrire depuis 2500 ans ! \\

\perso{Aristophane} — Soit !.. mais je suis toujours prêt à commuer cette peine en peine de mort. \\

\perso{Philologus} — Ô Aristophane, pardonne-moi ! Si à te servir, je ne suis d'aucun prix, j'aurais du moins l'honneur de l'avoir entrepris. Ma mauvaise pente t'aura causé du tort. Hélas ! mes péchés t'ont frappé. \\

\perso{Philologus $+$ chœur} — \allemand{Ach, meine Sünden \\
Haben dich geschlagen...}\note{« Hélas ! mes péchés \\
t’ont frappé. » 19.} \\

\perso{Philologus} — Aie pitié de moi.... aie pitié de moi... \\
\allemand{Erbarme dich, \\
Mein Gott, um meiner Zähren willen ! \\
Schaue hier, \\
Herz und Auge weint vor dir \\
Bitterlich.}\note{« Aie pitié de moi,\\
mon Dieu, au nom de mes larmes !\\
Vois ici,\\
mon cœur et mes yeux pleurent amèrement devant toi. » 39.}\\

\perso{Aristophane} — Quant à toi, Bibi le-bien-nommé, suis-moi : tu vas pouvoir t'abreuver jusqu'à la consommation des temps à la source de mes pièces ! \\

\perso{Bibi} — Vais-je y rencontrer Didyme et Hésychius conversant depuis bientôt deux millénaires sur tes œuvres ? \\

\perso{Aristophane} — Mieux que cela : je vais personnellement t'en révéler tous les secrets.\\

\begin{figure}[!h]
\centering
\includegraphics[scale=0.65]{fin.png}
%\caption{}
%\label{Encycle}
\end{figure}

\cleardoublepage

\end{document}